%% EXHALE: running update log, stage 3 (from 2026-09-23).
%% GENERATED from Update_EXHALE_stage3.md by src/utils/update_log_to_tex.py;
%% edit the Markdown, not this file.
%% Author: Kwang-Il Seon
\documentclass[english,a4paper]{my_memo}
\synctex=1
\usepackage{listings}
\usepackage{xcolor}
\usepackage{booktabs}
\usepackage{longtable}
\usepackage{array}
\usepackage{calc}
\usepackage{url}
\urlstyle{tt}
\usepackage{soul}
\usepackage{babel}
\providecommand{\tightlist}{\setlength{\itemsep}{0pt}\setlength{\parskip}{0pt}}
\providecommand{\pandocbounded}[1]{#1}
\providecommand{\passthrough}[1]{#1}
\newlength{\cslhangindent}
\newlength{\csllabelwidth}
\setcounter{secnumdepth}{0}
\sloppy
\begin{document}

\title{EXHALE: Update Log (stage 3, from 2026-09-23)}
\author{Kwang-Il Seon}
\date{Last updated: \docmoddate}
\maketitle

\tableofcontents
\bigskip

This log continues \texttt{md/Update\_EXHALE\_stage2.md} from 2026-09-23. The
earlier logs are kept under their own names:

\begin{itemize}
\tightlist
\item
  \texttt{docs/Update\_EXHALE\_stage0.\{tex,pdf\}}: the early phase (the ATES fork up to
  2026-08-15).
\item
  \texttt{docs/Update\_EXHALE\_stage1.\{md,tex,pdf\}}: sections 1-171, 2026-06 to
  2026-09-05; "Update\_EXHALE section N" in the source comments and the older
  memos points into that file.
\item
  \texttt{docs/Update\_EXHALE\_stage2.\{md,tex,pdf\}}: sections 1-17, 2026-09-05 to
  2026-09-22.
\end{itemize}

The entries of 2026-09-23 were first written into the stage 2 log as its
sections 18 to 24 and moved here the same day as sections 1 to 7; a record of
that day that cites "\texttt{md/Update\_EXHALE\_stage2.md} section 18" to "section
24" means section 1 to section 7 of this file (the references known at the
move were rewritten).

Each entry below is one dated section: what changed, why, which files, the
check that was run and its result. Section numbers restart at 1. The plan the
entries execute is \texttt{md/PLAN\_20260923\_rev2.md}.

\begin{center}\rule{0.5\linewidth}{0.5pt}\end{center}

\subsection{1. The Ritz residuals of the preconditioned spectrum (2026-09-23)}\label{1-the-ritz-residuals-of-the-preconditioned-spectrum-2026-09-23}

The spectrum hook \texttt{EXHALE\_PRECOND\_SPECTRUM=1} printed Ritz values with no
accuracy of their own (\texttt{md/ISSUES\_20260909.md} item (D8), plan section 16).
It now prints the residual \texttt{\textbar{}\textbar{}A\ y\ -\ theta\ y\textbar{}\textbar{}} of each.

\textbf{The residual.} For a Ritz pair \texttt{(theta,\ y\ =\ V\_k\ s)} with \texttt{\textbar{}\textbar{}s\textbar{}\textbar{}\ =\ 1} the
Arnoldi relation \texttt{A\ V\_k\ =\ V\_k\ H\_k\ +\ h(k+1,k)\ v(k+1)\ e\_k\^{}T} gives
\texttt{A\ y\ -\ theta\ y\ =\ h(k+1,k)\ s\_k\ v(k+1)}, so the residual is \texttt{\textbar{}h(k+1,k)\textbar{}\ \textbar{}s\_k\textbar{}}.
\texttt{ritz\_values\_of\_a\_matrix\_free\_operator} now hands back \texttt{h(k+1,k)}, and the
routine prints that residual and the residual over the value\textquotesingle s magnitude for
the six smallest values of every compression. That form assumes a linear
operator and an orthonormal basis, and the finite-difference action is linear
only to its truncation, so for the three smallest values of the whole operator
the residual is also measured by applying the operator once more to \texttt{y}, once
to each part of a complex pair. The operator action is now one routine,
\texttt{action\_of\_the\_preconditioned\_operator}, called by the recursion and by the
measured residual alike; the complex-pair packing of dgeev is undone in
\texttt{coefficients\_of\_the\_ritz\_vector}. Every existing \texttt{{[}diag\ 15{]}} line is printed
as before, and the residuals are new lines.

\textbf{Verification.} On the stored state of
\texttt{atomic\_scalar\_gj1132x0.01\_kzz1e9/HeH9.7}
(\texttt{g0002\_20260919T004843Z\_f7485b14}) with \texttt{EXHALE\_TRUST\_REGION=1},
\texttt{EXHALE\_ELEM\_DIAG=1}, \texttt{EXHALE\_PRECOND\_SPECTRUM=1} and \texttt{EXHALE\_JFNK\_MAXIT=1}, a
17-second run: the 30 \texttt{{[}diag\ 15{]}} lines of the binary before the change are
byte-identical after it, so extracting the action moved nothing in the
recursion; and the Arnoldi estimate and the measured residual agree to five or
six digits for all three measured pairs, the complex one included. MEASURED.
The hook runs only under \texttt{EXHALE\_PRECOND\_SPECTRUM=1} inside the trust-region
step, which no case of the regression matrix enables, so the matrix was not
rerun. Two unused local declarations of the same file (\texttt{islot}; \texttt{tv}, \texttt{net})
were removed on the way; they carried no value.

\textbf{What it shows.} At 200 products that spectrum is the one phase 9 of the
campaign quoted (largest 9.99385E-01, smallest 2.52128E-04, ratio 3.96381E+03,
reproduced to every digit), and its small end is not resolved:
\texttt{h(k+1,k)\ =\ 0.168501}, and the residual over the magnitude is 12.8, 3.58, 1.86,
1.86, 1.09 and 1.09 for the six smallest values, every one larger than one.

\textbf{The whole space.} \texttt{n\_ritz\_step} became \texttt{EXHALE\_RITZ\_STEPS=\textless{}k\textgreater{}} (default 200,
capped at the number of unknowns, below 2 refused). At k = 1500, the number of
unknowns of this operator, the Hessenberg is similar to the operator, and the
same fixture took 136 s. MEASURED: smallest 5.60717E-06, ratio 1.78560E+05,
43 values below 1E-4 and 319 below 1E-2, against 0 and 6 at 200 products. So
"nothing below 1E-4" is wrong on that checkpoint, and "there is no near-null
mode of the preconditioned operator" is refuted. What the recursion finds are
small eigenvalues of the operator the Krylov cycle sees: the hydrodynamic block
alone (three unknowns a cell; \texttt{nvar\_jac\ =\ 3\ +\ nspec\_row}, with no species rows
on this route), scaled and preconditioned. {[}Corrected later on 2026-09-23 by
\texttt{md/lhs1140b\_spectral\_checks\_20260923.md}: this read "and shifted by the
pseudo-time diagonal"; the shift is in the band preconditioner only, the hook
being handed \texttt{idtau\_leg\ =\ 0}.{]} They are not established as time scales of the
atmosphere, and they say nothing about the alternated composition-hydrodynamics
map. {[}Wording corrected 2026-09-23 after
the review of \texttt{md/PLAN\_20260923.md}: this paragraph read "it has a cluster of
slow modes".{]} Below about 1E-5 even the full recursion does not resolve them:
the residual measured on the operator is 9.72E-06, 9.88E-06 and 9.56E-06 for
the three smallest values (the first two one conjugate pair), larger than the
pair\textquotesingle s magnitude and 0.96 of the third\textquotesingle s, so none is resolved from zero at the
production probe length; the floor was later measured to be the rounding of
the forward difference (it falls as arc\^{}-0.97 with the probe length, spectral
checks memo). What the measurement bounds rigorously is the
smallest singular value of the action as applied: on the normalized Ritz
vector its norm is at most \textbar theta\textbar{} + \textbar\textbar r\textbar\textbar{} = 1.5E-05. Whether the 1E-5 floor
of the residual is truncation, rounding or nonadditivity of the
finite-difference action is not decided by these runs;
\texttt{measure\_the\_additivity\_of\_the\_jacobian\_action} is the narrower test. The
vectors (the real and imaginary parts of the pair, and the third) carry 53 to
61 per cent in the mass rows and peak in the wind around cells 93 to 169, not
at the base, but since they are not resolved their location is not
established either. The documents that quoted the 200-product numbers carry
a dated correction beside them. What stands is the direct measurement they
were offered beside: on the alternated route the Krylov cycle reaches its
tolerance in one to four products.

\subsection{2. Two records the campaign left open, closed (2026-09-23)}\label{2-two-records-the-campaign-left-open-closed-2026-09-23}

\textbf{\texttt{EXHALE\_RESIDUAL} now writes the configuration it measured under.} The
route stopped at \texttt{src/EXHALE\_main.f90:1335}, before \texttt{write\_setup\_report} and
\texttt{write\_resolved\_config} further down, so \texttt{EXHALE\_setup.out} was left empty and
\texttt{EXHALE\_resolved.out} absent. It now writes both before its stop, as the direct
steady route already did before its own. MEASURED on the stored state of
\texttt{atomic\_scalar\_gj1132x0.01\_kzz1e9/HeH9.7}: 62 and 48 lines written where there
were none, the printed residual byte-identical (9.806145E-01, 8.723527E-06,
1.071147E+00), and the resolved file identical to a normal run\textquotesingle s but for two
added lines, the base ghost count and its electrons as solved, present because
this route evaluates the residual, and so solves the ghost, before it writes.

\textbf{The conversion now asserts the charge.} \texttt{molecular\_seed\_from\_atomic\_state}
computed the positive charge density before and after the 2 H -\textgreater{} H2 transfer
and never compared the two. It now does, cell by cell, against
\texttt{inventory\_feasible\_slack} (eight units in the last place) times the charge
density, the bound \texttt{element\_inventory} already uses for two expressions of the
same nuclei; a neutral cell must come back exactly. A breach stops the run, and
the largest change is printed in units of \texttt{eps\ *\ n\_chg}. A relative change was
not used because a neutral cell carries no charge and the ratio is not defined
there. The \texttt{molecular\_seed} suite asserts it in each of its three modes:
29 of 29 pass, and the change is 0.00 units in every mode, so in memory the
transfer conserves the charge bitwise, and the 1.5 units the M2 conversion
audit measured were its file round trip, as that audit said. The failure path
was read, not exercised.

\subsection{3. Two statements of the source that were not true, corrected (2026-09-23)}\label{3-two-statements-of-the-source-that-were-not-true-corrected-2026-09-23}

Found while checking the review of \texttt{md/PLAN\_20260923.md} against the code;
the plan\textquotesingle s revision is \texttt{md/PLAN\_20260923\_rev1.md}. Neither change moves a
number: one is a comment and one is the text of a printed line.

\textbf{The certificate printed a withdrawn reason.} Its report says why the
species rows below \texttt{cert\_regime\_wind\_r} = 1.20 are reported and do not gate,
and for the band 1.10 to 1.20 it printed that the row\textquotesingle s discretization, 4.6e-4
at the binding radius, "stands above any tolerance one could set there". The
comment above \texttt{cert\_regime\_wind\_r} said the same, that a 1e-5 gate would ask
the residual to fall below the error of its own discrete equation. That
inference was withdrawn on 2026-09-11 (\texttt{md/code\_status\_20260910.md} section
3.3): the algebraic residual of the discrete equations is not bounded below by
their truncation error, and the same routine took a wind H2 row to 2.4e-13 on
a held background. The gating decision is unchanged; its stated reason is now
the one that document gives, that no state of this code has come near 1e-5
there and that the terms of the row are conserved only to 2e-3
(\texttt{certification.f90}, the comment of the two regimes and the report block).

\textbf{The absent-carrier threshold claimed a round-off argument.} The comment on
\texttt{carrier\_absent\_fraction} (\texttt{diffusive\_photochemistry.f90}) said that 1e-20 of
the free reservoir is "four decades below the round-off of a double against
the reservoir, so nothing at or under it is a resolved density at all", and
that such a carrier carries nothing any reported quantity depends on. A small
term vanishing from a rounded sum does not make it unresolved: stored as its
own double, 4e-17 stays 4e-17 with a relative spacing of 1.5e-16, and its
production, loss and transport stay resolved with it. The comment now says the
threshold is a choice and not a round-off limit, and that the test reads the
abundance alone, not the production or the incoming flux that could raise it.
The threshold itself is unchanged.

Verified: the build compiles (\texttt{EXHALE.x} md5
\texttt{69233669123c368b2d58aebc132eda85}; the printed string moved it from
\texttt{67646328ad476a976a7953aba65bdffc}), and one \texttt{stationary\ evaluate} of a stored
molecular state (the pass-6 state of \texttt{kzz1e9/HeH0.55} from the 2026-09-21
alternation audit) prints the corrected block and reproduces to four digits
the verdict that audit measured on it under its own binary (carrier balance H2
9.722E-04 at cell 500, elemental transport He/H 7.353E-04 at cell 336). The line counts of both files are
unchanged. The regression matrix was not rerun: the printed text enters no
output file the matrix compares, and a comment enters no object code.

\subsection{4. Four more source comments that were not true, corrected (2026-09-23)}\label{4-four-more-source-comments-that-were-not-true-corrected-2026-09-23}

Found while checking the second review of the LHS 1140 b plan
(\texttt{md/PLAN\_20260923\_review1.md}) against the code; the plan\textquotesingle s revision is
\texttt{md/PLAN\_20260923\_rev2.md}. All four are comments in
\texttt{diffusive\_photochemistry.f90}; the line count of the file is unchanged and
the rebuilt \texttt{EXHALE.x} has the same md5, \texttt{69233669123c368b2d58aebc132eda85},
so no object code moved.

\begin{itemize}
\tightlist
\item
  \textbf{The boundary section of the module header said "OUTER: zero gradient,
  as the element operator does"}, and that a run whose carriers are not
  negligible there is flagged in \texttt{output/Oxygen\_chemistry.txt}. The carrier
  operator imposes zero diffusive, eddy and settling flux on the outer face
  (\texttt{Jf(N)\ =\ 0}), which with a settling drift is not a zero gradient, and
  \texttt{Oxygen\_chemistry.txt} is written only by oxygen-chemistry runs and flags
  no carrier. The header now states the flux condition and points to
  \texttt{carrier\_face\_coefficients}.
\item
  \textbf{\texttt{carrier\_face\_coefficients} justified that closure by the outer cells
  being collisionless}, quoting \texttt{Kn\_bulk} = 1.4-2.4 at the domain top and an
  exobase at 18.8-25.4 R\_p. Those numbers come from runs made before the
  2026-09-13 well-balanced re-solve. MEASURED with
  \texttt{src/utils/collisional\_validity.py} on the current certified states and on
  the \texttt{kzz1e9/HeH0.55} audit state: \texttt{Kn\_bulk} at the top cell is
  0.055-0.142, the exobase lies above the domain, and the largest species
  value in the critical region is 0.41-0.59. The comment now says the outer
  cells are transitional and that the zero face flux rests on the other
  reason it gives, the measured inward settling flux of a ghost-built face.
  \texttt{md/collisional\_validity.md} carries the same measurement as a dated note.
\item
  \textbf{The row-scale floor was said to be converged by construction} ("a
  density that small cannot change any observable, so a row below it IS
  converged", and "every term of such a row is round-off"). The floor is the
  rate that would move a carrier by 1e-20 of its element in one
  signal-crossing time, a scale that a row is judged against; it proves
  nothing about the row.
\item
  \textbf{The absence test was said to separate round-off from a real row}. It
  reads the abundance alone, and a carrier far below the threshold can carry
  a real, resolved production; the comment now says that the test removes
  such a cell from the carrier gate without asking whether production or
  incoming flux could raise it, and that this is a choice.
\end{itemize}

Outside the source, the comment on the 2.26 km/s overlay shift in
\texttt{LHS1140b/make\_memo\_figures.py} said the equivalent width is unaffected. The
equivalent widths are computed from unshifted arrays, so the shift enters
none of them; applied to a model it would move the fixed-window EW by
1.2e-5 relative (MEASURED). The comment now says both.

Verified: \texttt{make} rebuilt \texttt{EXHALE.x} to the same md5; the Python file parses.
Nothing was run beyond that, since no executable statement changed.

\subsection{5. One equivalent-width operator for the observation and every model (2026-09-23)}\label{5-one-equivalent-width-operator-for-the-observation-and-every-model-2026-09-23}

The user\textquotesingle s decision on \texttt{md/PLAN\_20260923\_rev2.md} section 4: the observed
statistic stays as it is (1.1084293096 +/- 0.0294774727 \%A, working target
1.108 +/- 0.030), and every model is put through the same operator. A model
spectrum, already convolved to the instrument resolution, is interpolated
linearly onto the observation\textquotesingle s own vacuum samples inside 10832.60 to
10834.20 A (44 samples, 10832.622 to 10834.175 A) and integrated with their
trapezoid weights. Before, each model was integrated with a masked trapezoid
on its own grid inside the nominal window, so the two integrals started and
stopped at different wavelengths.

\textbf{One module, six readers.} \texttt{LHS1140b/he10830\_equivalent\_width.py} holds
the window, the vacuum/air ratio, the observation, and the operator
(\texttt{model\_equivalent\_width}, \texttt{transit\_file\_equivalent\_width},
\texttt{observed\_equivalent\_width}). The six places that computed an EW each on its
own now call it: \texttt{LHS1140b/models/status.py} (the results table and, since
this change, the EW column of the provenance table, which it used to copy
from each run\textquotesingle s \texttt{REPRODUCE.md}), \texttt{LHS1140b/make\_memo\_figures.py} (the
unbroadened and the broadened curves, and the observed value, whose variable
and printed label no longer call the released spectrum a digitization),
\texttt{LHS1140b/models/compare\_archive.py}, \texttt{LHS1140b/models/write\_reproduce.py}
(whose line also said the window was in air; it is in vacuum), and the inline
EW of \texttt{LHS1140b/models/run\_case.sh}. The unused regular expression that read
the EW out of \texttt{REPRODUCE.md} was removed.

\textbf{What moved, MEASURED.} Over the 177 EXHALE spectra of the catalog the
change is at most 1.45e-4 \%A, relative -1.1e-3 to +0.9e-3; 31 spectra move in
the fourth decimal, 15 of the 83 certified rows of the catalog table among
them (for example the certified molecular K\_zz = 1e9 He/H 2.13 and 9.7 read
1.5604 and 2.4060 where they read 1.5603 and 2.4061), and one crossing moves
(K\_zz = 0: 3.1011 to 3.1010). Broadened curves and p-winds curves move more,
up to about 0.6 per cent, because their wings reach the window edges that the
old model integral kept and the observation\textquotesingle s samples do not (for example
the matched-kernel EW of the well-mixed He/H 0.4 solution, 1.067 to 1.065).
\texttt{LHS1140b/MODELS.md} sections 7 and 8 were regenerated by
\texttt{models/status.py\ -\/-write}; \texttt{docs/lhs1140b\_catalog\_state.tex} carries
the new EWs, a paragraph on the operator and the recomputed comparison of its
section 4 (1.374869 to 1.356745 \%A, 0.60 standard deviations), rebuilt to
seven pages and checked; the memo \texttt{docs/lhs1140b\_exhale\_vs\_pwinds.tex} and
its figures are brought to the same operator in a separate pass
(\texttt{md/lhs1140b\_ew\_operator\_memo\_changes\_20260923.md}).

\subsection{6. The stage export of the stationary alternation (2026-09-23)}\label{6-the-stage-export-of-the-stationary-alternation-2026-09-23}

\texttt{EXHALE\_STAGE\_EXPORT=1} (default off) writes, for every pass of the
alternation, the state the driver holds at four named stages:
\texttt{hydro\_entry} (before the hydrodynamic solve; from the second pass on, the
state the previous composition update handed over), \texttt{hydro\_return} (the
accepted solve after the ionization sweep that follows it), \texttt{certification}
(the joint test, with the three residual rows), and \texttt{composition\_return}
(after the element and carrier updates and the rebuild of p and T at the
fixed conserved state). Each goes to \texttt{output/stage\_state\_p\textless{}pass\textgreater{}\_\textless{}stage\textgreater{}.txt}
with every row, ghosts included: r, rho, v, p, T, heat, cool, the particle
and electron densities, H I, H II, He I, He II, He III, He(2\^{}3S), H2, H2+,
H3+, HeH+, the conserved variables, and at the certification the residual
rows. At \texttt{hydro\_entry} and at the certification the driver also asks
\texttt{conservation\_budget} for the next assembly under the stage\textquotesingle s name
(\texttt{conservation\_budget\_request\_next}, new), so the energy row\textquotesingle s terms of that
state are exported without counting assemblies, which probes and trials also
consume. It is the instrument \texttt{md/PLAN\_20260923\_rev2.md} P1.5 and P1.6 (b)
call for (the four named-stage captures of review R4).

The writer calls nothing that sets a module variable. MEASURED on one pass
from the certified \texttt{molecular\_scalar\_gj1132\_kzz1e9/HeH2.13} g0008 state: the
data rows of \texttt{Hydro\_ioniz.txt} and \texttt{Ion\_species.txt} are bitwise the same with
the export off (twice), with it on, and with the previous binary
\texttt{69233669123c368b2d58aebc132eda85}; the log differs only by the export
lines. The two exported budget files carry their stage names. Binary
\texttt{EXHALE.x} md5 \texttt{b017fe42407c980068e112a8ca3d71d3}, copied to
\texttt{LHS1140b/models/EXHALE\_b017fe42.x} with its manifest. The regression matrix
was not rerun: every new call in the driver is guarded by the key, and the
budget module gained a subroutine and a header line that are reached only on
a request, so with the key off no executed statement of any route changed.
The bitwise check above covers the stationary alternation; the matrix\textquotesingle s
cases, which march, were not run.

\subsection{7. Three defects the family C audit found, fixed, and two small ones (2026-09-23)}\label{7-three-defects-the-family-c-audit-found-fixed-and-two-small-ones-2026-09-23}

The audit of the 13-decade H2 hole in
\texttt{molecular\_photochem\_gj1132\_kzzprofile/HeH9}
(\texttt{md/lhs1140b\_family\_c\_h2\_slab\_20260923.md}, item P1.7 of
\texttt{md/PLAN\_20260923\_rev2.md}) traced it to the molecular seed and found why no
solve removed it. Each finding was checked in the source before it was
changed.

\textbf{The element relaxation threw away its own kept steps}
(\texttt{binary\_element\_diffusion.f90}, \texttt{relax\_element\_composition}). Its retry
budget counts every rejected step of the relaxation, and when it ran out the
routine returned \texttt{element\_relaxation\_failed} and restored the entry
composition, even after 24 to 31 steps had been kept. The module\textquotesingle s outcome
table, and the driver that reads it, reserve that outcome for a relaxation
that kept no step; the driver then refuses the pass and skips the carrier
relaxation (\texttt{EXHALE\_main.f90}, the \texttt{outer\_element\_update\_refused} branch).
This is how the carrier operator never ran on the family C case in four
solves, and how the 2026-09-18 rung at He/H 1.411 ended at its second pass
("took 28 steps and found no admissible 29th", READ, L34c section 4). A
relaxation that has kept steps now ends as a partial advance
(\texttt{element\_relaxation\_step\_budget}) with the kept steps handed back, and says
so once in the log; one that kept none is still refused with its entry
composition. The retry budget still counts every rejection. Test: two new
rows of \texttt{src/tests/element\_operator/} on a test knob
(\texttt{element\_refuse\_after\_accepted\_for\_test}, -1 in every production run) that
lets three steps through and refuses the rest: the relaxation reports the
partial advance, three steps, and a moved composition; with no step let
through it reports the refusal and returns the entry composition bitwise.
Every row of the suite passes.

\textbf{A sweep could rewrite a transported H2} (\texttt{ionization\_equilibrium.f90}, the
extraction of the molecular partition). With the carriers transported the
sweep is handed H2 (\texttt{x\_h2\_fixed}) and solves the rest against it; a root
meets the imposed value through its own row, but a non-root acceptance
(class 4) or a clamp onto the element budget (class 3) returned another
x(4), and the extraction wrote that back as the cell\textquotesingle s H2. MEASURED on the
seed\textquotesingle s intermediate state of the family C case: the returned H2 differed from
the handed one at exactly the 21 non-root cells, by up to 82 per cent; with
the fix it differs nowhere by more than 7.8e-16. The imposed value is now
written and the atomic hydrogen takes the remainder of the budget.

\textbf{The molecular seed read a root off cells the sweep had not solved}
(\texttt{molecular\_seed\_from\_atomic\_state.f90}, the \texttt{local} mode). The seed takes the
root of each cell\textquotesingle s own H2 row at the ions of one sweep; at cells where that
sweep adopted a non-root with every ion at zero, the root was 4e-17 of the
hydrogen, thirteen decades below the neighbouring cells, and that was the
hole. The sweep now records the acceptance class of each cell
(\texttt{ieq\_cell\_acceptance\_class}, written and never read inside the sweep), and
the seed reads no root off a class-4 or class-6 cell: such a cell takes the
root of its nearest solved neighbours, interpolated in log x2 over r, and the
smaller of that and the fit. MEASURED on the same reproducer: the 21 cells go
from 4e-17 (and from up to 6.8e-4 where the rising slab was) to a smooth
1.8e-4 to 2.1e-3 between the solved cells 428 and 450; outside them the seed
moves by at most 8.2e-4 relative, through the H2 fix above. Two further
changes the audit proposed were tried and not kept: iterating sweep and root
to a common fixed point did not converge in 30 iterations and drove the base
layer to its unshielded branch (x2 = 6.9e-4 at 974 K), because a cell\textquotesingle s
shielding column is the H2 above it and the iteration is not a local one; and
adopting the root in place of the smaller of root and fit follows from the
same iteration and was withdrawn with it. The \texttt{molecular\_seed} suite passes
29 of 29.

\textbf{Smaller.} The comment on \texttt{nH\_free} in \texttt{diffusive\_photochemistry.f90} said
it excludes the hydrogen of the ion stages and molecular ions; it is the
cell\textquotesingle s hydrogen less the HeH+ nucleus (\texttt{carrier\_element\_totals}). And
\texttt{carrier\_row\_terms.txt} was written with \texttt{ES14.7}, which drops the \texttt{E} of a
three-digit exponent (\texttt{1.0332735-100}); it is now \texttt{ES15.7E3}.

\textbf{Not changed, reported.} Two provenance attributions in the state index of
the family C case are wrong (the audit memo, D4): the certificate of g0002
belongs to a later solve, and g0001\textquotesingle s manifest names another binary than its
own file header. The index is written by the catalog tools and was not edited
by hand.

Verified: the element-operator and molecular-seed suites pass; the H2 fix
and the seed on the family C reproducers as above; the stationary route with
\texttt{EXHALE\_STAGE\_EXPORT} unchanged. \texttt{EXHALE.x} md5
\texttt{45190de9ba3752eb720a728d54d3a1e3}. The one regression case that runs the H2
carrier operator, \texttt{mol\_carrier}, was run with the binary before and after the
element and H2 fixes (section 8).

\subsection{8. The regression check of section 7, a cell writer, and corrections from the spectral and On stall tests (2026-09-24)}\label{8-the-regression-check-of-section-7-a-cell-writer-and-corrections-from-the-spectral-and-on-stall-tests-2026-09-24}

\textbf{The one regression case that runs the H2 carrier operator.} \texttt{mol\_carrier}
was run from its input with \texttt{OMP\_NUM\_THREADS=1\ EXHALE\_MAXSTEPS=12000} twice,
once with the binary before the fixes of section 7 (\texttt{EXHALE\_b017fe42.x}) and
once after (\texttt{EXHALE\_af853fe7.x}), in copies
(\texttt{LHS1140b/models/.P1/d2d3\_mol\_carrier/}). MEASURED: the two
\texttt{Hydro\_ioniz.txt} differ by at most 1.3e-11 relative and the two
\texttt{Ion\_species.txt} by at most 9.4e-13; neither run met a non-root cell or the
retry budget of the element relaxation, so the difference is the H2 now
written exactly at the imposed value where a root had met it to the solver\textquotesingle s
tolerance. Below 0.1 per cent, it needs no golden refresh.

\textbf{A cell writer on the \texttt{EXHALE\_NEWTON\_TEST} route} (\texttt{EXHALE\_main.f90}). The
route printed only the row maxima of the solver\textquotesingle s vector F(Y), so where it
differed from the \texttt{EXHALE\_RESIDUAL} route the cell was not known (ISSUES D5).
It now writes \texttt{output/newton\_residual\_profile.txt} with the columns of
\texttt{output/residual\_profile.txt} (its temperature column is 0: the route forms no
temperature array). MEASURED on \texttt{molecular\_scalar\_gj1132\_kzz1e9/HeH2.13} g0008
with \texttt{EXHALE\_RELOAD\_EQ=0}: the two routes differ in cells 1 and 2 alone, where
the base face mass flux differs by 9 per cent, and nowhere else in 500 cells;
the mass-row maxima (5.10e-09 against 1.37e-09) are that base difference.

\textbf{From the checks of the preconditioned spectrum}
(\texttt{md/lhs1140b\_spectral\_checks\_20260923.md}):

\begin{itemize}
\tightlist
\item
  the cost line of the additivity measurement printed the run\textquotesingle s cumulative
  count of inner cell solves where it meant this measurement\textquotesingle s own
  (\texttt{nt\_calls}, now \texttt{nt\_calls\_here}, \texttt{steady\_newton.f90});
\item
  the comment that the Newton model\textquotesingle s products reuse the frozen WENO3
  weights is true on the line-search route and not on the trust-region route,
  where the merit evaluation sets \texttt{weno\_mode\ =\ 0} and nothing restores it
  before the leg and the spectral hook; the comment now says so, and which of
  the two the leg should use is left open;
\item
  the base boundary\textquotesingle s comment that the states "sit decades from" the
  thresholds of the wind window now records a state that sits one decade away
  and a probe that crosses the switch (\texttt{base\_boundary.f90});
\item
  the records that called the spectrum\textquotesingle s operator "shifted by the pseudo-time
  diagonal" (ISSUES D8, stage 3 section 1, the p2\_base note, the plan rev2
  section 3.6) were corrected: the shift is in the band preconditioner only.
\end{itemize}

\textbf{From the On stall test} (\texttt{md/lhs1140b\_on\_stall\_end\_to\_end\_20260923.md}):

\begin{itemize}
\tightlist
\item
  \texttt{Ionization\ transport:\ True} was refused with \texttt{Coupled\ carrier\ solve:\ True}
  but not with \texttt{On\ stall}, which hands the same registry the same unknowns
  once the alternation stalls; it is now refused with either
  (\texttt{input\_read.f90});
\item
  the comment on \texttt{EXHALE\_JFNK\_MAXIT} said a named cap makes the run take one
  stationary solve; it forbids a second entry into the stationary route and
  caps the solve of every pass inside the one entry, which the test measured
  (\texttt{EXHALE\_main.f90}); the plan rev2 carried the same error and is corrected.
  The handover itself fails: the conditions are met and the transition begins,
  and the re-evaluation of the snapshot\textquotesingle s residual stops the run on the face
  mass flux guard, the carrier balance also reading a background the snapshot
  does not hold. That is a change of the solver\textquotesingle s control flow and is reported
  for a decision, not fixed.
\end{itemize}

\textbf{The update log.} The entries of 2026-09-23 were moved from the stage 2 log
(its sections 18 to 24) into this file on the user\textquotesingle s instruction, and
\texttt{src/utils/update\_log\_to\_tex.py} now takes the stage as its argument (\texttt{2} or
\texttt{3}; the code-size appendix stays with stage 2). \texttt{docs/Update\_EXHALE\_stage2.tex}
was regenerated for sections 1 to 17 (310 pages, no error or undefined
reference) and this log\textquotesingle s twin built (7 pages at the time).

Verified: the element-operator and molecular-seed suites pass; the builds
compile. \texttt{EXHALE.x} md5 \texttt{46d5b9f4c0a62c1ab0cdf9bb165839ef}.

\begin{center}\rule{0.5\linewidth}{0.5pt}\end{center}

\subsection{9. Three defects the transient measurements found, repaired (2026-09-24)}\label{9-three-defects-the-transient-measurements-found-repaired-2026-09-24}

The P4 measurements of the plan (\texttt{md/lhs1140b\_p4\_transients\_20260923.md},
whose worker reported them and left them in place) found three defects; each
is repaired and its path checked.

\textbf{The default CFL number was a single-precision literal.} \texttt{real*8\ ::\ CFL\ =\ 0.6} stored 0.6000000238 (\texttt{parameters.f90}); every run that took the default
marched at a step 3.97e-8 longer than the one it declared (MEASURED: the first
step on a certified molecular state is 0.54736376538957865 s with the default
and 0.54736374363929763 s with \texttt{CFL:\ 0.6}). Now \texttt{0.6d0}. The same kind of
literal set the Voronov (1997) collisional-ionization coefficients of all
seventeen rows of \texttt{Cool\_coeff.f90} (11.26 stored as 11.2600002), and the
physics probe \texttt{voronov\_carbon\_ionization.f90} mirrored it on purpose to
compare exactly; both are now double-precision literals, and the probe passes
at its 1e-12 tolerance (MEASURED). The rounding was below 1e-7 relative, four
decades under the precision of the fits; a declared constant is still the
constant it declares.

\textbf{\texttt{EXHALE\_UPDATE\_MAP} could stop before its last requested step.} The
diagnostic switched off the flux-spread stop, the residual gate and the JFNK
hand-off, but not the step-change stop or the stall stop, and on a stationary
state a step of a quarter of the CFL step already moves the state by less
than \texttt{dtu\_th}: a six-step map ended after three blocks. While the map runs,
every stop flag is now cleared after the stop tests of each step
(\texttt{EXHALE\_main.f90}), so the map ends on its own step count and the staged
secondary ionization cannot change the physics between two of its steps; the
stop line says "the update map wrote its N requested step(s)". MEASURED on a
copy of the failing run without \texttt{Force\ start}: six blocks, dt from 6.6296e-05
to 6.6296e-11 of the code unit.

\textbf{\texttt{EXHALE\_DUMP\_IC=1} wrote placeholders as the lower ghost rows.} \texttt{load\_IC}
fills the primitive ghost rows with cell 1\textquotesingle s state until the first \texttt{Apply\_BC},
which writes the ghosts of the conserved array only; the dump wrote the
primitive array. It now takes the ghost rows of the primitive array the
boundary derivation held (\texttt{bc\_W\_hold}). MEASURED: the two lower ghost rows
read rho = 1.302e14 and 1.242e14, T = 233.0998 and 226.0275 K, against the
placeholder 4.857e13 and 546.06 K before, and the file\textquotesingle s 233.1004 and 226.0281
K (the 2.5e-6 is the ghost composition the boundary re-derives from the loaded
column).

The P4 memo now carries its command record (section 7, which its worker had
not written) and the repairs beside each defect.

\subsection{10. The elemental flux records ride on the Riemann face mass flux; a row writer for the helium element (2026-09-24)}\label{10-the-elemental-flux-records-ride-on-the-riemann-face-mass-flux-a-row-writer-for-the-helium-element-2026-09-24}

\textbf{What was wrong.} The element transport operator is handed the face mass
flux of the Riemann solve (\texttt{face\_mass\_flux\_of\_state}), the one the mass row
differences. Two records of it were not: the base budget entry
(\texttt{element\_base\_flux}) and \texttt{output/element\_flux\_profile.txt} built their face
mass flux as the face mean of the cell-centred density and velocity. Near the
base the cell-centred velocity carries a collocated two-cell odd-even mode
(project record P44), so on the certified \texttt{molecular\_scalar\_gj1132\_kzz1e9/HeH0.55}
state the file read 4 pi r\^{}2 F = -1.07e9, +1.16e8, -3.8e7 and +5.9e7 g/s at
faces 1 to 4, while the Riemann flux is 3.934813e7 g/s through every face
(MEASURED). The "standing base sound wave" the comments named was that mode,
not a wave of the solution.

\textbf{What the records are now} (\texttt{binary\_element\_diffusion.f90},
\texttt{certification.f90}, \texttt{EXHALE\_main.f90}):

\begin{itemize}
\tightlist
\item
  \texttt{write\_element\_flux\_profile} takes the face mass flux of the state it
  describes and is called by the certification of that state, with the flux
  the element row was measured on, whenever \texttt{EXHALE\_DIFFUSION\_CHECK=1} or a
  lower-atmosphere profile is in use; the element step no longer writes it.
  Before, the file held the profile of the last relaxation, taken before the
  last hydrodynamic solve, and a stationary evaluate never wrote it.
\item
  The base budget entry uses the flux of its own evaluation: the relaxation\textquotesingle s
  \texttt{Frho\_in}, or on the marching path the step flux \texttt{F0/6\ +\ F1/6\ +\ 2\ F2/3} of
  the three SSP-RK3 stages (the flux whose divergence is the step\textquotesingle s density
  change), read once per step; with neither, it says "not measured".
\item
  \texttt{face\_mass\_flux\_cgs} converts a code face flux with n0 mu v0 and the face
  mean of msum.
\item
  The velocity is no longer an argument of \texttt{element\_diffusion\_step} or
  \texttt{relax\_element\_composition} (nothing read it).
\item
  \texttt{composition\_residual} and \texttt{trace\_composition\_residual} reduced their norm
  with \texttt{max}, whose result with a NaN is compiler dependent; a non-finite row
  now sets the norm to huge.
\item
  \texttt{element\_transport\_residual}, documented as side-effect free, added its
  face reconstructions to the run\textquotesingle s face counters; it now restores them.
\end{itemize}

\textbf{The helium row, term by term.} \texttt{EXHALE\_ELEMENT\_ROW\_TERMS=1} (default off)
writes \texttt{output/element\_row\_terms.txt} from the certification\textquotesingle s own evaluation:
the diffusive and advective terms of each face, the residual, the floor, the
scale and the measure, with the face F\_rho, X\_face and J. Its measure agrees
with the certification\textquotesingle s full-precision anchor to the sixteenth digit.

\textbf{What it measured} (\texttt{md/lhs1140b\_element\_row\_terms\_20260924.md}): on the
certified He/H 0.55 and 1.8 states the helium flux through r \textless{} 1.10 varies by
2.2e-3 and 7.8e-5 of its median, nearly all of it the step across cell 2 (the
first solved cell, row 6.1e-5 and 2.3e-5, ungated); every other cell below
r = 1.2 stands at or below 1.5e-7 and 2.7e-7. At cell 2 the residual would change
the helium density in 9e9 and 3e10 s, against an advective cell time of 2e6 s
and a diffusive one of 47 s. The comment in \texttt{certification.f90} that gave
"element fluxes not conserved, spread 9.8" as the reason for not gating
r \textless{} 1.10 had measured the face-mean artifact; it and the printed line now say
the band is reported by decision, and the gating itself is unchanged. The
Phase-E statement that the LHS 1140 b overlap window is structurally empty
rested on the same quantity and is marked stale where it stands
(\texttt{md/phase\_e\_flux\_closure\_design.md}, \texttt{src/utils/element\_flux\_closure.py},
\texttt{md/certification\_tolerance\_anchoring\_20260910.md}); it was not re-measured.
\texttt{src/utils/collisional\_validity.py} and \texttt{md/collisional\_validity.md} used
the same wording for the base cells; corrected (the computation, whose
velocity scale binds four decades below its threshold there, is unchanged).

\textbf{Checks.} The solved state does not move: an evaluate of the certified He/H
1.45 state by the new binary writes \texttt{Hydro\_ioniz.txt}, \texttt{Ion\_species.txt} and
both \texttt{\_adv} files identical to the 46d5b9f4 evaluate apart from the
provenance line, CERTIFIED, and its \texttt{element\_flux\_profile.txt} carries
7.0353671e7 g/s from face 1 outward (MEASURED; the worker measured the same on
the He/H 0.55 and 1.8 states). Suites: element operator 65 of 65 (its base
record rows assert the new definition and fail against the old objects by
1.54 and 1.31), \texttt{diffusion\_tests.x} 35 of 35, certification 123 PASS,
steady species rows 202, species face flux 18 (the last two by the worker).
The dead velocity arrays of \texttt{diffusion\_tests.f90} and
\texttt{certification\_contexts.f90} are removed, and the test T6, which claimed a
uniform mixture "under an arbitrary velocity field" that the operator never
read, now says what it tests.

\subsection{11. The molecular K\_zz = 1e9 ladder certified across the observed equivalent width (2026-09-24)}\label{11-the-molecular-k_zz--1e9-ladder-certified-across-the-observed-equivalent-width-2026-09-24}

Items P1.1 to P1.3 of the plan, run under the binary \texttt{EXHALE\_46d5b9f4.x}
(the tree of section 8), every rung a stationary alternation from a certified
parent mapped onto the new reservoir with the pressure held
(\texttt{map\_state\_to\_grid.py\ -\/-ic\ -\/-reservoir\ He/H}, the base mixing ratio from the
same nucleus fraction f), 24 passes and a 3 h ceiling allowed. Run
directories under \texttt{LHS1140b/models/.P1/}; driver \texttt{.P1/make\_rung.sh}.

{\def\LTcaptype{none} % do not increment counter
{\small
\begin{longtable}[]{@{}>{\raggedright\arraybackslash}p{\dimexpr(\linewidth-14\tabcolsep)/7\relax}>{\raggedright\arraybackslash}p{\dimexpr(\linewidth-14\tabcolsep)/7\relax}>{\raggedright\arraybackslash}p{\dimexpr(\linewidth-14\tabcolsep)/7\relax}>{\raggedright\arraybackslash}p{\dimexpr(\linewidth-14\tabcolsep)/7\relax}>{\raggedright\arraybackslash}p{\dimexpr(\linewidth-14\tabcolsep)/7\relax}>{\raggedright\arraybackslash}p{\dimexpr(\linewidth-14\tabcolsep)/7\relax}>{\raggedright\arraybackslash}p{\dimexpr(\linewidth-14\tabcolsep)/7\relax}@{}}
\toprule\noalign{}
He/H & parent & passes & wall of the passes & outcome & EW {[}\%A{]} & log10 Mdot \\
\midrule\noalign{}
\endhead
\bottomrule\noalign{}
\endlastfoot
0.55 & the 2026-09-21 audit state (pilot) & 12, then 4 & 4.30 h, then 0.07 h & certified in the continuation & 0.1015 & 7.59 \\
1.8 & certified 2.13 (descent pilot) & 12, then 3 & 1.04 h, then 0.04 h & certified in the continuation & 1.3584 & 7.88 \\
1.45 & certified 1.8 & 17 & 1.15 h & certified & 1.0964 & 7.85 \\
1.6 & certified 1.8 & 21 & 1.22 h & certified & 1.2155 & 7.86 \\
1.52 & certified 1.45 & 14 & 0.88 h & certified & 1.1532 & 7.86 \\
0.689 & certified 0.55 & 24 (budget) & 2.85 h & not certified: elemental row 9.9e-4, still falling & 0.274 (uncertified state) & 7.67 \\
\end{longtable}
}
}

MEASURED. Every certified state was re-entered four times in a chain of
\texttt{Restart\ intent:\ stationary\ evaluate}, each entry reading the state the one
before wrote: CERTIFIED at every entry, the largest relative change of any
field after four entries 1.2e-12 (the 1.8 chain; 3.5e-13 or less for the others). The EW is the catalog\textquotesingle s measurement
operator (section 5) on the spectrum \texttt{EXHALE\_transit.py} synthesized from the
published state. The two adjacent certified rungs 1.45 and 1.52 straddle the
observed 1.108 +/- 0.030 \%A, and \texttt{models/status.py} puts the crossing of the
log-log chord at He/H = 1.4643. The interval is 0.0205 dex, below the plan\textquotesingle s
0.03 dex floor, so the refinement stopped there.

\textbf{Connectivity} (\texttt{.P1/connectivity/compare\_rungs.py}, figure
\texttt{docs/figures/lhs1140b\_kzz1e9\_molecular\_rungs.pdf}): from 1.45 to 2.13 the
temperature, velocity, density, H II and He 2\^{}3S profiles, the mass-loss rate
(7.847 to 7.902), the outer Mach number (0.515 to 0.508) and the H2 front
(x(H2) = 1e-2 at 1.200 to 1.161 R\_p) move monotonically, and neighbor
differences scale with the step (about 3 per cent in T for 0.02 dex). The
0.55 state is a different wind: x(H2) above 0.5 to 1.13 R\_p and near 1e-2 in
the far wind, a peak temperature of 3400 K against 5200 K and more, He 2\^{}3S
thirty times lower. The interval between 0.55 and 1.45 is not sampled
(the 0.689 rung is its first point), so whether the two are one branch, and
whether a second crossing lies between them, is not established. The
0.689 rung\textquotesingle s last state (uncertified) is intermediate: x(H2) above 0.5 to
1.10 R\_p as at 0.55, but 1.5e-4 at 10 R\_p, near the 1.45 value.

\textbf{The line shape} (\texttt{docs/figures/lhs1140b\_kzz1e9\_molecular\_ew.pdf}): the
crossing is of the integrated EW alone. At He/H = 1.45 the model red line is
4.07 per cent deep, 3.3 times the depth Cherubim et al. (2026) fit
(1.24 +0.22/-0.23 per cent), and correspondingly narrower; the molecular
solutions share the line-shape disagreement of the unbroadened atomic
spectra.

\textbf{Published} (\texttt{models/run\_case.sh} with \texttt{SEED=} the certified run, which
solves nothing new from a certified entry, evaluates, synthesizes the spectrum
and writes \texttt{REPRODUCE.md}): \texttt{molecular\_scalar\_gj1132\_kzz1e9/HeH0.55} g0007 and
the new cases \texttt{HeH1.45}, \texttt{HeH1.52}, \texttt{HeH1.6}, \texttt{HeH1.8} (g0002 each), added to
\texttt{models/make\_models.py}. The new cases\textquotesingle{} \texttt{input.inp} is the certified
\texttt{HeH2.13} input with the He/H changed, the input the states were solved under;
\texttt{make\_models.py} itself would write \texttt{Restart\ intent:\ stationary\ equilibrate}
for a molecular case, which the older \texttt{HeH0.083} carries and the certified
cases do not (a difference of route, reported for a decision). The catalog
inventory reads 157 state indexes, 100 certified. \texttt{LHS1140b/MODELS.md}
sections 7 and 8 and \texttt{models/CASE\_INVENTORY.\{md,json\}} were rewritten by
\texttt{status.py\ -\/-write}; \texttt{docs/lhs1140b\_catalog\_state.tex} carries the new rows
(section 13) and \texttt{docs/lhs1140b\_exhale\_vs\_pwinds.tex} its corrected
molecular paragraph and campaign table (both rebuilt, no error, pages checked).
A verification binary copied into \texttt{models/} had become the "newest
EXHALE\_.x" that \texttt{status.py} takes as the binary of record; it was moved to
the run directory that used it.

\textbf{The physical record of the straddling states} (plan section 3.1, with the
writer of section 10, \texttt{.P1/element\_row\_terms/record\_straddle.txt}): the H2
carrier row is worst inside the gated wind (5.16e-6 and 7.90e-6), and the
helium element row is worst at cell 2, ungated (2.04e-5 and 1.27e-4), the
next worst below r = 1.2 being 1.07e-6 (cell 6) and 5.4e-6 (cell 3), both in
the first cells above the base; at cell 2 the residual would change the helium density in 2.9e10 and
4.6e9 s, against an advective cell time of 1.5e6 s, and the region below 1.2
passes the helium flux on to within 9.3e-5 of its throughput.

\textbf{P3} (\texttt{md/lhs1140b\_p3\_low\_xuv\_20260923.md}, run 2026-09-23 under
\texttt{EXHALE\_69233669.x}): the atomic He/H 9.7 branch certifies at 0.0675, 0.0662
and 0.065 of the fiducial XUV (EW 0.0239, 0.0172, 0.0132 \%A) and refuses at
0.06 and 0.0624 on the elemental row; each certified rung re-entered four
times, CERTIFIED throughout.

\subsection{12. The regression of sections 9 and 10, and the validation of the crossing (2026-09-24)}\label{12-the-regression-of-sections-9-and-10-and-the-validation-of-the-crossing-2026-09-24}

\textbf{The matrix} (\texttt{make\ check}, binary \texttt{ac232b34}, the tree with section 9;
sixteen cases run, against the goldens of 2026-09-20). Verdict REGRESSION
FAIL, for three reasons, MEASURED and attributed with the binary of the day
before (\texttt{EXHALE\_46d5b9f4.x}, the tree without section 9) run on copies of the
cases:

\begin{enumerate}
\def\labelenumi{\arabic{enumi}.}
\tightlist
\item
  \textbf{The CFL literal of section 9.} The control reproduces the goldens
  (\texttt{mol\_lyman\_werner} identical, \texttt{mol\_carrier} within 1.3e-11), so every
  movement is today\textquotesingle s, and in these marching cases the only change that acts
  is the step 3.97e-8 shorter. The state files move by at most 9.1e-5
  (\texttt{wasp\_he23off}; 1.6e-5 or less in the rest), below the 1e-3 tolerance.
  What exceeds it is the \texttt{adv\_mass\_row} column of \texttt{Hydro\_ioniz\_adv.txt}, the
  mass-row measure the post-processing writes per row, in the four
  molecular 12000-step snapshots that are far from stationary (du near 0.8):
  1.2e-3 to 4.3e-3 relative (\texttt{mol\_lyman\_werner}, \texttt{mol\_diffusion},
  \texttt{mol\_sec\_ion}, \texttt{mol\_carrier}), and 0 against 5.6e-15 in one row of
  \texttt{wasp\_full\_newton}. \texttt{hydrostatic\_column} moves one ghost-row ion density
  from 4.1e-185 to 1.2e-170 (the control is identical to its golden). None
  of these is a change of a physical quantity above the tolerance.
\item
  \textbf{\texttt{oxygen\_chemistry} stops before its first step}, with the control as
  well: "ghost composition is not a fixed point", the map applied 20 times
  and still moving by 5.7e-7 against an accuracy of 1e-11, after the
  constrained molecular roots of the cold start were clamped onto the
  simplex. The golden (2026-09-20 13:53) predates the ghost fixed point
  (commit of 2026-09-20 16:18); the case has not run since. The refusal is
  the boundary\textquotesingle s designed behavior, so what to do with it is reported for a
  decision, not changed.
\item
  Two files of \texttt{iontrans\_metals} have no reference (MISS, not a failure).
\end{enumerate}

\textbf{Section 10 does not move them.} The five cases that exercise the changed
paths (element operator and its records, the certification\textquotesingle s element block)
were re-run with the tree of section 10 (\texttt{76c3ce17}) against the outputs of
the matrix: \texttt{mol\_diffusion}, \texttt{lower\_profile}, \texttt{iontrans\_metals},
\texttt{mol\_carrier}, \texttt{hydrostatic\_column}, REGRESSION PASS, all byte-identical
(tolerance 0). The other cases run without element diffusion and do not
reach those paths.

\textbf{The goldens were not refreshed.} The refresh of the sixteen cases that
section 9 moves was refused by the session\textquotesingle s permission check; the outputs of
the matrix stand in the case directories and are what a refresh would
snapshot.

\textbf{The validation of the crossing of section 11} (plan order 12):

\begin{itemize}
\tightlist
\item
  A reverse continuation: He/H 1.52 reached downward from the certified 1.6
  certifies (pass 13) and agrees with the 1.52 reached upward from 1.45 to
  5.3e-6 in every hydrodynamic field and 1.4e-4 in every species, EW 1.15317
  against 1.15316 \%A (\texttt{LHS1140b/models/.P1/reverse\_HeH1.52\_from1.6/}).
\item
  A refined grid (\texttt{LHS1140b/models/.P1/refine\_grid/}): the two straddling
  states on 1000 cells (the base block and every other spacing halved, the
  outer radius unchanged) certify (passes 16, 15) and give EW 1.0640 and
  1.1196 \%A (against 1.0964 and 1.1532 on 500 cells), log10 Mdot 7.84 and
  7.85. They still straddle 1.108; the chord moves from He/H 1.4643 to 1.505.
  The discretization shift, 3 per cent in EW, is of the size of the
  observation\textquotesingle s working error; a 2000-cell pair is running to read the order.
\item
  The physical record of section 11 and the operator of section 5.
\item
  The domain test waits for a grid option (a decision).
\end{itemize}

\subsection{13. The width of the line on the molecular crossing, and the catalog summary brought up to date (2026-09-24)}\label{13-the-width-of-the-line-on-the-molecular-crossing-and-the-catalog-summary-brought-up-to-date-2026-09-24}

\textbf{The broadening the line needs} (user instruction: after the EW, match the
width with an assumed turbulent Gaussian as \texttt{docs/lhs1140b\_exhale\_vs\_pwinds.tex}
does). \texttt{LHS1140b/models/.P1/connectivity/broaden\_molecular.py} restates the
memo\textquotesingle s operation (the excess absorption of the instrument-convolved spectrum
convolved with a Gaussian of FWHM f in velocity, the three-Gaussian extractor
of \texttt{he\_line\_metrics.py}, the EW operator of section 5, f matched where the
red-pair FWHM reaches the measured 0.841 A) and first reproduces the memo on
its atomic well-mixed He/H 0.40 state: f = 22.22 km/s, depth 1.229, blue
0.161, ratio 7.65, EW 1.0645 (the memo prints 1.065). MEASURED on the
certified molecular K\_zz = 1e9 states:

{\def\LTcaptype{none} % do not increment counter
{\small
\begin{longtable}[]{@{}>{\raggedright\arraybackslash}p{\dimexpr(\linewidth-16\tabcolsep)/8\relax}>{\raggedright\arraybackslash}p{\dimexpr(\linewidth-16\tabcolsep)/8\relax}>{\raggedright\arraybackslash}p{\dimexpr(\linewidth-16\tabcolsep)/8\relax}>{\raggedright\arraybackslash}p{\dimexpr(\linewidth-16\tabcolsep)/8\relax}>{\raggedright\arraybackslash}p{\dimexpr(\linewidth-16\tabcolsep)/8\relax}>{\raggedright\arraybackslash}p{\dimexpr(\linewidth-16\tabcolsep)/8\relax}>{\raggedright\arraybackslash}p{\dimexpr(\linewidth-16\tabcolsep)/8\relax}>{\raggedright\arraybackslash}p{\dimexpr(\linewidth-16\tabcolsep)/8\relax}@{}}
\toprule\noalign{}
He/H & cells & EW, f = 0 & matched f {[}km/s{]} & red {[}\%{]} & blue {[}\%{]} & red/blue & EW at f \\
\midrule\noalign{}
\endhead
\bottomrule\noalign{}
\endlastfoot
0.55 & 500 & 0.1015 & 22.32 & 0.113 & 0.014 & 7.89 & 0.0981 \\
1.45 & 500 & 1.0964 & 22.21 & 1.224 & 0.160 & 7.64 & 1.0600 \\
1.52 & 500 & 1.1532 & 22.20 & 1.287 & 0.169 & 7.63 & 1.1149 \\
1.6 & 500 & 1.2155 & 22.19 & 1.357 & 0.178 & 7.61 & 1.1752 \\
1.8 & 500 & 1.3584 & 22.16 & 1.516 & 0.200 & 7.58 & 1.3135 \\
2.13 & 500 & 1.5604 & 22.12 & 1.741 & 0.231 & 7.55 & 1.5087 \\
9.7 & 500 & 2.4060 & 21.80 & 2.684 & 0.359 & 7.47 & 2.3264 \\
1.45 & 1000 & 1.0640 & 22.21 & 1.188 & 0.155 & 7.65 & 1.0288 \\
1.52 & 1000 & 1.1196 & 22.19 & 1.250 & 0.164 & 7.64 & 1.0825 \\
observed & & 1.108 +/- 0.030 & & 1.24 +0.22/-0.23 & 0.25 +0.14/-0.12 & 6.7 +12.7/-3.1 & 1.108 +/- 0.030 \\
\end{longtable}
}
}

At He/H 1.45 the width is reached at 22.21 km/s (sigma 9.43 km/s) and the
depth, the blue component and the ratio then fall inside the observed 1 sigma
intervals without being tuned. The kernel does not depend on composition or
chemistry (22.12 to 22.32 over 0.55 to 2.13; the memo\textquotesingle s atomic families 22.20
to 22.40). The Gaussian wings carry 3.3 per cent of the EW out of the
window, so the broadened EW crosses the working target 1.108 at He/H 1.5112
between the certified 1.45 and 1.52 on the catalog grid; on 1000 cells the
broadened EW of 1.52 is 1.0825, and the chord extended past 1.52 reaches the
target at 1.55 (DERIVED, an extrapolation). On the advection-corrected
profiles weighted by n(2 3S) r over r, c\_s = 4.33 km/s in the line-forming
region of both straddling states, so the kernel is Mach 2.2 (sigma), 3.1
(sqrt 2 sigma, the v\_turb of Lampon et al. and p-winds) and 3.8 (sqrt 3
sigma). Figure \texttt{docs/figures/lhs1140b\_kzz1e9\_molecular\_broadened.pdf}.

\textbf{A correction to the memo found on the way.} The memo\textquotesingle s c\_s of the
line-forming region of the well-mixed He/H 0.55 state, 4.1 km/s, did not
follow from its own T = 2499 K and mu = 1.94; recomputed with the memo\textquotesingle s
weighting (n(2 3S) r integrated over r on the \texttt{\_adv} profiles, which
reproduces T = 2498 K and mu = 1.93) it is 4.21 km/s, and the Mach numbers of
the matched kernel move from 2.3, 3.3, 4.0 to 2.2, 3.2, 3.9 (summary and
section of the memo corrected, marked in place).

\textbf{The catalog summary.} Renamed on the user\textquotesingle s instruction to
\texttt{docs/lhs1140b\_catalog\_state.tex} (no date in the name; it was
\texttt{lhs1140b\_catalog\_state\_20260922.tex}, and the references in
\texttt{md/du\_stop\_vs\_stationary\_20260920.md}, \texttt{md/session\_handoff\_20260920.md}
and the stage 2 and 3 logs now point to the new name; the stage 2 twin
\texttt{Update\_EXHALE\_stage2.\{tex,pdf\}} still carries the old name and was not
regenerated). Brought up to date rather than annotated: the counts (157
indexes, 100 certified; 103 models, 91 certified), the table by group, the
table of certified models (eight new rows) and of uncertified ones (the
molecular 0.55 row out, the He/H 9.7 rungs at 0.06 and 0.0624 in), the
mechanisms (and the count of refusing rows, which the old text stated
wrongly: the energy row, not the mass row, appeared most often), the
operator-difference statement re-measured over the 91 certified rows (still
at most 1.5e-4 \%A, 15 rows in the fourth decimal), and a new section 5 on the
molecular crossing, its validation and the broadening. Rebuilt (11 pages, no
error), every changed page checked.

\subsection{14. The user\textquotesingle s decisions of 2026-09-24 on the open items, and what each gave}\label{14-the-users-decisions-of-2026-09-24-on-the-open-items-and-what-each-gave}

The list of section 6 of \texttt{md/lhs1140b\_plan\_rev2\_execution\_20260924.md},
reordered by what it changes, was put to the user; every item was decided
"as recommended". Items 2 and 3 (the rungs below He/H 1.45, and the grid of
the catalog) are still running and are recorded in the next section.

\textbf{Item 1, the width of the line.} The added Gaussian of section 13 is an
empirical parameter of the comparison, fixed by the measured width alone;
every composition read from the full line is stated as conditional on it
(\texttt{docs/lhs1140b\_catalog\_state.tex} section 5, the execution record).

\textbf{Item 4, which cells gate the species rows: compared, not changed.} A
default-off diagnostic, \texttt{EXHALE\_CERT\_GATE\_DIAGNOSTIC=1} (\texttt{certification.f90}),
prints for every carrier and element row what the verdict would be if every
cell from cell 3 outward were gated against the wind tolerance (1e-5), cells 1
and 2 reported. MEASURED on all 99 ladder cases re-evaluated with it
(\texttt{md/lhs1140b\_gating\_diagnostic\_20260924.md}): every stored verdict
reproduces; under the candidate rule 38 of the 91 certified cases stay
certified and 53 would be refused, 52 of them by the helium element row in
the first cells above the base (cell 3 in 45 cases, up to 9.1e-2) and one by
the H2 row (\texttt{molecular\ kzz1e9/HeH9.7}, 4.76e-5 at r = 1.127); no uncertified
case would become certified. In 27 atomic K\_zz cases the element row exceeds
1e-5 in 76 to 205 cells up to r = 1.09 to 1.19 (on the atomic kzz1e9 He/H 1.50
and 1.60 states 1.5e-3 and 1.8e-3 at cell 3, 184 and 182 cells out to
r = 1.16). The premise of the item, that below r = 1.2 only the first solved
cell stands above 1e-5, holds on the six certified molecular K\_zz = 1e9
states and not on the atomic ladders. The gate was not changed. The same
worker\textquotesingle s compiler warnings showed two truncations in the certification report
(\texttt{System\_HeH\_mol\_metals} printed as \texttt{System\_HeH\_mol\_metal}, the carrier row\textquotesingle s
scale text cut at 76 characters); both fields were widened (log text only).
The diagnostic shares the gate\textquotesingle s reduction with the verdict
(\texttt{species\_row\_distance\_over\_gated\_cells}); the certification suite passes 143
of 143, and with the variable unset every output is byte-identical to the
build before it (MEASURED by the worker on four states).

\textbf{Item 6, a grid option that keeps the interior.} \texttt{Outer\ shells\ {[}r\_face,cells{]}:\ \textless{}r\_face\textgreater{}\ \textless{}cells\textgreater{}} (K33d of \texttt{md/input\_schema.md}, default
off): the grid is built as before, and \texttt{\textless{}cells\textgreater{}} cells are appended beyond its
outer face, their widths in geometric progression with the ratio that puts
the last face exactly on \texttt{r\_face}; every constructed center up to the first
outer ghost, every face, width and volume and the window indices are bitwise
those without the key (new test \texttt{grid\_and\_gates/outer\_shells\_base\_identity},
40 of 40), \texttt{N} and everything sized on it follow the extended grid, and a
state written with shells carries them in its \texttt{\#\ grid} line, so it loads only
into a run with the same shells. With the key absent every output is
byte-identical to the control build (MEASURED on an LHS 1140 b evaluate and on
\texttt{hydrostatic\_column}).

\textbf{Item 5, the domain test (P1.4 (c)) and the outer closure}
(\texttt{md/lhs1140b\_outer\_shells\_domain\_test\_20260924.md}). The straddling
molecular states on the 500-cell grid plus 43 shells to 60 R\_p (spacing ratio
continuous to 4e-4) both certify (13 and 10 passes). MEASURED: the wind turns
transonic, with a sonic point at 40.3 R\_p (Mach 0.51 at the old 29 R\_p edge
becomes 0.77); the EW falls by 1.74 per cent at both compositions (1.0964 to
1.0773, 1.1532 to 1.1330), log10 Mdot moves by 2e-4, and the profiles move by
at most 0.13 per cent inside 2 R\_p and by up to 20 to 30 per cent at the old
edge. The EW crossing moves from He/H 1.4643 to 1.4886; with the matched
kernel (unchanged, 22.2 km/s) both extended states lie below 1.108. The H2
carrier\textquotesingle s last-cell defect stays two decades below the gate on both domains
(5.6e-7 at 60 R\_p): the zero diffusive outer face does not decide these
states; the subsonic outflow boundary at 30 R\_p holds the outer wind on a
subsonic branch and acts through the hydrodynamics (tentative).

\textbf{Item 7, the composition update holding the pressure.} \texttt{Composition\ update\ holds:\ energy\textbar{}pressure} (K15j, default \texttt{energy}): with \texttt{pressure}, the
composition update of the alternation keeps rho, v and p of every cell,
re-closes the eliminated species at the temperature of the new composition
and rebuilds the energy from the caloric equation of state
(\texttt{rebuild\_energy\_at\_held\_pressure}, \texttt{EXHALE\_main.f90}), printing the energy
it adds. MEASURED (\texttt{md/lhs1140b\_composition\_update\_pressure\_20260924.md}):
it removes what it was built for, the base pressure step (cells 1 to 3 held
exactly; the base face mass flux handed to the next solve 4.46e-8 against the
wind\textquotesingle s 4.44e-8, where the default hands -2.84e-5), and the next hydrodynamic
solve still fails in all three tests: on the well-mixed He/H 0.083 and 0.55
cases on the energy row at the H2 front (the carrier update raises the
cooling there by 31 to 44 per cent at the fixed wind, under either value),
and on the K\_zz = 1e9 rung 1.8 to 1.6, where the default recovers and
certifies at pass 21, it fails from pass 2 on and never recovers. The
obstacle behind the base step is the front\textquotesingle s energy balance under a
composition change at fixed wind. The default is unchanged; the option stays,
default off, with its negative result recorded. With the key absent every
output is byte-identical to the control build (MEASURED on the P1.6 entry and
on \texttt{hydrostatic\_column}).

\textbf{Item 8, one route for every molecular catalog case.} \texttt{make\_models.py}
wrote \texttt{Restart\ intent:\ stationary\ equilibrate} for a molecular case; the
certified molecular cases carry \texttt{stationary}, and were solved under it from
an atomic-to-molecular conversion as well as from molecular continuations
(He/H 2.13 and 9.7, runs of 2026-09-19, READ). So instead of a rule by the
kind of seed (the recommendation), every molecular case now takes
\texttt{stationary}: \texttt{make\_models.py} writes it, with \texttt{Coupled\ carrier\ solve:\ False}
and the base-grid comment line the certified cases carry, and its output now
equals the \texttt{input.inp} of every molecular K\_zz = 1e9 and well-mixed case
byte for byte (the two photochemical molecular cases differ only in the order
of one line); \texttt{molecular\_scalar\_gj1132\_kzz1e9/HeH0.083/input.inp}, the one
case with \texttt{equilibrate}, was changed to match.

\textbf{Item 9, the \texttt{On\ stall} handover:} deferred.

\textbf{Item 10, the goldens.} Refreshed for the sixteen cases the matrix of
section 12 ran (all but \texttt{oxygen\_chemistry}), from the outputs of that matrix;
the previous goldens are kept in \texttt{backup/regression/golden\_pre\_cfl\_20260924/}.

\textbf{Item 11, \texttt{oxygen\_chemistry}.} The case now starts from its own golden of
2026-09-20 (the 1000-step snapshot of the cold start, kept in its \texttt{IC/}) and
marches 1000 steps from there (\texttt{Load\ IC?\ True}, \texttt{Restart\ intent:\ relaxation}; its \texttt{README.md} records why); the load passes the ghost fixed
point and the run ends at du = 20.9, log10 Mdot 9.39 (MEASURED with the
control build). Its golden was snapshotted from that run.

\textbf{The smaller items.} The Phase-E overlap window is entry (AH) of
\texttt{md/TO\_BE\_DONE.md}, deferred to the next time the LHS 1140 b profile runs
are handled. \texttt{docs/Update\_EXHALE\_stage2.\{tex,pdf\}} was regenerated for the
edit of its section 17 (310 pages, no error). \texttt{src/utils/map\_state\_to\_grid.py}
wrote \texttt{nan} as the center shift in its \texttt{\#\ mapped:} line between grids of
different cell counts; it now says the shift is not defined there.

\subsection{15. The grid of the crossing (decision item 3) and the regression of sections 13 and 14 (2026-09-24)}\label{15-the-grid-of-the-crossing-decision-item-3-and-the-regression-of-sections-13-and-14-2026-09-24}

\textbf{Three grids, three compositions.} The certified molecular K\_zz = 1e9
states at He/H 1.45, 1.52 and 1.6 were mapped onto grids of 1000 and 2000
cells (the base block and every other spacing halved at each step, the outer
radius unchanged; target grids from cold-start dumps) and solved again with
the rung environment; all six certify (1000 cells: 16, 15, 15 passes; 2000
cells: 13, 13, 19). MEASURED through the catalog\textquotesingle s evaluation, synthesis and
EW operator (\texttt{LHS1140b/models/.P1/refine\_grid/},
\texttt{.P1/connectivity/grid\_convergence.txt}):

{\def\LTcaptype{none} % do not increment counter
{\small
\begin{longtable}[]{@{}>{\raggedright\arraybackslash}p{\dimexpr(\linewidth-12\tabcolsep)/6\relax}>{\raggedright\arraybackslash}p{\dimexpr(\linewidth-12\tabcolsep)/6\relax}>{\raggedright\arraybackslash}p{\dimexpr(\linewidth-12\tabcolsep)/6\relax}>{\raggedright\arraybackslash}p{\dimexpr(\linewidth-12\tabcolsep)/6\relax}>{\raggedright\arraybackslash}p{\dimexpr(\linewidth-12\tabcolsep)/6\relax}>{\raggedright\arraybackslash}p{\dimexpr(\linewidth-12\tabcolsep)/6\relax}@{}}
\toprule\noalign{}
He/H & EW, 500 cells & 1000 cells & 2000 cells & continuum (DERIVED) & log10 Mdot 500 / 1000 / 2000 \\
\midrule\noalign{}
\endhead
\bottomrule\noalign{}
\endlastfoot
1.45 & 1.0964 & 1.0640 & 1.0475 & 1.030 & 7.85 / 7.84 / 7.84 \\
1.52 & 1.1532 & 1.1196 & 1.1025 & 1.085 & 7.86 / 7.85 / 7.85 \\
1.6 & 1.2155 & 1.1804 & 1.1625 & 1.144 & 7.86 / 7.86 / 7.86 \\
\end{longtable}
}
}

The EW converges at first order (successive-difference ratio 1.96, p = 0.97,
at every composition), and the continuum column is the Richardson
extrapolation at that order. The crossing of the working target 1.108 moves
from He/H 1.4643 (500 cells) to 1.5054 (1000), 1.5273 (2000) and 1.552
(continuum); from the broadened line of section 13, from 1.5112 to 1.5547,
1.5780 and about 1.60 (the last just above the 1.6 rung); every other value is
interpolated between certified rungs. The matched kernel does not move with
the resolution (22.18 to 22.21 km/s on all nine states). The 500-cell catalog
overestimates the EW by about 6 per cent, twice the observation\textquotesingle s working
error; a formally second-order scheme converging at first order in the EW
points to a first-order error in the line-forming region, whose source is not
established. Combined with the domain factor of section 14 (DERIVED, the two
measured separately), the crossing stands near 1.58 from the EW and near 1.63
from the broadened line. Which grid the catalog should use is left to the
user with these numbers.

\textbf{The regression of sections 13 and 14.} After all the code of this series
had landed (the diagnostic of item 4, the grid option of item 6, the
composition option of item 7, the widened report fields, and the mapper
record), the tree was rebuilt (\texttt{EXHALE.x} md5
\texttt{2df491cb020e8200be247e2469964a63}, no warning) and the whole matrix run
against the goldens refreshed in section 14: REGRESSION PASS, all seventeen
cases byte-identical (66 files; the two \texttt{\_adv} files of \texttt{iontrans\_metals} have
no reference, as that route writes none), \texttt{oxygen\_chemistry} included on its
warm start.

\subsection{16. The user\textquotesingle s decisions of 2026-09-24 on the four remaining items}\label{16-the-users-decisions-of-2026-09-24-on-the-four-remaining-items}

Put to the user after sections 14 and 15; decided the same day.

\textbf{The catalog\textquotesingle s grid: 500 cells, the compositions quoted from the grid
continuum.} The catalog keeps its 500-cell grid (every certified state, the
chains and the ladders stay as they are); a composition read from it is
quoted as the continuum value of section 15, with the difference between the
500-cell and the continuum value as the grid uncertainty of every catalog
number: from the EW, He/H 1.55 (the 500-cell crossing is 1.46, 0.025 dex
lower, because the 500-cell EW is about 6 per cent high), from the broadened
line about 1.60 (1.51 on 500 cells); with the domain factor of section 14, 1.58 and 1.63 (DERIVED, the two
effects measured separately). \texttt{docs/lhs1140b\_catalog\_state.tex} section 5
states the rule.

\textbf{The gating candidate: not adopted.} The species rows keep their current
gate; the diagnostic \texttt{EXHALE\_CERT\_GATE\_DIAGNOSTIC=1} stays as a default-off
report.

\textbf{The pressure-holding composition update: kept, default off.} \texttt{Composition\ update\ holds:\ pressure} stays in the code with its negative result (section
14, item 7); the default remains \texttt{energy}.

\textbf{The energy imbalance at the H2 front: explored.} The obstacle item 7
exposed, the energy row at the H2 front failing after a composition change at
fixed wind, is being diagnosed (\texttt{LHS1140b/models/.P1/h2\_front\_energy/}); its
result is recorded in a later section.

\subsection{17. The atomic-data audit of 2026-09-24, applied}\label{17-the-atomic-data-audit-of-2026-09-24-applied}

\texttt{md/EXHALE\_ATOMIC\_DATA\_AUDIT\_2026-09-24.md} (an audit of the atomic data
against the sibling code MoCHII) was checked item by item in the code and,
where it could be obtained, in the original source, and applied; the four
memos carry every verdict, source, size and command:
\texttt{md/atomic\_data\_audit\_applied\_core\_20260924.md} (the H/He cooling core),
\texttt{md/atomic\_data\_audit\_metal\_tables\_20260924.md} (the CHIANTI metal
tables), \texttt{md/charge\_exchange\_detailed\_balance\_20260924.md} (the
charge-exchange table) and \texttt{md/atomic\_data\_audit\_applied\_docs\_20260924.md}
(comments, unreachable code, the detailed-balance test). Every item the
audit raised was found to be real; the corrections are physics corrections
and are on by default.

\textbf{The H/He cooling core.}

\begin{itemize}
\tightlist
\item
  Recombination cooling of every ion is k T alpha (3/2 + dln alpha/dln T)
  on the coefficient the balance itself uses (B1, B7, B12). Hui \& Gnedin
  (1997) print the He III fit with a factor 8 where the hydrogenic scaling
  gives 2: He III removed 3.05 kT per capture at 1e4 K and now 0.76 kT. The
  relation reproduces Hummer (1994) Table 1 to 0.1 per cent.
\item
  The post-process temperature root and \texttt{eval\_cool} now call one cell
  assembly, \texttt{radiative\_cooling\_of\_cell} (B2); the two coolings differed by
  a factor 1.04-1.60 on the \texttt{wasp\_full} post-process and now agree to
  1.3e-14.
\item
  \textbf{The 2\^{}3S -\textgreater{} 1\^{}1S electron-impact de-excitation (B3)}, the
  detailed-balance reverse of q13, is in every place the triplet balance is
  written, with the 19.82 eV returned to the gas. Oklopcic \& Hirata (2018)
  leave it out as small against 2\^{}3S -\textgreater{} 2\^{}1S, 2\^{}1P; that holds at 1e4 K
  (ratio 0.055) and not below: 0.15 at 5e3 K, 0.57 at 3e3 K, 2.8 at 2e3 K
  (MEASURED), the temperatures of the LHS 1140 b line-forming region.
\item
  Voronov (1997) fits take Voronov\textquotesingle s fitted dE again (B4: Mg II x1/1.47, C I
  x1/1.10 at 5e3 K); n = 2 recombination is the balance\textquotesingle s case-B coefficient
  split with Draine\textquotesingle s 2s fraction (B6, +2.6 per cent at 1e4 K); He I, He II
  (A3) and H I (B8) collisional-excitation cooling are CHIANTI 11.0.2 sums
  with the observed level energies (x1.11, x1.29 and x1.11 at 1e4, 1e4 and
  5e3 K), H I on one rate set shared with the n = 2 model; single-precision
  literals are gone from the four files (B10); \texttt{use\_2lev\_cool},
  \texttt{lambda\_2level}, \texttt{melem\_Z}, \texttt{cool\_FeII\_func} and the other unreachable
  wrappers are removed (B11).
\end{itemize}

\textbf{The metal line cooling (A2).} Mg II took the CHIANTI \texttt{.scups} scaling
energy (4.27 eV) as its transition energy where the observed g-weighted
4.4300 eV belongs, 1.15 times too high at 1e4 K and 1.39 at 5e3 K; it is
refitted with observed energies and every bound channel (0.2 per cent).
Ca II and Mg I carried the resonance line alone; they now take (T, n\_e)
statistical-equilibrium tables with every channel (at n\_e = 1e9 cm\^{}-3 the
old coefficient was 1.7-4 times low for Ca II and 5.8-480 times low for Mg I
over 1e4-3e3 K). Fe I stays on its Boltzmann (LTE-limit) table, now stated
as an upper limit at most 2.0 times the low-density value; the data for a
density-resolved Fe I solve are not available here.

\textbf{Charge exchange.} The reverse rows of Huang et al. (2023) Table 4 depart
from detailed balance with internal partition functions by factors 0.18 to
2.9. Each measured or computed direction is kept as the original paper gives
it, and the other is computed from it by detailed balance with NIST
partition functions; four reverse rows are not carried (A2, A10, C4, C6),
their forward channel ending in an excited or radiative product; three
transcriptions are corrected to the originals (A12 Si2+ + H: 1.23e-9 and
the sign of its exponent; C5: Zhao et al. 2004 eq. 19 in full; A15: Stancil
et al. 1998 Table 1 directly). Ten pairs whose originals could not be
obtained (mostly Physical Review A) stay as printed and are marked
KNOWN\_BLOCKED in the new test \texttt{src/tests/charge\_exchange\_detailed\_balance};
among them N + H+ (its endothermic direction falls with T) and Na + H+ (both
directions below 1e-60 cm\^{}3 s\^{}-1, so no Na-H charge exchange acts).

\textbf{Other corrections of the same series.} Comments that contradicted the
code or the source (C1, C4 to C9, C11) are corrected; \texttt{rk\_R21\_H\_Hep\_cx},
\texttt{rk\_R22\_Hp\_He\_cx} and the three unreachable Lyman-Werner fits are removed
after a call-graph search; the Lyman-Werner warning of
\texttt{ionization\_equilibrium.f90} and \texttt{write\_output.f90} no longer cites a DB96
fit the band share has not used since 2026-09-06 (\texttt{NH2\_db96\_max} removed);
\texttt{src/tests/a2\_m1} compiles again against a built tree and its saved output
was regenerated (the 0 K photolysis thresholds; \texttt{md/a2\_reaction\_audit.md}
carries a correction block); the \texttt{Makefile} flag stamp now hashes the first
line of \texttt{\$(FC)\ -\/-version}, whose conda-forge package build number
\texttt{-dumpfullversion} drops: the previous \texttt{build/} held 35 objects of gcc
16.2.0-4 and 92 of 16.2.0-5 linked into one binary.

\textbf{Checks.} The tree was rebuilt from \texttt{make\ distclean} (\texttt{EXHALE.x} md5
\texttt{551dae520985d63c9ea8b15c21106589}, no warning; copied to
\texttt{LHS1140b/models/EXHALE\_551dae52.x} with \texttt{BINARY\_MANIFEST\_551dae520985.txt});
\texttt{physics\_probe}, \texttt{charge\_exchange\_detailed\_balance} (65 pass, 10 known
blocked) and \texttt{charge\_exchange\_rows} pass on it. The regression matrix and
the remaining suites are recorded in the next sections.

\textbf{Effect (MEASURED by the workers, on builds of the same physics).}
\texttt{wasp\_full}: log10 Mdot -2e-4 (cooling core), +4e-4 (charge exchange), -2.7e-3
(metal tables, a du-stopped case, so below its path spread); Fe III x0.48
at 1.2 R\_p. LHS 1140 b: the He 10830 EW of the certified molecular states
falls by 14.9 to 16.7 per cent on evaluation, almost all from B3, and by 20
per cent after one stationary re-solve (He/H 1.45: 1.0964 -\textgreater{} 0.8775 \%A;
1.52: 1.1532 -\textgreater{} 0.9180), because n(2\^{}3S)/n(He) falls by 41-58 per cent at
5-10 R\_p, where n\_e q31g is 3-100 times A31. No certificate of the catalog
holds under the corrected physics. \textbf{Every LHS 1140 b number of this log
before this section, and of \texttt{docs/lhs1140b\_catalog\_state.tex}, is of the
physics before the audit.}

\subsection{18. The H2-front energy imbalance, the H3+ non-LTE interpolant, and the descent below He/H 1.45 (2026-09-24)}\label{18-the-h2-front-energy-imbalance-the-h3-non-lte-interpolant-and-the-descent-below-heh-145-2026-09-24}

\textbf{The H2 front (decision of section 16).} Diagnosed with the tree binary of
section 15 and no change to the source
(\texttt{md/lhs1140b\_h2\_front\_energy\_imbalance\_20260924.md}, figure
\texttt{docs/figures/lhs1140b\_h2\_front\_energy\_imbalance.pdf}). MEASURED: across the
composition update one term of the front\textquotesingle s energy row moves, the H3+
infrared cooling (99.6 per cent of the cooling there): on the well-mixed
He/H 0.083 update H2 at the front rises 3.2 per cent, H3+ answers with +8.6
per cent (+17.5 with the pressure held), and the cooling of the spike cell
rises by 31 (44) per cent against a heating that rises by 1 per cent. The
code path is consistent (the next solve assembles the relaxation\textquotesingle s cooling
to 1e-7 and 3e-5). About two thirds of the jump came from the interpolant of
the H3+ non-LTE factor (below). With the composition frozen a hydrodynamic
root exists and the trust-region solve reaches it (\textbar\textbar R\textbar\textbar{} 1.5e-7 and 1.6e-8),
where the line search cannot take its first step; the well-mixed He/H 0.083
alternation still does not converge with the pressure held and the trust
region (6 passes, the front walking outward, Mdot -26 per cent). The
trust region inside the alternation is a solver-control change and was not
made.

\textbf{The H3+ non-LTE factor is now a power law between the columns of Miller et
al. (2013) Table 6} (\texttt{h3p\_cooling.f90}, \texttt{s\_table}). The emission per H3+
ion can grow at most as fast as the collision rate that populates the
emitting levels, so the local index d ln s / d ln n(H2) lies between 0 and 1,
and the table\textquotesingle s node-to-node indices do (0.00 to 0.89). The interpolant
linear in s over log n gave 12.6 just above the 1e8 cm\^{}-3 column and 0.17
just below it at 2000 K (DERIVED), and over r = 1.25-1.40 R\_p of the
well-mixed He/H 0.083 state an H3+ cooling 1.8 to 5.6 times the power law\textquotesingle s.
In T the interpolant stays linear; its form there is not fixed by any bound
and is written at the code with the row ratios. The correction is on by
default and is in the binary of section 17.

\textbf{The descent below He/H 1.45 (decision item 2 of section 14)} ran under the
P1.2 rules on \texttt{EXHALE\_46d5b9f4.x} (the physics before section 17;
\texttt{LHS1140b/models/.P1/descend\_from\_1.45\_record.txt}): 1.152 from 1.45 not
certified (18 passes, 3 h); 1.2924 from 1.45 certified (16 passes); 1.152 from
1.2924 not certified (24 passes, the H2 carrier row 3.1e-4 at cell 247);
1.2202 from 1.2924 not certified (24 passes, the He/H element row 8.6e-5 at
cell 234 and falling by about 10 per cent a pass); the driver stopped on two
consecutive failed refinements. The interval 0.55 to 1.29 stays unsampled,
and the one new certified state is of the physics before section 17.

\textbf{The molecular family under the corrected physics (user decision of
2026-09-24: this family first).} The certified He/H 1.45, 1.52, 1.6, 1.8
and 2.13 states are being re-solved from themselves with
\texttt{EXHALE\_551dae52.x} (\texttt{LHS1140b/models/.P1/newphys/}, the rung environment,
24 passes, 4 h); their result, the new crossing and what follows for the
rest of the catalog are recorded in a later section.

\subsection{19. The second round of the atomic-data audit (2026-09-24/25)}\label{19-the-second-round-of-the-atomic-data-audit-2026-09-2425}

The audit\textquotesingle s author rechecked the applied tree from MoCHII
(\texttt{\textasciitilde{}/MoCHII/MoCHII\_v1.00/md/EXHALE\_AUDIT\_RECHECK\_2026-09-24.md} and
\texttt{...\_SECOND\_RECHECK\_2026-09-24.md}): every item of section 17 was confirmed as
applied correctly, and the open ones were listed. Each was verified in the
code and the original source before it was applied; memos
\texttt{md/atomic\_data\_audit\_round2\_core\_20260924.md},
\texttt{md/charge\_exchange\_detailed\_balance\_20260924.md} section 10 and
\texttt{md/chianti\_descaling\_spline\_20260924.md}.

\textbf{Recombination and the helium network.}

\begin{itemize}
\tightlist
\item
  The ground-state capture coefficient alpha\_1 is the Milne relation on the
  photoionization cross section the transfer uses, for H I, He I and He II
  (0.998 of Hummer 1994 for H I, 1.006-1.011 of Hummer \& Storey 1998 for
  He I); case B = Badnell alpha\_A minus that alpha\_1 (the Mao \& Kaastra fit,
  2.7 per cent low for He I at 1e4 K, is removed).
\item
  The Oklopcic \& Hirata (2018) alpha\_1 is the capture to 1\^{}1S alone; with the
  singlet/triplet coupling off the excited-singlet captures were missing.
  Case B is now split into triplet and singlet captures with Hummer \& Storey
  (1998) Table 5 in both configurations.
\item
  \textbf{H ground captures escape where the gas is thin} (new key \texttt{H\_rec\_escape},
  default True, K14d2 of \texttt{md/input\_schema.md}): alpha\_HII = alpha\_B + y\_HI
  alpha\_1 with the local escape weight of the He II construction, case A in
  the thin outer wind and case B in the thick base; \texttt{False} is case B
  everywhere. A correction of the case-B assumption where it does not hold,
  hence on by default.
\item
  The He III recombination cascade (ground capture, direct capture to n = 2,
  He II Ly-alpha, the two-photon continuum of Nussbaumer \& Schmutz 1984, the
  2s-2p mixing of He II) is absorbed on the spot under \texttt{He\_rec\_coupling}.
\item
  The collisional feed of 2\^{}3S through the higher triplets (feed/q13 = 0.018,
  0.101, 0.328 at 5e3, 1e4, 2e4 K, CHIANTI 11.0.2) is added to q13 wherever
  the triplet balance reads it, and so are the 2\^{}3S -\textgreater{} n\^{}1L (n \textgreater= 3)
  excitations.
\item
  A(2s -\textgreater{} 1s) = 8.2249 s\^{}-1 (Nussbaumer \& Schmutz 1984; was 8.26); the 2s-2p
  l-mixing of H by protons, He+ and He2+ from Pengelly \& Seaton (1964) with
  their own cut-off rule (the Lamb shift or fine structure for 2s, not the
  Debye radius); the 2\^{}3S electron-impact ionization as Black (1981) eq. 12
  prints it; \texttt{coll\_rate\_prefactor} = 8.629132e-6 from CODATA, one parameter;
  the 2\^{}3S triplet cooling is a CHIANTI sum (Black\textquotesingle s 13179 K was his fitted
  threshold, not an energy); the (T, n\_e) cooling tables above their top
  density continue with the last slope clipped to {[}-1, 0{]}; q31g is recorded
  in the diagnostic states (state 28 of \texttt{stage\_channels.txt}).
\end{itemize}

\textbf{Energy terms that were missing.} The metal recombination energy (k T
alpha (3/2 + dln alpha/dln T)) and the metal collisional-ionization energy
(potential times the Voronov rate) are in the cooling; the charge-exchange
energy defects heat the gas (new heat channel 20, \texttt{heat\_charge\_exchange}:
per reaction the defect less the energy radiated by an excited product,
product states from the sources read; the ledger closes to 1.4e-8 eV an
event). On \texttt{wasp\_full} the charge-exchange heat is 18 per cent of the
heating at the base, 3.9 at 1.1 R\_p, 0.4 at 1.3 R\_p; the base share is
mostly Fe + H+, whose product state could not be read (Rutherford \& Vroom
1972) and is taken as ground, an upper bound.

\textbf{Charge exchange.} N + H+ / N+ + H are now Kingdon \& Ferland (1996) Table 1
(Butler \& Dalgarno 1979) with the reverse by detailed balance (the printed
Huang rows were 69 and 1.5e7 times low at 1e4 K); Na + H+, Na+ + H and K + H+
are not carried (no reading of the printed Table 4 form gives a physical
rate; the originals, Dutta et al. 2001 and Watanabe et al. 2002, could not be
obtained); D26 and D30 are derived by detailed balance from the source
directions; the gate now also requires every carried rate to obey the
threshold law (205 pass, 6 known blocked, 2 known defects).

\textbf{CHIANTI descaling.} CHIANTI prescribes the natural cubic spline through
the \texttt{.scups} points (\texttt{descale\_scups.pro}, CHIANTI Technical Reports 4 and 13);
the scripts of \texttt{cooling\_data/} used the not-a-knot spline (as ChiantiPy).
Every CHIANTI-based coefficient was regenerated with the natural spline: He I
cooling moves by -2 to +0.5 per cent over 3e3-3e4 K (Upsilon(1\^{}1S-2\^{}1P) by
-16 per cent at 1e4 K), H I by +-0.7 per cent, O I fine-structure Upsilon by
+1.3 to +6.7 per cent, most others below 1 per cent.

\textbf{Other repairs of the series.} The Lyman-Werner band share the H2O/OH
continua see is now summed cell by cell along the path at each cell\textquotesingle s own
(T, n\_H) (it was one uniform slab at the local T: the band\textquotesingle s photon ledger
was 18 per cent short on the warm-started \texttt{oxygen\_chemistry} state; only
\texttt{Oxygen\ chemistry:\ True} runs reach it); six test suites were repaired
(\texttt{md/test\_suite\_repairs\_20260924.md}).

\textbf{Build and checks.} Tree rebuilt from \texttt{make\ distclean}: \texttt{EXHALE.x} md5
\texttt{b059a4be0459d65fdc3dff9894ba5d2f}, no warning, copied to
\texttt{LHS1140b/models/EXHALE\_b059a4be.x} with \texttt{BINARY\_MANIFEST\_b059a4be0459.txt}.
All 42 suites run on it (MEASURED): 38 pass; \texttt{grid\_and\_gates} fails its three
\texttt{stationary\_evaluate\_*} rows on the pinned fixture
\texttt{backup/regression/wasp\_full\_newton/IC/} (solved under the old physics; a
re-solve under this physics certifies, \texttt{info\ =\ 0}, \textbar\textbar R\textbar\textbar{} 2.3e-9, and waits in
\texttt{LHS1140b/models/.P1/suite\_repairs/wfn\_IC\_r2/} for the golden refresh);
\texttt{residual\_determinism} fails its known closure-spread row (1.17e4, the
solver item open since 2026-09-18); the two \texttt{steady\_*} suites need the log of
a finished steady solve as input.

\textbf{Effect (MEASURED by the workers).} Of the core items alone: \texttt{wasp\_full}
log10 Mdot 13.35 -\textgreater{} 13.32, the H ionization front outward (x(H+) 0.36 lower
at 1.27 R\_p), T up to +11 per cent; \texttt{mol\_metals} -0.004. Of the
charge-exchange items: \texttt{wasp\_full} +4.2e-3, N II down by up to 1.8 dex. On
the LHS 1140 b He/H 1.45 state the core items lower the He 10830 EW by 7 per
cent on evaluation, mostly through the triplet/singlet split of case B
(0.041 \%A), \texttt{H\_rec\_escape} (0.020) and the Milne alpha\_1 (0.011).

\textbf{The LHS 1140 b molecular family between the two rounds} (binary
\texttt{EXHALE\_551dae52.x}, \texttt{LHS1140b/models/.P1/newphys/}): re-solved from the
certified states, no rung certified in 24 passes (the H2 carrier row at the
front falling by 0.72-0.86 a pass), then all five certified in 7-11 more
passes. Their EWs (MEASURED, evaluate + synthesis + the EW operator): He/H
1.45 0.9813, 1.52 1.0185, 1.6 1.0587, 1.8 1.1500, 2.13 1.2765 \%A; the 1.108
crossing near He/H 1.70. These are of the physics of section 17 only; the
states under this section\textquotesingle s physics are being solved from them.

\subsection{20. The molecular K\_zz = 1e9 family under the audited physics (2026-09-25)}\label{20-the-molecular-k_zz--1e9-family-under-the-audited-physics-2026-09-25}

The user\textquotesingle s decision of 2026-09-24 was to re-solve this family first. The
five states He/H 1.45, 1.52, 1.6, 1.8 and 2.13 were re-solved from the
states certified under the first round (section 19) with
\texttt{EXHALE\_b059a4be.x} (\texttt{LHS1140b/models/.P1/newphys/HeH\textless{}x\textgreater{}\_r2/}, the rung
environment, \texttt{resolve.sh}): all five certified in 13 passes (the H2 carrier
row at the front falling from 0.5 to 6e-6 - 9e-6). Each was then evaluated,
synthesized and measured with the EW operator (\texttt{evaluate\_ew.sh},
\texttt{eval\_r2/}; every evaluation certified). MEASURED:

{\def\LTcaptype{none} % do not increment counter
{\small
\begin{longtable}[]{@{}>{\raggedright\arraybackslash}p{\dimexpr(\linewidth-10\tabcolsep)/5\relax}>{\raggedright\arraybackslash}p{\dimexpr(\linewidth-10\tabcolsep)/5\relax}>{\raggedright\arraybackslash}p{\dimexpr(\linewidth-10\tabcolsep)/5\relax}>{\raggedright\arraybackslash}p{\dimexpr(\linewidth-10\tabcolsep)/5\relax}>{\raggedright\arraybackslash}p{\dimexpr(\linewidth-10\tabcolsep)/5\relax}@{}}
\toprule\noalign{}
He/H & EW, audit round 2 & round 1 & before the audit & log10 Mdot, round 2 \\
\midrule\noalign{}
\endhead
\bottomrule\noalign{}
\endlastfoot
1.45 & 0.9387 & 0.9813 & 1.0964 & 7.899 \\
1.52 & 0.9776 & 1.0185 & 1.1532 & 7.906 \\
1.6 & 1.0198 & 1.0587 & 1.2155 & 7.913 \\
1.8 & 1.1159 & 1.1500 & 1.3584 & 7.927 \\
2.13 & 1.2497 & 1.2765 & 1.5604 & 7.945 \\
\end{longtable}
}
}

The 1.108 crossing of the log-log chord moves from He/H 1.4643 (before the
audit) to \textbf{1.7833, between the certified 1.6 and 1.8}, on the 500-cell
grid; the EW falls by 14 to 20 per cent at fixed He/H and log10 Mdot rises
by about 0.05. The grid continuum factor and the domain factor of sections
15 and 14 were measured under the old physics and are not re-measured here.
Figure \texttt{docs/figures/lhs1140b\_kzz1e9\_audit\_stages.pdf}
(\texttt{.P1/newphys/compare\_physics\_stages.py}): the EW against He/H at the three
stages, the line at He/H 1.8, and T, v, n(2\^{}3S), x(H+) and x(H2) at He/H 1.8.
At He/H 1.8 (MEASURED on the \texttt{\_adv} profiles, before the audit / round 1 /
round 2): the line keeps its shape and loses depth; the peak temperature is
5439 / 5315 / 5600 K; the velocity at the outer edge 1.21 / 1.40 / 1.36
km/s; x(H+) at 20 R\_p 0.059 / 0.056 / 0.050 (the ground captures escape in
the thin far wind); the metastable density at its peak (1.35-1.38 R\_p) 33.3 /
33.3 / 31.8 cm\^{}-3, and its radial integral over 1-1.5, 1.5-3 and 3-30 R\_p
11.1 / 10.9 / 10.2, 28.8 / 28.3 / 27.6 and 28.4 / 24.2 / 23.4 cm\^{}-3 R\_p. The
loss is in the cool outer wind beyond 3 R\_p (-18 per cent), where the
superelastic 2\^{}3S -\textgreater{} 1\^{}1S collisions outrun the radiative decay, and the EW
falls by 18 per cent with it.

\subsection{21. The code-size appendix moves to this log (2026-09-25)}\label{21-the-code-size-appendix-moves-to-this-log-2026-09-25}

On the user\textquotesingle s instruction the hand-written code-size appendix
\texttt{docs/Update\_EXHALE\_appendix.tex} is now input by this log\textquotesingle s typeset twin
instead of the stage 2 twin (\texttt{src/utils/update\_log\_to\_tex.py}, \texttt{STAGES});
\texttt{docs/Update\_EXHALE\_stage2.\{tex,pdf\}} were regenerated without it. Its
references to "this log" now name the stage 2 log where they meant it, and
it states that its numbers (measured 2026-09-22) do not yet include the code
of stage 3; they are re-measured with \texttt{src/utils/codesize.py} when the
atomic-data series is closed.

\subsection{22. The originals supplied on 2026-09-25, read and applied}\label{22-the-originals-supplied-on-2026-09-25-read-and-applied}

The user supplied the papers the two audit rounds could not obtain; each was
read and the code checked against it (memos
\texttt{md/charge\_exchange\_detailed\_balance\_20260924.md} sections 11 and 11.8,
\texttt{md/atomic\_sources\_hhe\_20260925.md}).

\textbf{Charge exchange.}

\begin{itemize}
\tightlist
\item
  Na + H+ (Dutta et al. 2001, non-radiative, product H(n=2), endothermic by
  1.682 eV) and Na, K + H+ (Watanabe et al. 2002, radiative): neither paper
  gives a fit and Huang et al.\textquotesingle s Table 4 forms are in neither; the rate is
  now the Maxwellian average of Dutta\textquotesingle s cross sections (their own rate
  figure violates the threshold law below about 5000 K) plus Watanabe\textquotesingle s
  radiative values: 8.4e-12 at 1e4 K against 4.4e-171 as coded before. The
  reverse is not carried (it starts from H(n=2)). K + H+ from Watanabe.
\item
  Fe + H+, O+, N+ (Rutherford \& Vroom 1972 Table I): the 300-1200 K values
  in log-log, held above 1200 K; the measured cross sections start at 8 eV,
  so every wind rate is an extrapolation, bracketed by 1.4e-9 and 4.0e-9 at
  1e4 K (DERIVED; the held value is the source\textquotesingle s own extrapolation to 6 per
  cent; Pequignot \& Aldrovandi 1986 infer 2.7e-9). The product state is not
  determined in the paper; the near-resonance reading (Fe+ z4F/z4D, 1.18 eV
  of heat an event instead of the full 5.70 eV defect) is the one Pequignot
  \& Aldrovandi\textquotesingle s Table 3 also takes; the range is stated at the code.
\item
  S + H+ (Zhao et al. 2005 Table III), Si + He+ (Satta et al. 2013 eq. 8,
  Table 1), S + C+ (Chenel et al. 2010 Table 2), Na + Ca+ (Kwolek et al.
  2019, measured at 0.015-0.09 eV, the Langevin value kept as an upper bound)
  replace the Table 4 transcriptions; the reverses of channels whose forward
  ends in an excited product (Fe rows, Si + He+, Na + Ca+) are not carried.
  C + Si+ (D1/D2) stays blocked: the paper Table 4 cites for it does not
  contain it. The detailed-balance gate: 194 pass, 0 fail, 1 known blocked.
\item
  Effect: \texttt{wasp\_full} log10 Mdot -3e-4; the charge-exchange heat at the base
  falls from 18 to 6.7 per cent of the heating; Na I D2 transit depth 1.169
  -\textgreater{} 1.156 per cent.
\end{itemize}

\textbf{Hydrogen and helium.}

\begin{itemize}
\tightlist
\item
  The He I collision strengths of the metastable network (q13, q31g, q31a,
  q31b) were exponential fits the ATES code carried, not coefficients of
  Lampon et al. (2020), who take the values from Bray et al. (2000). Bray
  tabulates 5.6e3-5.6e5 K; the fits passed through the first four table
  values and were an unsupported extrapolation below 5.6e3 K (1.10-1.16
  times the new strengths at 1.5e3 K). Now: Bray\textquotesingle s table inside its range
  (PCHIP in log-log; one misprint of Table 2 corrected), the CHIANTI he\_1
  temperature dependence joined continuously outside it; the feed through
  the higher triplets, the 2\^{}3S cooling, the 2\^{}3S -\textgreater{} n\^{}1L excitations and
  the He I cooling sums were rebuilt on Bray. On the LHS 1140 b He/H 1.8
  state the He 10830 EW rises by 2.3 per cent on evaluation.
\item
  H 2s-2p mixing by electrons: the coded rate (Janev et al. 2003, from
  Chibisov 1969) was ten times below Seaton\textquotesingle s (1955) first-order dipole
  constant; now Seaton (1955) eq. (55), approximation V (5.78e-5 cm\^{}3 s\^{}-1 at
  1e4 K against his Table 1 5.7e-5; Osterbrock \& Ferland Table 4.10 agree).
\item
  A(2s -\textgreater{} 1s) = 8.22461 s\^{}-1 (H) and 526.823 s\^{}-1 (He II), Drake (1986)
  eq. (27) (section 19 wrote 8.2249, the Nussbaumer \& Schmutz 1984 value it
  replaced).
\item
  The case-B share of recombinations into H 2s is now Pengelly\textquotesingle s (1964,
  Table I) hydrogenic fraction, the one the He II branch already used; the
  quadratic in T that section 17 (item B6) called "Draine\textquotesingle s 2s fraction" is
  Christie et al.\textquotesingle s (2013) fit, and Draine (2011) gives only Tables 14.2 and
  14.3.
\item
  The atomic-mode (no 2\^{}3S) He I case-B photons are now sent at their own
  exits (584 A, 19.8 eV, two-photon) split by the network\textquotesingle s rates, instead
  of one averaged energy (3.10 instead of 4.12 eV of heat a capture at high
  density; z = 0.966 and 0.660 at low and high density against Draine\textquotesingle s
  0.96 and 0.67).
\item
  Voronov (1997): all 20 rows equal the published Table I; its stated range
  (1 eV-20 keV) is written at the code with what the fit does below it.
  Anderson et al. (2000) confirm the CHIANTI H I strengths to 2.5 per cent at
  0.5-5 eV; Yan \& Babb (2023) published values equal the arXiv ones; Dere et
  al. (1997) and Del Zanna et al. (2015) say nothing about the descaling
  spline, which rests on the CHIANTI IDL code.
\end{itemize}

\textbf{Build.} Tree rebuilt from \texttt{make\ distclean}: \texttt{EXHALE.x} md5
\texttt{50aa18d79bcd25f219b5b59baa31b8b8}, no warning, copied to
\texttt{LHS1140b/models/EXHALE\_50aa18d7.x} with its manifest. The manifests of
2026-09-24/25 listed files as changed that were not (the header lines of the
previous manifest matched the file-name search); the three lists were rebuilt
from the md5 lines of consecutive manifests. The refined-grid runs of the
round-2 physics that were still running (He/H 1.6 on 2000 cells, 2.13 on 1000
and 2000) were stopped as superseded; the molecular family, the regression
fixture, the regression matrix and the suites are being re-run on this binary.

\textbf{The texts the corrections made stale} (\texttt{md/stale\_texts\_20260925.md}):
the physics document (\texttt{docs/physics\_overview/*.tex}, the PDF rebuilt, 70
pages), \texttt{docs/transmission\_spectrum.\{tex,md,pdf\}} and
\texttt{md/atomic\_data\_EXHALE\_vs\_MoCHII.md} were brought to the code. The transit
post-processor built the n = 2 populations of the Balmer lines with Christie
et al.\textquotesingle s rates while claiming to mirror the solver; it now uses the solver\textquotesingle s
rates, one transcribed function each (agreement with the Fortran 3.1e-13,
MEASURED), and keeps only its own Ly-alpha and Balmer fields (which the
profiles do not carry). H-alpha peak depth of HD 209458 b -1.73 per cent,
of \texttt{wasp\_full} +0.14 per cent (MEASURED), almost all from the proton
l-mixing, which lowers n(2s)/n(2) at 1.3-3 R\_p from 0.53-0.69 to 0.26-0.34.
The n = 2 threshold cross section 1.4e-17 cm\^{}2 is Seaton (1959) at the n = 2
edge (1.378e-17 DERIVED), not Osterbrock \& Ferland.

\subsection{23. The Markdown documents move to md/ (2026-09-25)}\label{23-the-markdown-documents-move-to-md-2026-09-25}

On the user\textquotesingle s instruction the 350 Markdown files directly under \texttt{docs/} moved
to the new top-level \texttt{md/}; \texttt{docs/} keeps the LaTeX and PDF documents, the
figures and the dated audit directories, whose 13 Markdown files stayed with
the files they describe. Done with \texttt{mv}, not \texttt{git\ mv}: git sees the 35 tracked
files as deleted under \texttt{docs/} and new under \texttt{md/}. Paths that followed:

\begin{itemize}
\tightlist
\item
  \texttt{.gitignore}: the rule that keeps the development record off the public
  repository (\texttt{docs/*.md} with 35 exceptions) is now \texttt{md/*.md} with the same
  35 exceptions under \texttt{md/}.
\item
  73 relative links of 24 moved files into \texttt{docs/} subdirectories now read
  \texttt{../docs/...}.
\item
  3079 references \texttt{docs/\textless{}name\textgreater{}.md} in 416 maintained files (sources and
  their comments, tests, scripts, \texttt{README.md}, \texttt{README\_HOWTO.md}, the \texttt{md/}
  files themselves, the \texttt{docs/} LaTeX, the regression README and scripts) now
  read \texttt{md/\textless{}name\textgreater{}.md}; four of them are strings the binary prints (the setup
  report, one lower-atmosphere message, one test message), so the tree binary
  changed (\texttt{EXHALE.x} md5 \texttt{b30e56f63335431dcb22b9fe3f80f948}, no warning); the
  state files the regression compares carry none of them.
\item
  \texttt{src/utils/update\_log\_to\_tex.py} reads the logs from \texttt{md/}.
\item
  Left as they were: run records and products (inputs, outputs, state files,
  \texttt{REPRODUCE.md}, logs, the \texttt{.P1} run areas, archived regression snapshots),
  whose text records the paths of their day and whose bytes other records
  hash; and the one string \texttt{LHS1140b/models/make\_models.py} writes into the 94
  stored catalog inputs, so that it still reproduces them byte for byte.
\item
  The workspace \texttt{CLAUDE.md} and the session memory point to \texttt{md/}.
\item
  The 17 \texttt{docs/} documents whose text names these paths were rebuilt on the
  user\textquotesingle s instruction (the two log twins regenerated first): every build
  without error or undefined reference, the new paths in the PDFs.
\end{itemize}

\subsection{24. The last items of the MoCHII comparison, and the molecular crossing under the corrected physics (2026-09-25/26)}\label{24-the-last-items-of-the-mochii-comparison-and-the-molecular-crossing-under-the-corrected-physics-2026-09-2526}

\texttt{md/EXHALE\_MOCHII\_ATOMIC\_COMPARISON\_2026-09-25.md} (a side-by-side of the
atomic processes of EXHALE and MoCHII, written from the MoCHII side) found the
recombination construction, the recombination cooling, the He 2\^{}3S network
and the H 2s-2p mixing now the same in both codes to four digits at 1e4 K,
and listed five open items for EXHALE; each was verified and applied (memos
\texttt{md/charge\_exchange\_detailed\_balance\_20260924.md} section 12 and
\texttt{md/photoionization\_584\_auger\_20260925.md}).

\begin{itemize}
\tightlist
\item
  \textbf{N2+ + H -\textgreater{} N+ + H+} (Barragan et al. 2006 Table 2, 9.70e-10 cm\^{}3 s\^{}-1 at
  1e4 K, the product N+ 2s2p\^{}3 3D0, so no reverse; heat 4.57 eV an event;
  \texttt{metals.inp:\ cx\_N2p\_H\ \textless{}scale\textgreater{}}, default 1). On \texttt{wasp\_full} N III falls by up
  to 5.4 decades.
\item
  \textbf{He2+ + H -\textgreater{} He+ + H+ + photon} (Kingdon \& Ferland 1996 Table 1, 1.0e-14,
  radiative, no reverse and no heat), in every system that carries He III;
  He III falls by 0.69 dex at the \texttt{wasp\_full} base. The original of the rate
  (West, Lane \& Cohen 1982, Phys. Rev. A 26, 3164) could not be obtained.
\item
  \textbf{Density-dependent detailed balance.} Where the source names the product
  level, the reverse is the thermal reverse times the population-weighted
  factor of the metastable levels at the local (T, n\_e) (five-level
  statistical equilibrium of N I, S II, C I from CHIANTI 11.0.2), and the
  heat of the forward channel removes the forbidden-line energy a metastable
  product radiates at low density; elsewhere the thermal reverse is exact at
  every density. \texttt{wasp\_full} (n\_e 8e8-3.5e9) moves by less than 0.2 per cent.
  A 0/0 below 60 K in the metastable Boltzmann weights (one LHS 1140 b base
  cell sits at 40.8 K) was guarded.
\item
  \textbf{He I 584 A.} The line is resonantly scattered by He0 inside the cell,
  and each scattering ends in absorption by H I (or H2, or a metal ion), in
  conversion to 2\^{}1S through the 2.06 um branch (1/912) and the existing
  two-photon exit, or in escape (Wood, Mathis \& Ercolano 2004 sect. 4.1,
  eqs. 15-17, generalized to every continuum absorber of the cell and to
  escape); their eq. 17 carries He/H = 0.1 in its coefficient.
\item
  \textbf{Inner-shell photoionization.} Each subshell absorption heats with the
  photoelectron plus the Auger energy, the fluorescence and the further
  ionization energy are not heat, and the neutrals of C, N, O, Mg, Si, Ca and
  Fe are carried straight to X++ where the vacancy autoionizes (Kaastra \&
  Mewe 1993 Tables 2 and 3; Ca II 3s is treated as radiative, since the
  measured energies forbid its autoionization). At 98 eV the electrons of a
  Mg II absorption receive 0.333 of the photon energy instead of 0.847.
\end{itemize}

\textbf{Build.} \texttt{EXHALE.x} md5 \texttt{f4da21f29579d87370efec2a4220f4c4}, no warning,
copied to \texttt{LHS1140b/models/EXHALE\_f4da21f2.x} with its manifest. On LHS 1140
b the two He items move the He 10830 EW by +0.05 per cent (He/H 1.8), but a
state of the previous physics does not certify under them, so the molecular
states are re-solved on this binary.

\textbf{The molecular K\_zz = 1e9 family under the physics of section 22}
(\texttt{EXHALE\_50aa18d7.x}; \texttt{.P1/newphys/}, \texttt{.P1/refine\_grid\_r3/}). He/H 1.8 and
2.13 re-solved from the round-2 states; 1.6 and 1.52 as rungs from 1.8 (1.45
failed its rung: the H2 carrier row held at 1e-1 through four passes).
MEASURED, every evaluation certified except the one marked:

{\def\LTcaptype{none} % do not increment counter
{\small
\begin{longtable}[]{@{}>{\raggedright\arraybackslash}p{\dimexpr(\linewidth-10\tabcolsep)/5\relax}>{\raggedright\arraybackslash}p{\dimexpr(\linewidth-10\tabcolsep)/5\relax}>{\raggedright\arraybackslash}p{\dimexpr(\linewidth-10\tabcolsep)/5\relax}>{\raggedright\arraybackslash}p{\dimexpr(\linewidth-10\tabcolsep)/5\relax}>{\raggedright\arraybackslash}p{\dimexpr(\linewidth-10\tabcolsep)/5\relax}@{}}
\toprule\noalign{}
He/H & EW, 500 cells & 1000 cells & 2000 cells & continuum (DERIVED) \\
\midrule\noalign{}
\endhead
\bottomrule\noalign{}
\endlastfoot
1.52 & 1.0038 & & & \\
1.6 & 1.0475 & 1.0187 & 1.0039 & 0.988 \\
1.8 & 1.1468 & 1.1162 & 1.1003 (a) & 1.083 \\
2.13 & 1.2853 & 1.2516 & 1.2343 & 1.216 \\
crossing of 1.108 & 1.721 & 1.783 & 1.819 & 1.860 \\
\end{longtable}
}
}

(a) the re-evaluation of the certified state refuses the base mass row by a
factor 1.9 (1.4e-8 against 7.4e-9); its EW is quoted for reference.

The EW converges at first order with the grid (successive-difference ratios
1.93-1.95, p = 0.95-0.97), as under the physics before the audit, and the
crossing moves from He/H 1.46 (500 cells) and 1.55 (continuum) before the
audit to 1.72 and 1.86. The 60 R\_p domain factor of section 14 was measured
before the audit and is not applied here.

\textbf{The documents brought to the code} (\texttt{md/user\_manual\_update\_20260926.md}):
\texttt{docs/EXHALE\_user\_manual.tex} (rebuilt, 103 pages, no error or undefined
reference, every changed page looked at), \texttt{README.md} and \texttt{README\_HOWTO.md}
now describe the input keys (\texttt{H\_rec\_escape}, \texttt{He\_rec\_coupling}, the
\texttt{metals.inp} keys \texttt{cx\_N2p\_H}, \texttt{cx\_O2p\_H}, \texttt{cx\_full}), the 24 columns of
\texttt{Heating\_breakdown.txt} (charge-exchange heat in column 24), the
\texttt{stage\_channels.txt} layout and the atomic data of sections 17-24. Found in
the code and fixed: \texttt{H\_rec\_escape} was parsed but missing from the list of
recognized keys (a spurious "unrecognized input line" warning, and a second
statement not refused); the missing-\texttt{metals.inp} message promised built-in
abundances that are zero; comments said \texttt{cx\_full} toggles the He-H pair and
that O2+/N2+ + H are dropped; the \texttt{stage\_channels.txt} header omitted the
metal charge-exchange channels 27-32. \texttt{md/input\_schema.md} K14e (still
"Koskinen 2013") and its \texttt{metals.inp} table (no group-E keys) were brought up
to date. Not changed: \texttt{H\_rec\_escape} and \texttt{He\_rec\_coupling}, like the other
physics switches of their group, are neither echoed in the setup report nor
restart option tokens.

\textbf{Checks of the binary of this section} (\texttt{EXHALE\_f4da21f2.x}). The suites:
\texttt{physics\_probe} did not compile one new probe (\texttt{inner\_shell\_photoabsorption\_ledger}
called \texttt{cx\_set\_cell} without the electron density the density-dependent
reverses added; the two items were written at the same time); fixed, and the
suite passes. \texttt{grid\_and\_gates} fails its fixture rows plus
\texttt{outer\_iteration\_ending\_is\_the\_stagnation\_one} and
\texttt{restart\_lower\_ghost\_is\_the\_derived\_boundary} (7.9e-14 against 1e-14), all on
the pinned fixture of the previous physics; \texttt{residual\_determinism} keeps its
known row; the two \texttt{steady\_*} suites need a solve log. The regression matrix
ran all 17 cases; every one differs from the goldens by the physics
(\texttt{wasp\_full} log10 Mdot and the base inflow: 214 inflow cells against 257).
\textbf{The user held the golden refresh (2026-09-26).} The stationary re-solve
of the fixture from the previous one did not converge (info = 2: the energy
row 3.1e-2 at cell 406, the closure of \texttt{System\_HeH\_TR\_metals} 4.6e-3 at cell
340), while the matrix case \texttt{wasp\_full\_newton} (marching, then Newton)
certified (info = 0, \textbar\textbar R\textbar\textbar{} 4.4e-9); that output is the candidate fixture.
The tree was rebuilt after the comment and message corrections of this
section (\texttt{EXHALE.x} md5 \texttt{0a7532d5087f14df41142fb499b68240}; it differs from
\texttt{EXHALE\_f4da21f2.x} in printed strings and comments only).

\textbf{The molecular states on \texttt{EXHALE\_f4da21f2.x}} (\texttt{.P1/newphys/HeH\textless{}x\textgreater{}\_r4},
re-solved from the certified states of the previous physics; certified in
9, 9, 8 and 7 passes; every evaluation certified). He 10830 EW, 500 cells:
He/H 1.52 1.0047, 1.6 1.0483, 1.8 1.1477, 2.13 1.2862 (+0.06-0.09 per cent
against the table above); the 1.108 crossing 1.7194 (1.7213 before). The
refined grids were not re-solved on this binary; the crossing\textquotesingle s grid factor
above stands within the size of this change.

\subsection{25. The helium-hydrogen charge exchange from its originals (2026-09-26)}\label{25-the-helium-hydrogen-charge-exchange-from-its-originals-2026-09-26}

The He2+ + H original the user supplied (West, Lane \& Cohen 1982) and an
analysis made on the MoCHII side of the He+ + H exchange
(\texttt{\textasciitilde{}/MoCHII/MoCHII\_v1.00/md/heii\_h\_charge\_exchange/}) were checked against
the published papers and applied (\texttt{md/charge\_exchange\_detailed\_balance\_20260924.md}
section 12.9). Every rate below is DERIVED from the cited cross sections
unless marked READ.

\begin{itemize}
\tightlist
\item
  \textbf{He2+ + H -\textgreater{} He+(1s) + H+ + photon.} The 1.0e-14 cm\^{}3 s\^{}-1 of Kingdon \&
  Ferland (1996) is not a result of any calculation: it is the value they set
  for every reaction without a significant channel (READ, their sect. 2).
  Now the Maxwellian average of the West et al. (1982) radiative cross
  sections (Figs. 4-6, with the resolved resonances): 2.0e-13 at 1e3 K,
  1.65e-13 at 1e4 K, 1.57e-13 at 3e4 K, 16-20 times the old value; a
  calculation above 1e3 K, an extrapolation below.
\item
  \textbf{He+ + H -\textgreater{} He + H+, non-radiative.} The Kingdon \& Ferland fit reproduces
  the singlet-conditional cross sections of Zygelman et al. (1989) without
  the 1/4 of the spin statistics, and those cross sections are themselves
  hundreds of times those of Kimura et al. (1993) and Loreau et al. (2014).
  Now the Loreau et al. (2014) H+ + He cross section by microscopic
  reversibility (weights 4:1, threshold 13.65 eV/u, zero below 15 eV/u):
  4.4e-17 at 1e4 K instead of 9.0e-15, 1.5 per cent of the radiative channel;
  the band at the code (up to 2.9x at 1e4 K from the threshold treatment;
  0.4-2.5x above 3e4 K). The reverse He + H+ (B1) follows by detailed balance
  (5.1e-22 at 1e4 K instead of 1.0e-19); the heat of an event is unchanged.
\item
  \textbf{He+ + H, radiative.} Stancil, Lepp \& Dalgarno (1998) Table 1 row 19
  prints 1.20e-15 (READ); the code carried the 1.25e-15 of Glover \& Jappsen
  (2007). Now 1.20e-15.
\end{itemize}

Effect (MEASURED): \texttt{wasp\_full} log10 Mdot -0.002 (inside the marching spread),
He III down to 0.94 dex at the base, He II and He 2\^{}3S +48 and +52 per cent at
1.1 R\_p, the charge-exchange share of the heating 5.6 -\textgreater{} 0.9 per cent at 1.05
R\_p. \textbf{LHS 1140 b He/H 1.8 (evaluation of the certified state): the old
He+ + H exchange supplied up to 10.7 per cent of the heating; without it the
peak temperature falls from about 5600 to 5150 K, He II and He 2\^{}3S rise by
14 and 15 per cent at 1.5 R\_p, and the He 10830 EW rises from 1.1477 to 1.1818
\%A (+3.0 per cent).} The certificate does not hold under the corrected rates
(energy row 8.4e-2), so the molecular states are re-solved on the binary of
this section (\texttt{EXHALE.x} md5 \texttt{59c6be36820b8ba5c958c294ad9e6f02}, copied to
\texttt{LHS1140b/models/EXHALE\_59c6be36.x} with its manifest).

\textbf{The molecular states re-solved on \texttt{EXHALE\_59c6be36.x}} (\texttt{.P1/newphys/HeH\textless{}x\textgreater{}\_r5},
from the certified states of section 24): He/H 1.6, 1.8 and 2.13 certified in
25, 23 and 21 passes, every evaluation certified. He 10830 EW, 500 cells
(MEASURED): 1.6 1.0577, 1.8 1.1550, 2.13 1.2902 (+0.9, +0.6 and +0.3 per cent
against section 24; the +3.0 per cent of the evaluation at fixed wind is
mostly given back by the solved wind); the 1.108 crossing moves from 1.7194
to 1.7026 (DERIVED). He/H 1.52 stopped in its second pass: a JFNK trial
iterate with \textbar\textbar R\textbar\textbar{} near 2 left a lower ghost whose base H2 handoff partition
did not close in 30 passes (last move 2.0e-6 against 1e-10), and the guard of
\texttt{ionization\_equilibrium.f90} (a boundary nothing stands behind is refused)
ends the run with \texttt{error\ stop} even on a trial state. Whether a trial should
instead be refused and the step shortened is a solver-control question put
to the user.

\subsection{26. A trial state no longer stops the run at an unclosed lower ghost (2026-09-26)}\label{26-a-trial-state-no-longer-stops-the-run-at-an-unclosed-lower-ghost-2026-09-26}

User decision: as recommended (\texttt{md/ghost\_closure\_trial\_refusal\_20260926.md}).
The guard of \texttt{ionization\_equilibrium.f90} that stops the run when the lower
ghost\textquotesingle s base H2 handoff partition does not close against its own ionization
balance (tolerance 1e-10, 30 passes) now distinguishes a state a caller can
refuse from one it cannot: the stationary JFNK\textquotesingle s trial evaluations
(\texttt{eval\_residual} of \texttt{steady\_newton.f90} when it screens admissibility: line
search, trust region, Levenberg-Marquardt and pseudo-time trials, the
finite-difference probes) and the relaxation\textquotesingle s trial chemistry
(\texttt{equilibrate\_chemistry\_at\_fixed\_conserved\_state}) declare it before the
sweep, and there an unclosed ghost refuses the trial (logged as \texttt{CANDIDATE\ REFUSED}, counted in the ledger, refusal code 7, closure reason 5), the
solver then shortening the step as for its other refusals; an accepted state,
the state written, the marching sweep and every other path keep the stop. The
same rule now covers the ghost-composition fixed point one step later. MEASURED:
the He/H 1.52 re-solve of section 25 passes the trial that stopped it (the
lambda = 1 trial of JFNK solve 3 iteration 7; the halved trials close) and
certifies at pass 24; an evaluate of He/H 1.6 and a copy of \texttt{mol\_carrier}
are byte-identical to the build before (\texttt{EXHALE\_59c6be36.x}) in every data
line; the suites pass, the two failing ones (\texttt{residual\_determinism},
\texttt{grid\_and\_gates} fixture rows) failing identically on the build before.
\texttt{EXHALE.x} md5 \texttt{7db5afe9b1d80ce757fcc07d1b499330}, copied to
\texttt{LHS1140b/models/EXHALE\_7db5afe9.x} with its manifest.

The He/H 1.52 state on this physics: EW 1.0149 \%A (certified evaluation); with
1.6, 1.8 and 2.13 of section 25 the 500-cell ladder is 1.0149, 1.0577,
1.1550, 1.2902 and the crossing of 1.108 stays at 1.7026.

\textbf{The grid of the crossing on this physics} (\texttt{.P1/refine\_grid\_r5/}, binary
\texttt{EXHALE\_59c6be36.x}, the 1000-cell He/H 1.8 and 2.13 continued on
\texttt{EXHALE\_7db5afe9.x}, identical physics; every state and evaluation
certified). He 10830 EW, MEASURED; continuum DERIVED (Richardson at the
measured order):

{\def\LTcaptype{none} % do not increment counter
{\small
\begin{longtable}[]{@{}>{\raggedright\arraybackslash}p{\dimexpr(\linewidth-10\tabcolsep)/5\relax}>{\raggedright\arraybackslash}p{\dimexpr(\linewidth-10\tabcolsep)/5\relax}>{\raggedright\arraybackslash}p{\dimexpr(\linewidth-10\tabcolsep)/5\relax}>{\raggedright\arraybackslash}p{\dimexpr(\linewidth-10\tabcolsep)/5\relax}>{\raggedright\arraybackslash}p{\dimexpr(\linewidth-10\tabcolsep)/5\relax}@{}}
\toprule\noalign{}
He/H & 500 cells & 1000 cells & 2000 cells & continuum \\
\midrule\noalign{}
\endhead
\bottomrule\noalign{}
\endlastfoot
1.6 & 1.0577 & 1.0297 & 1.0150 & 0.999 \\
1.8 & 1.1550 & 1.1249 & 1.1094 & 1.093 \\
2.13 & 1.2902 & 1.2571 & 1.2403 & 1.223 \\
crossing of 1.108 & 1.703 & 1.764 & 1.797 & 1.838 \\
\end{longtable}
}
}

First order again (successive-difference ratios 1.91-1.96, p = 0.93-0.97).
Against the physics before the audit (1.464 on 500 cells, 1.552 in the
continuum) the crossing has moved to 1.703 and 1.838; the continuum value is
interpolated between the certified continuum points of 1.8 and 2.13.

\subsection{27. The molecular family under the corrected physics: catalog publication and the width of the line (2026-09-26)}\label{27-the-molecular-family-under-the-corrected-physics-catalog-publication-and-the-width-of-the-line-2026-09-26}

\textbf{Catalog.} The certified He/H 1.52, 1.6, 1.8 and 2.13 states of sections 25
and 26 were published into \texttt{LHS1140b/models/molecular\_scalar\_gj1132\_kzz1e9/}
(\texttt{run\_case.sh}, \texttt{SEED=} the certified run, \texttt{EXHALE\_7db5afe9.x}): every re-
measurement "original claim REPRODUCED", every evaluation certified; EW 1.0149,
1.0577, 1.1550, 1.2902 \%A, log10 Mdot 7.89, 7.90, 7.91, 7.93. The catalog\textquotesingle s
He/H 0.55, 1.45 and 9.7 molecular states are of the physics before the audit
and were not re-solved.

\textbf{The added Gaussian} (\texttt{.P1/connectivity/broaden\_molecular\_final.py}, the
operations of section 13; output \texttt{broaden\_molecular\_final.txt}, figure
\texttt{docs/figures/lhs1140b\_kzz1e9\_molecular\_broadened\_final.pdf}). The memo check
(atomic well-mixed He/H 0.40) reproduces f = 22.22 km/s. MEASURED: the matched
kernel is 22.10-22.19 km/s FWHM on every state and grid (sigma 9.39-9.42
km/s), and at it the red depth, blue depth and ratio of the He/H 1.8 state
lie inside the observed 1 sigma (red 1.289, blue 0.170, ratio 7.60). The
broadened EW (MEASURED on each grid, continuum DERIVED):

{\def\LTcaptype{none} % do not increment counter
{\small
\begin{longtable}[]{@{}>{\raggedright\arraybackslash}p{\dimexpr(\linewidth-10\tabcolsep)/5\relax}>{\raggedright\arraybackslash}p{\dimexpr(\linewidth-10\tabcolsep)/5\relax}>{\raggedright\arraybackslash}p{\dimexpr(\linewidth-10\tabcolsep)/5\relax}>{\raggedright\arraybackslash}p{\dimexpr(\linewidth-10\tabcolsep)/5\relax}>{\raggedright\arraybackslash}p{\dimexpr(\linewidth-10\tabcolsep)/5\relax}@{}}
\toprule\noalign{}
He/H & 500 cells & 1000 cells & 2000 cells & continuum \\
\midrule\noalign{}
\endhead
\bottomrule\noalign{}
\endlastfoot
1.6 & 1.0227 & 0.9955 & 0.9813 & 0.966 \\
1.8 & 1.1167 & 1.0876 & 1.0726 & 1.057 \\
2.13 & 1.2474 & 1.2155 & 1.1992 & 1.182 \\
crossing of 1.108 & 1.781 & 1.851 & 1.890 & 1.933 \\
\end{longtable}
}
}

As turbulence the kernel is supersonic: on the He/H 1.6 and 1.8 states the
line-forming region (n(2\^{}3S) r weighting) has T = 2650 and 2756 K, mu = 1.94
and 2.03, c\_s = 4.33 km/s, so sigma = 9.4 km/s is Mach 2.2 as a line-of-sight
dispersion, 3.1 as sqrt(2) sigma and 3.8 as an isotropic speed, as before the
audit. From the EW alone the crossing is 1.703 (500 cells) and 1.838
(continuum); from the broadened line 1.781 and 1.933.

\subsection{28. The catalog summary brought to the tree\textquotesingle s state, and why the molecular winds do not certify (2026-09-26)}\label{28-the-catalog-summary-brought-to-the-trees-state-and-why-the-molecular-winds-do-not-certify-2026-09-26}

\texttt{docs/lhs1140b\_catalog\_state.tex} (13 pages, rebuilt, pages 5 to 8 checked
by eye):

\begin{itemize}
\tightlist
\item
  The certified table carries the molecular K\_zz = 1e9 rows the catalog now
  publishes, READ from each case\textquotesingle s \texttt{run.log} and \texttt{models/status.py}: 1.52,
  1.6, 1.8 and 2.13 under \texttt{EXHALE\_7db5afe9.x} (note a); 1.45 under
  \texttt{EXHALE\_b059a4be.x}, the second round of corrections, EW 0.9387 (note b;
  under \texttt{EXHALE\_50aa18d7.x} neither its re-solve nor its rung from 1.52
  certified within the time ceiling, and it was not attempted under the
  binary of record); 9.7 (note c, below). The rows of the table before this
  entry were the pre-audit values (1.6: EW 1.2155, now 1.0577).
\item
  A new subsection of section 3, "Why the molecular winds do not certify",
  assembled from \texttt{md/lhs1140b\_stationary\_L34c\_20260918.md},
  \texttt{md/lhs1140b\_alternation\_audit\_20260921.md},
  \texttt{md/lhs1140b\_coupled\_trial\_20260921.md} and \texttt{md/TO\_BE\_DONE.md} L22: the
  composition half of the partitioned route refuses (the H2 carrier row in
  all six cases), by two measured mechanisms (the composition update takes
  the wind out of solution in the well-mixed cases; a drift under a
  saturated movement bound in \texttt{kzz1e9/HeH0.083}); the coupled route fails
  in its Krylov solve; the open item is a preconditioner that holds the
  species rows. Every one of those measurements predates the atomic-data
  corrections and none was repeated under \texttt{EXHALE\_7db5afe9.x}.
\item
  The two molecular cases that published no state
  (\texttt{molecular\_scalar\_gj1132\_wellmixed/HeH2.13},
  \texttt{molecular\_photochem\_gj1132\_kzzprofile/HeH2.09}) are now named.
\end{itemize}

Found and fixed: \texttt{LHS1140b/models/status.py} printed the \texttt{not\_solved.md}
reason beside \texttt{molecular\_scalar\_gj1132\_kzz1e9/HeH9.7}, which is certified.
The certified generation \texttt{g0008\_20260920T052941Z\_c3c466a4} is the
2026-09-20 evaluation of the state written on 2026-09-19; the
\texttt{not\_solved.md} is that of a re-solve made later the same day
(\texttt{g0006\_20260920T000506Z\_5b2061aa}, \texttt{composition\_refusal}, the H2 carrier
row 5.13e-4 at cell 262 after 40 passes), which did not replace it. A
certified row with a \texttt{not\_solved.md} beside it now reads "a later re-solve,
not the state shown: ..." (MEASURED on the table output). \texttt{MODELS.md} was
not rewritten.

\subsection{29. Errors fixed and models solved first: the plan of 2026-09-26}\label{29-errors-fixed-and-models-solved-first-the-plan-of-2026-09-26}

User instruction (2026-09-26): re-checking solutions that move little
against the observational error is not the priority; fixing errors and
solving the models that do not solve is. Plan \texttt{md/PLAN\_20260926.md}; the
user approved its recommendations (goldens and fixture refreshed; the
composition-side solver work reopened if continuation does not close the
cases; \texttt{diffusion\_check} taken out of the model count). New physics only
where it can change results seriously; the rest is recorded, not applied.

\subsubsection{29.1 An OpenMP data race in the metal sweep (fixed)}\label{291-an-openmp-data-race-in-the-metal-sweep-fixed}

\texttt{meg\_g02} (the inner-shell Auger share of neutral photoionization to X++,
added with binary f4da21f2 on 2026-09-25) was missing from the private list
of the He + metals cell sweep (\texttt{ionization\_equilibrium.f90}, then \textasciitilde2137, a
\texttt{default(shared)} loop), so threads overwrote each other\textquotesingle s value between the
write and the two reads. MEASURED: route-B stationary restart of
\texttt{wasp\_full\_newton} failed at 8 threads (composition elimination never
contracts, \texttt{info\ =\ 2}) and certified at 1 thread; two 8-thread runs of one
input differed before the solve started (\texttt{md/metal\_sweep\_openmp\_race\_20260926.md}).
Every metals-on run at more than one thread since f4da21f2 was affected;
the regression matrix (1 thread) and every metals-on catalog certificate
(last published 2026-09-20) were not. Fix: \texttt{meg\_g02} private, and all four
cell loops of \texttt{ioniz\_eq} converted to \texttt{default(none)} with explicit shared
lists; an audit of every parallel region of \texttt{radiation/} and
\texttt{lower\_atmosphere/} found one more, formal, race (two table-ready flags read
without synchronization; now \texttt{omp\ atomic\ read/write\ seq\_cst}). New suite
\texttt{src/tests/thread\_count\_identity/} (200 steps of \texttt{wasp\_full} at 1 and 8
threads, data lines bitwise): 7db5afe9 FAIL (3.2e-3 Hydro, 5.4e-2 Ion),
fixed binary PASS. The r0 case (route B, 8 threads) now certifies,
\texttt{info\ =\ 0}, \textbar\textbar R\textbar\textbar{} 6.50e-9. No 1-thread run moved.

\subsubsection{29.2 Fix batch (errors confirmed in the source)}\label{292-fix-batch-errors-confirmed-in-the-source}

F1 free-free Gaunt factor: van Hoof et al. (2014) eq. (21) uses the
infinite-mass Rydberg (READ, their sect. 2), so Ry/k = 157887.51 K, not
157807 (the Hui \& Gnedin 1997 constant, kept in that fit); F2 the carrier
diffusivity read by outputs and diagnostics was zero on a stationary path
(MEASURED D\_H2 = 0 at every cell; now 7.9e6 - 8.2e9 cm\^{}2/s); F3 the
determinism replay now passes \texttt{u} to \texttt{cell\_row\_scales}; F4 pressure formed
before the composition refresh at 3 of 9 sites (coupled-block entry, after
\texttt{solve\_steady\_*}, trial end; the other 6 verified not stale); F5 the CO
column, dissociation rate and shielding factor joined the step checkpoint;
F6 an inert \texttt{Coupled\ carrier\ solve} is reported; F7 five wrong comments;
F8 the \texttt{\_adv} headers carry every provenance field; F10 seven callerless
\texttt{Cool\_coeff} wrappers removed and the Voronov H/He fit defined once
(rates bitwise unchanged). Effect on the regression inputs (single thread):
F1 alone, at most 1.4e-4 relative (\texttt{wasp\_full} velocity), log10 Mdot
unchanged to 4 digits.

\subsubsection{29.3 The interdiffusion enthalpy flux (new physics, serious)}\label{293-the-interdiffusion-enthalpy-flux-new-physics-serious}

The energy equation of a diffusing mixture carries q\_d = sum\_i h\_i J\_i
(Cook 2009, Phys. Fluids 21, 055109, eqs. 11-13; Schunk 1977, Rev. Geophys.
Space Phys. 15, 429, eqs. 15b, 16, 20c; READ). EXHALE moved He relative to H
and carried no such term in either energy equation. MEASURED size on LHS
1140 b He\_diffusion states: 13-26 \% of the local energy budget at
1.1-5 R\_p (\texttt{md/interdiffusion\_enthalpy\_flux\_20260926.md}). Implemented in
both routes (binary model: q\_d = J\_He (h\_He - h\_H), sensible enthalpies
with the electrons each ion carries; formation energy verified to be
carried already); key \texttt{Interdiffusion\ enthalpy\ flux:} default True with
\texttt{He\_diffusion} (the term belongs to the equation), False only to match
published models that omit it (Koskinen 2013, 2022; Yelle 2004); restart
token \texttt{interdiff\_enth} (older states restart under
\texttt{Restart\ option\ change:\ interdiff\_enth}). Suite
\texttt{src/tests/interdiffusion\_enthalpy\_flux/} 18/18. Effect on the molecular
He/H 1.8 state (24 + 24 passes, carrier row 1.6e-5 at the end, not yet
certified): T -166 K median over 1.3-2 R\_p (largest -236 K, -4.4 \%),
He/H -3.1 \%, the n(2\^{}3S) column -1.3 \%, log10 Mdot +0.002.

\subsubsection{29.4 He+ + H non-radiative charge exchange near threshold}\label{294-he--h-non-radiative-charge-exchange-near-threshold}

Section 25 took the Loreau et al. (2014) H+ + He cross section forward by
reciprocity with sigma\_r = 0 below 15 eV/u, where their calculation starts.
That closes an open channel: the forward (exothermic) rate then fell as
exp(-E/kT) at low T (5e-34 at 300 K). Adopted, as in MoCHII (the user\textquotesingle s
decision relayed by the MoCHII session), sigma\_r(x) = sigma\_r(15)
sqrt((x - x\_th)/(15 - x\_th)) between x\_th = 13.655 eV/u and 15 eV/u, i.e.
sigma\_f v constant at low energy (ion-induced-dipole capture with a
constant transfer probability; an assumption no calculation tests).
Rate: 9.64e-17 at 1e4 K (was 4.41e-17), 9.56e-17 at 1e3 K (was 1.5e-21);
3.3 \% of the radiative channel at 1e4 K. Generator
\texttt{cooling\_data/heii\_h\_charge\_exchange\_rate.py} (reproduces the old table to
0.4 \% with the old cut).

The radiative channel, the user asked, is one coefficient in both codes:
Stancil, Lepp \& Dalgarno (1998) Table 1 row (19), 1.20e-15 (T/300)\^{}0.25,
the Zygelman et al. (1989) Table I DIRECT column times the 1/4 of the
singlet approach (READ; 1.17e-15 at 300 K from the original). MoCHII moved
from the Glover \& Jappsen (2007) R26 1.25e-15, which cites the same data
and cannot be reproduced from it (2026-09-26); EXHALE unchanged. A
digitization of Courtney et al. (2021) Fig. 3 agrees with this fit within
2 \% over 1e3 - 1e4 K and lies 1.9x above it at 100 K; a fit to their
tabulated rate needs the data file (University of Georgia Molecular
Opacity Project, not reachable from here).

\subsubsection{29.5 Goldens, fixture, suites}\label{295-goldens-fixture-suites}

Goldens re-snapshotted from the tree binary 4e275f1d (the physics of
7db5afe9) after a full matrix (all 17 cases differed from the round-3
goldens, as expected); backup \texttt{golden\_pre\_round45\_20260926}. Fixture
\texttt{wasp\_full\_newton/IC} replaced by the 1-thread route-B certified state of
7db5afe9 (backup \texttt{IC\_pre\_round45\_20260926}). Every suite on the integrated
binary: all pass except grid\_and\_gates (3 rows: the fixture is not a root of
the fixed physics, to be re-made with the final binary), residual\_determinism
(known row), the two suites that need run logs.

Closing the series (2026-09-27): the matrix on the final binary 6844f188
(sections 29.1-29.4, 29.7, 29.9) differs from the goldens of 7db5afe9 in
every case beyond 1e-3, largest in the helium columns of the lower cells
(3.7 \% at row 1 of \texttt{wasp\_full}, the He+ + H rate at low T) and in the
molecular cases (the ghost start, the progress control); goldens refreshed
from it (backup \texttt{golden\_pre\_final\_20260927}). Fixture re-solved on 6844f188
at 8 threads (\texttt{info\ =\ 0}, \textbar\textbar R\textbar\textbar{} 4.11e-9; backup \texttt{IC\_pre\_final\_20260927});
\texttt{grid\_and\_gates} 345/345. Every suite passes except the known
\texttt{residual\_determinism} row and the two suites that need run logs.

\subsubsection{29.6 Models solved}\label{296-models-solved}

MEASURED, \texttt{.P1/s1\_remeasure/}, \texttt{.P1/s2\_continuation/}: every failure of the
catalog\textquotesingle s unsolved molecular and atomic cases stands under 7db5afe9 (six
passes each); continuation then solved:

\begin{itemize}
\tightlist
\item
  \texttt{molecular\_scalar\_gj1132\_kzz1e9/HeH1.45}: rung from the certified 1.52,
  certified in 13 passes.
\item
  \texttt{molecular\_scalar\_gj1132\_wellmixed/HeH2.13}, 1.8, 1.4: the well-mixed
  molecular cases differ from the K\_zz = 1e9 ones ONLY by \texttt{He\_diffusion}
  off (their \texttt{base.inp} carries \texttt{Kzz\_base\ 1.0e9}, which the carrier
  transport reads). From the certified kzz1e9/2.13 state, He/H set uniform
  (new \texttt{map\_state\_to\_grid.py\ -\/-uniform}) and diffusion off: certified in 10
  passes; then rungs 1.8 (10 passes) and 1.4 (14 passes) with
  \texttt{make\_rung.sh} (\texttt{UNIFORM=1}). The direct re-solve of the catalog\textquotesingle s
  well-mixed 2.13 state is refused by the progress control at pass 5, so
  the continuation is what solves it.
\item
  \texttt{molecular\_photochem\_gj1132\_kzzprofile/HeH9}: on the final binary of
  the series (2dd9a865, face at the level), mapped from its 40-pass state
  on 6844f188 (\texttt{.P1/s2\_continuation/photo9\_f}, row 7.4e-5 falling 0.77 a
  pass) and continued 40 passes: certified (H2 carrier row 9.2e-6 at cell
  217), \texttt{.P1/s2\_continuation/photo9\_v4} (2026-09-27). Its earlier stops were
  the lower-ghost closure defect of section 29.9.
\item
  Ionization-stage transport on the molecular 1.8 state stalls on the He++
  row under the progress control; the cause is the control (29.7).
\end{itemize}

\subsubsection{29.7 A progress-control defect of the composition relaxation}\label{297-a-progress-control-defect-of-the-composition-relaxation}

Diagnosed (\texttt{.P1/iontr\_diag/}): each pass of the carrier relaxation ends on
its particle-count cap, the state moves a fixed amount a pass, the relative
gated measure stays flat because residual and scale fall together, and the
outer control reads that as no progress, halves the cap and ends the loop.
With the cap held (diagnostic override) the same state certified in 14
passes.

Fixed (\texttt{EXHALE\_main.f90} progress control, binary d4b571e8): a pass whose
consumed carrier update ended on its movement bound, whose solve converged
and whose joint distance did not rise by more than 1.1x is not judged by the
relative measures and does not halve the bound; the bound doubles after such
a pass up to 0.5 and halves when a bound-limited update raises the distance
of a converged wind or precedes the first failed solve; omega follows the
same rule and no longer drops on hydro entry failures;
\texttt{EXHALE\_CARRIER\_TRUST\_HOLD} still freezes the bound. Also fixed: the scale
\texttt{x\_ref} of the carrier fixed-point test was the H2 maximum alone, so with
transported stages and no molecules it sat at its 1e-30 floor and the test
could never pass (logged displacement \textasciitilde1.7e28, MEASURED on
\texttt{iontrans\_metals}). MEASURED: the hot-Uranus stagnation fixture certifies at
pass 42 (it stagnated; the suite row that needs that ending now holds the
bound); \texttt{iontrans\_metals} accepts at pass 17 instead of 48, the two
certified states within 7.6e-5; the molecular He/H 1.8 state with transported
stages certifies in 14 passes from a starting bound of 0.5
(\texttt{EXHALE\_CARRIER\_TRUST=0.5}) but not from the default 1e-2, where the hydro
solve fails at pass 2 under the old control and the new one alike; with the
interdiffusion flux on, that state\textquotesingle s hydro solve fails on entry and a slow
mode on the He front (0.86-0.93 a pass) remains. REGRESSION on
\texttt{kzz1e9/HeH0.083}: a growing bound makes the H2 row at the outer cells
498-500 overshoot (up to 1.5e-1) and the bound saws; neither control
certifies that case (the old one stalls on the elemental row at cell 338).
\texttt{mol\_carrier} bitwise unchanged (the outer loop is not entered). Open: the
starting bound for transported stages, the H2 outer-cell overshoot, the
hydro entry failures with the interdiffusion flux.

\subsubsection{29.8 Incidents of this series}\label{298-incidents-of-this-series}

\begin{itemize}
\tightlist
\item
  A \texttt{sed} insertion with an empty line address put a three-line comment
  before every line of \texttt{src/tests/boundary\_state/molecular\_seed\_of\_an\_atomic\_state.sh};
  the 153 blocks were removed and the file checked by diff against the
  afternoon copy (identical except the intended line).
\item
  The tree binary was rebuilt while a regression matrix was running on it;
  the matrix was stopped and is re-run after the series.
\end{itemize}

\subsubsection{29.9 The lower-ghost H2 closure stalled with a zero move (fixed)}\label{299-the-lower-ghost-h2-closure-stalled-with-a-zero-move-fixed}

The molecular photochemical-column cases (He/H 9 and 2.09) stopped at
passes 1-3 with "the base handoff partition of a lower ghost did not close":
the closure\textquotesingle s own move exactly zero and its residual 9-41 times the 1e-10
tolerance. Diagnosed (\texttt{.P1/ghost\_stop/}): each closure pass after the first
started the molecular solve from the previous pass\textquotesingle s accepted state, whose
x(4) held the OLD imposed H2 partition; the pinned row x(4) - x\_h2\_fix then
entered hybrd1 carrying the whole secant move, and hybrd1 (xtol = sqrt(eps),
diag = 1) can return a move below 2 xtol \textbar\textbar x\textbar\textbar{} \textasciitilde{} 3e-8 untouched with
info = 1 (hybrd.f90: first delta = pnorm, halved on the rejected step, then
delta \textless= xtol xnorm). The write-back then imposed the value, the next pass\textquotesingle s
map saw no move, and 28 identical solves followed (27 of 27 traced events;
41 of 41 refusals and stops in the S1 logs, residuals 8.2e-10 - 2.8e-8, all
below that floor). Fix: \texttt{impose\_transported\_fractions} sets every unknown a
transport solve owns (H2 partition, oxygen carriers, ionization stages) to
its imposed value before both the molecular and the atomic hybrd1 solves;
the stop report prints the pin error. MEASURED with the fix (worker build):
0 of 127092 ghost closures unclosed, pin error exactly 0, no refusal and no
stop; the He/H 9 case then fails on its H2 carrier row (0.88), a separate
problem. Binary 6844f188. Every molecular base-handoff regression case is
reached (ghost composition moves by 1e-12 - 1e-9).

\subsubsection{\texorpdfstring{29.10 The \texttt{\_adv} post-process made second order}{29.10 The \_adv post-process made second order}}\label{2910-the-_adv-post-process-made-second-order}

User decision (2026-09-27): the catalog quotes the He 10830 EW of the \texttt{\_adv}
post-process. That post-process was a first-order upwind recursion in r
(backward Euler of dx/dr = R/v with the rates at r\_j and the residence time
at r\_\{j-1\}), the main source of the first-order grid convergence of the EW
(\texttt{md/ew\_grid\_order\_20260926.md}). Now the variable-step BDF2,
y\_j - a1 y\_\{j-1\} + a2 y\_\{j-2\} = g h\_j F(y\_j, r\_j) with w = h\_j/h\_\{j-1\},
a1 = (1+w)\^{}2/(1+2w), a2 = w\^{}2/(1+2w), g = (1+w)/(1+2w) (derived at
\texttt{post\_process\_adv.f90}, coefficients as eq. (1.4) of Li \& Liao,
arXiv:2201.00527; L-stable, so a population that equilibrates within a cell
does not ring as under the trapezoid rule), for every composition system
and the energy equation, rates and 1/v at the same point; backward Euler on
the first step, after a refused cell, for w \textgreater{} 1 + sqrt(2), and as a retake
of a step with a negative population or no temperature root (0 retakes on
the nine states). MEASURED on the nine certified grid states (He/H 1.6, 1.8,
2.13 on 500, 1000, 2000 cells), with the state held fixed: the post-process
share of the EW error falls from +3.0 \% / +1.0 \% (p = 1.00) to below
3e-4, and the \texttt{\_adv} 2\^{}3S column converges at p = 1.98 - 2.14 (1.00 - 1.01
before); a closed-form column test gives p = 1.989. The EW itself changes
by -3.9 \% (500 cells) to -0.97 \% (2000 cells) at He/H 1.6 (1.05692 -\textgreater{}
1.01571 on 500 cells) and still converges at first order: what remains is
the solved state\textquotesingle s first-order base closure (being fixed) and the transit
integration (p \textasciitilde{} 1.4, opposite sign).

Also fixed: the h div(rho v) term of the \texttt{\_adv} energy equation was built
from differences of the cell-centred rho v r\^{}2, which carries the odd-even
pattern of the cell-centred velocity at the base (1.8e12 to 3.4e12 between
neighbours while the face mass flux changed by 1e-10), giving 805 / 4008 /
579 K at the first corrected cell on 500 / 1000 / 2000 cells; it is now the
state\textquotesingle s own mass row, the face-flux divergence. Below about 1.13 R\_p the
molecule-free post-process lacks the molecular heating and H3+ cooling that
set the run\textquotesingle s energy balance there; by the user\textquotesingle s decision (2026-09-27) the
run T is kept there with status 3 "unsupported" wherever the omitted heating
and cooling exceed the carried ones (196 / 389 / 776 cells on 500 / 1000 /
2000 cells, up to 1.134 - 1.137 R\_p; EW change at most 1.2e-4 relative);
the user manual\textquotesingle s status table says so. Binary ca651768.
Suites: adv\_static\_limit 53, transit\_state 12, transit\_census 21,
physics\_probe 2199, thread\_count\_identity 9, all pass. Every \texttt{\_adv} golden
moves (5 short cases: T up to 5.9 \%, n(2\^{}3S) up to 9.4 \%). The catalog EWs
predate this change by -1 \% to -4 \%. Found, not changed:
\texttt{thermal\_damkohler\_number} is the net rate \textbar H - C\textbar/u, near zero at a
balanced cell, not the relaxation rate its comment describes.
Docs: the transit tool now says a stationary SOLVE writes no \texttt{\_adv} products
and its evaluation does; the project CLAUDE.md says the matrix compares the
\texttt{\_adv} pair as well.

\subsubsection{29.11 The lower face at the reservoir level}\label{2911-the-lower-face-at-the-reservoir-level}

The reservoir (the \texttt{base.inp} handoff or the Photochem profile\textquotesingle s matching
level) is stated at r = 1, but the domain\textquotesingle s lower face sat at
r\_edg(0) = 1 + dr/2, and \texttt{base\_boundary.f90} carried the reservoir pressure
up to that face along the reservoir\textquotesingle s own isentrope: a half-cell slab of
226 K gas whose thickness scales with dr lay between the level and the
domain (MEASURED: the lower layer offset by -2.2 \% / -1.15 \% in density on
500 / 1000 cells, \texttt{md/ew\_grid\_order\_20260926.md} section 3). Now the face
is the level in every grid mode (Uniform, Stretched, Mixed:
\texttt{define\_grid.f90}, r(0) = 2 - r(1), refusals if a ghost reaches r \textless= 0 or
the face is not 1), the face state is the stated reservoir with no
transport along any isentrope, the ghosts hold the reservoir continued
below the level, and the Photochem matching level is the face on every
grid; boundary model id \texttt{...\_face\_at\_level\_v4}. MEASURED on the molecular
He/H 1.8 state, 500 / 1000 / 2000 cells, all certified: the pressure gap
between the interior at the level and the reservoir falls from 2.93 /
1.44 / 0.70 \% to 0.114 / 0.067 / 0.038 \%; the lower-layer density offset
falls 20-40x; the wind quantities stay first order with errors 1.6-1.9x
smaller (log10 Mdot 7.9147 / 7.9113 / 7.9094 -\textgreater{} 7.9116 / 7.9098 / 7.9086;
\texttt{\_adv} EW 1.1548 / 1.1248 / 1.1091 -\textgreater{} 1.1482 / 1.1214 / 1.1074 \%A, measured
before the \texttt{\_adv} change of 29.10). The remaining first-order error is
likely the thermal layer at the contact, which the grid does not resolve
(cell-1 T 800 / 695 / 604 K against the 226 K reservoir, 13 \% lower at every
halving; not shown to drive the wind). Two other first-order terms were
found and left: the element-diffusion Dirichlet He/H is imposed at the
cell-1 center, dr/2 above the level (\texttt{binary\_element\_diffusion.f90}
946-959; the He/H change near the base is 1.5e-4), and the conduction bath
at the ghost center r(0) (\texttt{viscous\_conduction.f90} 485; conduction off in
these runs). Every cell center moves by dr/2, so the grid guard refuses
every state written before: such a state loads only as a seed made with
\texttt{map\_state\_to\_grid.py\ -\/-ic} onto the new grid (a 1-step cold run gives the
target grid). The regression fixtures and goldens are migrated in the same
series (MEASURED, 2026-09-27: every fixture of the default matrix and every
fixture a suite reads mapped with \texttt{map\_state\_to\_grid.py\ -\/-ic} onto a 1-step
cold-run grid; \texttt{wasp\_full\_newton/IC} re-solved and certified, info = 0,
\textbar\textbar R\textbar\textbar{} 1.43e-9, and \texttt{iontrans\_metals/IC} a copy of it; the \texttt{hp\_*} seeds,
whose recipe from \texttt{mol\_base\_handoff} is now refused by the \texttt{carrier} restart
token, mapped from their previous pairs with the seeds kept (x(H II) 0 and
1e-12); \texttt{jfnk\_hd189}, a named case outside the matrix whose legacy pair has
no \texttt{\#\ columns} header, not migrated; the two suites that read catalog states
now map them into their own scratch; every suite passes, and the standing
\texttt{residual\_determinism} row now passes (5.9e-2 against \textless{} 1); goldens
refreshed from the matrix run on 2dd9a865, backup
\texttt{golden\_pre\_facelevel\_20260927}); the LHS 1140 b catalog states stay
records of the previous
boundary (their EWs move by -0.6 \% at 500 cells, below the observational
error). Binary 2dd9a865.

\subsubsection{29.12 A restartable state after every outer pass}\label{2912-a-restartable-state-after-every-outer-pass}

The stationary route wrote its state only when the run ended, so a run
stopped from outside (a wall-clock \texttt{timeout}) lost every pass it had made:
MEASURED on 2026-09-28, three LHS 1140 b continuations lost 35, 13 and 9
passes. Now each outer pass that does not end the iteration writes the state
the next pass starts from, \texttt{output/Hydro\_ioniz\_last\_pass.txt} and
\texttt{output/Ion\_species\_last\_pass.txt}, overwritten each pass, by the same writer
and headers as the final state files, rebuilt from (u, f\_sp) in the order of
\texttt{rebuild\_state\_from\_checkpoint} with every side effect on the caloric mixture
undone; both files carry \texttt{certified=F\ cert\_reason=pass\_snapshot\_p\textless{}n\textgreater{}} and a
\texttt{\#\ pass\_snapshot} line (the movement bound and omega of that pass, which a
restart does not inherit: pass them by \texttt{EXHALE\_CARRIER\_TRUST} /
\texttt{EXHALE\_DIFF\_OMEGA}), so a pair mixed from two passes is refused by
\texttt{load\_IC} and no tool reads a snapshot as certified. Each file is written to
\texttt{\textless{}name\textgreater{}.part} and renamed. MEASURED: 0.03 - 0.05 s a pass on 500 cells
(passes of 90 - 135 s); a run killed during pass 3 left the complete pass-2
pair, which restarted and continued as the uninterrupted run did; the final
state files of a normal run are bitwise those of the previous binary.
\texttt{LHS1140b/models/.P1/s2\_continuation/config\_rung.sh} continues from the pair
when the parent has no final state. Binary a7255a03. The regression matrix
on it, strict (\texttt{REGRESSION\_REL\_TOL=0}): all 17 cases byte-identical to the
goldens of 2dd9a865 (MEASURED).

\subsubsection{29.13 The unsolved molecular cases by continuation (state of 2026-09-28 16:30)}\label{2913-the-unsolved-molecular-cases-by-continuation-state-of-2026-09-28-1630}

All runs under \texttt{LHS1140b/models/.P1/s2\_continuation/} (scripts \texttt{config\_rung.sh},
\texttt{make\_rung.sh} in \texttt{.P1/}, \texttt{chain.sh}, \texttt{facelevel\_continue.sh}; env as the
catalog: \texttt{EXHALE\_PTC\_DTAU0=1.0}, 8 threads). Certified states (MEASURED,
\texttt{certified=T}, the JFNK \textbar\textbar R\textbar\textbar{} of the last solve 6.5e-9 - 2.9e-8):

{\def\LTcaptype{none} % do not increment counter
{\small
\begin{longtable}[]{@{}>{\raggedright\arraybackslash}p{\dimexpr(\linewidth-10\tabcolsep)/5\relax}>{\raggedright\arraybackslash}p{\dimexpr(\linewidth-10\tabcolsep)/5\relax}>{\raggedright\arraybackslash}p{\dimexpr(\linewidth-10\tabcolsep)/5\relax}>{\raggedright\arraybackslash}p{\dimexpr(\linewidth-10\tabcolsep)/5\relax}>{\raggedright\arraybackslash}p{\dimexpr(\linewidth-10\tabcolsep)/5\relax}@{}}
\toprule\noalign{}
family & He/H & passes & binary & note \\
\midrule\noalign{}
\endhead
\bottomrule\noalign{}
\endlastfoot
kzz1e9 & 1.45 & 14 & 7db5afe9 & rung from 1.52 \\
kzz1e9 & 1.30 & 5 (after 24) & 7db5afe9 & \\
kzz1e9 & 1.223 & 9 (after 40 + 40) & 2dd9a865 & moved to the face-at-level grid \\
kzz1e9 & 1.186 & 24 (after 30) & 2dd9a865 & \\
kzz1e9 & 1.168 & 20 & 2dd9a865 & \\
kzz1e9 & 1.15 & 46 & a7255a03 & carrier bound held at 1e-2 \\
wellmixed & 2.13, 1.8, 1.4, 1.1, 0.85, 0.7603, 0.68, 0.6116 & 10 - 21 each & 7db5afe9 & path finding, previous grid \\
wellmixed & 0.6116, 0.55, 0.4 & 16, 18, 17 & 2dd9a865 & on the face-at-level grid \\
photochem kzzprofile & 9 & 41 & 2dd9a865 & after 40 passes on 6844f188 \\
\end{longtable}
}
}

Not solved yet: kzz1e9 below 1.15 (1.12 running, gated row falling to
1e-4 then flat at 1.3e-4 over passes 14-16); wellmixed 0.3347 and below;
photochem kzzprofile 2.09 (gated row oscillating 0.08 - 0.54, hydro info = 2
on most passes, about 40 min a pass); kzz1e9 0.083 and the wellmixed
0.083 are the far ends of the two chains.

What was learned on the way (MEASURED):

\begin{itemize}
\tightlist
\item
  The well-mixed molecular cases are the kzz1e9 ones with \texttt{He\_diffusion}
  off (their \texttt{base.inp} carries \texttt{Kzz\_base\ 1.0e9}, which the carrier
  transport reads); a kzz1e9 state with He/H set uniform
  (\texttt{map\_state\_to\_grid.py\ -\/-uniform}) certifies with diffusion off in 10
  passes, which is how the family was entered.
\item
  Runs stopped by a wall-clock \texttt{timeout} lost every pass (35, 13 and 9
  passes), which section 29.12 fixed for every later run; and
  \texttt{config\_rung.sh} ran under \texttt{set\ -e}, so a run that timed out never wrote
  its \texttt{EXIT.txt} and a waiting continuation never started (fixed).
\item
  The growing carrier bound of section 29.7 drives a sawtooth in the H2
  cases below He/H \textasciitilde1.2 (kzz1e9 1.15: the gated row between 5e-3 and 2e-1
  over 35 passes with hydro info alternating 0/2) and in the well-mixed
  0.3347 rung (between 2e-4 and 1e-2, the refusing cell moving between the
  front near cell 357 and the outer cell 500). Holding the bound at 1e-2
  (\texttt{EXHALE\_CARRIER\_TRUST\_HOLD=1e-2}) certified kzz1e9 1.15 in 46 passes, and
  1.12 then fell 2.6x a pass to 1e-4; it did not stop the well-mixed 0.3347
  oscillation (1e-3 - 2.7e-3 over 19 passes). The rule that grows the bound
  is therefore not right for the H2 carrier; an open item.
\end{itemize}

Documents brought to the code of this series (2026-09-28): the user manual
(input keys \texttt{Interdiffusion\ enthalpy\ flux}, \texttt{Outer\ shells}, \texttt{Composition\ update\ holds}, \texttt{H2\ double\ ionization}, \texttt{H2\ neutral\ dissociation} and the
retired \texttt{Oxygen\ transport} added from \texttt{known\_keys}; the \texttt{\_last\_pass} output
files; the face at the level; the BDF2 post-process; the progress control;
the thread-count note; stale statements corrected, among them the tokens a
\texttt{Restart\ option\ change:} line may not name, four not six), the physics and
algorithms document (the interdiffusion enthalpy flux in the energy
equations, the lower face, the grid constructions, the He+ + H pair, Ry/k,
the post-process step, the progress control and the snapshot; stale
statements corrected, among them the default of \texttt{H2\ double\ ionization},
\texttt{chung80}, and the claim that every extension defaults to off) and
\texttt{README.md}. Two source comments corrected (\texttt{base\_boundary.f90} model id,
\texttt{steady\_residual.f90} energy row scale; comments only), and
\texttt{src/tests/interdiffusion\_enthalpy\_flux/run.sh} made executable: \texttt{make\ test}
had skipped it silently.

\appendix
%% Appendices of the EXHALE update log.
%%
%% HAND-MAINTAINED, unlike Update_EXHALE.tex, which src/utils/update_log_to_tex.py
%% generates from Update_EXHALE.md.  The generator emits \appendix and then
%% %% Appendices of the EXHALE update log.
%%
%% HAND-MAINTAINED, unlike Update_EXHALE.tex, which src/utils/update_log_to_tex.py
%% generates from Update_EXHALE.md.  The generator emits \appendix and then
%% %% Appendices of the EXHALE update log.
%%
%% HAND-MAINTAINED, unlike Update_EXHALE.tex, which src/utils/update_log_to_tex.py
%% generates from Update_EXHALE.md.  The generator emits \appendix and then
%% \input{Update_EXHALE_appendix} after the converted body, so this fragment
%% carries no \appendix, no \documentclass and no \begin{document} /
%% \end{document} of its own.  Edit this file directly, then regenerate the
%% typeset log.
%%
%% Every number below is reprinted by
%%     python3 src/utils/codesize.py
%% from the repository root, together with the byte-identical list of
%% Table tab:untouched.  Rerun it and update this file in the same pass.
%%
%% Moved here on 2026-09-11 from docs/Update_EXHALE_stage1.tex, where the two
%% sections were the appendices of the stage-1 log.

% ===================================================================== %
\section{Code-size summary: EXHALE vs.\ the original ATES (2026-06-07; updated 2026-09-11)}
% ===================================================================== %
A line-by-line comparison of the Fortran source against the pristine upstream code
(\texttt{ATES/ATES-Code-main}) quantifies the extent of the in-house extension.
\emph{Pure code lines (SLOC)} exclude blank lines and comments (the comment
stripper is string-aware, so a \texttt{!} inside a quoted string is not treated as
a comment); the line-by-line diff additionally normalizes whitespace and case, so
reformatting and comment-only edits are \emph{not} counted as changes. The
measurement script is \texttt{src/utils/codesize.py}, which reprints every
number of this appendix and the byte-identical list of
Table~\ref{tab:untouched}. Measured 2026-09-11, i.e.\ after the whole of
stage 1 and the stage-2 series through item P19 of section 8: the acceptance
and certification contracts, the attempted-step controller, the partitioned
stationary solver, the oxygen and carbon chemistry, the molecular infrared
bands and the test suites those series brought with them.

\paragraph{What the script counts as \emph{core}.} Every \texttt{.f90} file
under \texttt{src/} except those in \texttt{src/modules/wind\_ae/}, which is
the ported Wind-AE IC generator and is listed separately. That rule puts the
standalone drivers under \texttt{src/tests/} in the core figure, and they are
now a large part of it, so the table splits the core row into the production
build and those drivers. Only the production files are compiled into
\texttt{EXHALE.x}; each test suite builds its own driver from
\texttt{src/tests/<suite>/run.sh}.

\begin{center}\small
\begin{tabular}{@{}lcc@{}}
\toprule
 & Fortran files & code lines (SLOC) \\
\midrule
Original ATES (\texttt{ATES-Code-main}) & 47 & 3\,514 \\
EXHALE core (excl.\ Wind-AE port)       & 169 & 71\,503 \\
\quad of which the production build     & 95 & 50\,762 ($\approx14.4\times$) \\
\quad of which the \texttt{src/tests/} drivers & 74 & 20\,741 \\
\midrule
net change (core) & $+122$ & $+67\,989$ ($\approx20.3\times$) \\
\midrule
Wind-AE IC-generator port (\texttt{wind\_ae/}) & $+36$ & $+4\,756$ \\
EXHALE total & 205 & 76\,259 ($\approx21.7\times$) \\
\bottomrule
\end{tabular}
\end{center}

\noindent
\textbf{Line-by-line diff, core} (comments + blanks excluded; whitespace/case
normalized): \textbf{69\,179} code lines added in EXHALE (new modules + the new
side of modified lines) and \textbf{1\,190} removed from the original (deleted + the
old side of modified lines), for \textbf{70\,369} total changed lines. At the file
level, of the 47 original files \textbf{38 were modified}, \textbf{8 are
unchanged} (all eight byte-for-byte; the list is Table~\ref{tab:untouched}) and
\textbf{one was removed}: \texttt{flux/\allowbreak speed\_\allowbreak estimate\_\allowbreak HLLC.f90}, Toro's
adaptive $p_\star$ estimate, was \texttt{use}d by \texttt{Num\_\allowbreak Fluxes.f90} and
never called, because the HLLC branch builds its own Davis/Einfeldt speeds
inline. \textbf{123 new core modules} were added. Of those, 49 are part of the
simulation build:

{\footnotesize
\begin{longtable}{@{}>{\raggedright\arraybackslash}p{0.36\linewidth} >{\raggedright\arraybackslash}p{0.50\linewidth} r@{}}
\toprule
New module & Purpose & SLOC \\
\midrule
\endfirsthead
\toprule
New module & Purpose & SLOC \\
\midrule
\endhead
\texttt{time\_\allowbreak step/\allowbreak steady\_\allowbreak newton.f90} & stationary solve of the hydrodynamic rows: residual $F(Y)=0$ on the physical cells, Newton/Krylov with pseudo-transient continuation, line search, best-iterate ledger & 9435 \\
\texttt{lower\_\allowbreak atmosphere/\allowbreak diffusive\_\allowbreak photochemistry.f90} & vertical transport of the molecular carriers (H$_2$, OH, H$_2$O, CO) solved together with the chemistry that makes and destroys them & 2485 \\
\texttt{lower\_\allowbreak atmosphere/\allowbreak molecular\_\allowbreak infrared\_\allowbreak data.f90} & band and line data of the molecular infrared coolant (generated by \texttt{cooling\_data/molecular\_infrared\_bands.py}) & 2182 \\
\texttt{functions/\allowbreak binary\_\allowbreak element\_\allowbreak diffusion.f90} & conservative binary H/He element transport (stage-resolved pair coefficients, donor-cell drift, metals slaved to H) & 1898 \\
\texttt{time\_\allowbreak step/\allowbreak attempted\_\allowbreak step.f90} & the attempted-step controller: one checkpoint, one trial, one adoption boundary, and the rollback contract & 1608 \\
\texttt{time\_\allowbreak step/\allowbreak certification.f90} & one evaluator of the equations a run actually solves, and the contexts that read its verdict & 1273 \\
\texttt{nonlinear\_\allowbreak system\_\allowbreak solver/\allowbreak constrained\_\allowbreak chemical\_\allowbreak equilibrium.f90} & chemical equilibrium of the molecular H/He network in logarithmic number densities, element conservation as explicit residual rows & 1173 \\
\texttt{lower\_\allowbreak atmosphere/\allowbreak h2\_\allowbreak self\_\allowbreak shielding\_\allowbreak table.f90} & H$_2$ Lyman--Werner dissociation cross section with overlapping lines, and its two branching ratios & 952 \\
\texttt{time\_\allowbreak step/\allowbreak energy\_\allowbreak semi\_\allowbreak implicit.f90} & implicit source step for the thermal energy of a cell, at fixed composition and heating & 545 \\
\texttt{time\_\allowbreak step/\allowbreak steady\_\allowbreak residual.f90} & finite-volume stationary residual $R = \mathrm{d}u/\mathrm{d}t$ and its norms, shared by the diagnostic, the convergence monitor and the stationary solver & 509 \\
\texttt{time\_\allowbreak step/\allowbreak viscous\_\allowbreak conduction.f90} & Navier--Stokes viscous momentum force, its dissipation, and thermal conduction (Crank--Nicolson) & 504 \\
\texttt{radiation/\allowbreak excited\_\allowbreak hydrogen.f90} & non-LTE H($n{=}2$) 2s/2p equilibrium (Christie, Arras \& Li 2013) and the Balmer feedback & 474 \\
\texttt{lower\_\allowbreak atmosphere/\allowbreak element\_\allowbreak inventory.f90} & one stoichiometric map of the H, He, O and C nuclei of a cell, and the three constraints a composition state is judged by & 448 \\
\texttt{radiation/\allowbreak charge\_\allowbreak exchange.f90} & charge-exchange network (Huang et al.\ 2023, Table 4) & 446 \\
\texttt{files\_\allowbreak IO/\allowbreak lower\_\allowbreak atmosphere\_\allowbreak profile.f90} & lower-atmosphere profile reader: log-$p$ interpolation, elemental reservoirs, $K_{zz}(p)$ & 367 \\
\texttt{states/\allowbreak caloric\_\allowbreak eos.f90} & caloric equation of state: internal energy density, pressure and temperature, with the H$_2$ level ladder & 340 \\
\texttt{lower\_\allowbreak atmosphere/\allowbreak molecular\_\allowbreak reaction\_\allowbreak heat.f90} & net chemical heat of the collisional reactions of the H$_2$/He network & 339 \\
\texttt{lower\_\allowbreak atmosphere/\allowbreak oxygen\_\allowbreak rates.f90} & rate coefficients and thermodynamic data of the oxygen-hydrogen chemistry, with their validity ranges & 316 \\
\texttt{functions/\allowbreak element\_\allowbreak census.f90} & one definition of the elemental and charge content of a composition state, and the invariant checks built on it & 307 \\
\texttt{nonlinear\_\allowbreak system\_\allowbreak solver/\allowbreak System\_\allowbreak HeH\_\allowbreak mol.f90} & coupled H/He $+$ molecular (H$_2$/H$_2^+$/H$_3^+$/HeH$^+$) system, turnover-scaled rows & 303 \\
\texttt{radiation/\allowbreak electron\_\allowbreak energy\_\allowbreak degradation.f90} & degradation of a fast photoelectron in a partly ionized H/H$_2$/He gas: the heat fraction and the ionizations of each target & 287 \\
\texttt{radiation/\allowbreak lya\_\allowbreak rt.f90} & Ly$\alpha$ transfer by the escape-probability closure, evaluated in line each step & 257 \\
\texttt{lower\_\allowbreak atmosphere/\allowbreak molecular\_\allowbreak infrared\_\allowbreak cooling.f90} & infrared exchange of the molecular layer: H$_2$O and CO vibration-rotation bands, H$_2$ quadrupole and magnetic-dipole lines & 231 \\
\texttt{states/\allowbreak base\_\allowbreak boundary.f90} & the lower boundary as a characteristic condition at the face, the number of conditions set by the Mach number there & 199 \\
\texttt{functions/\allowbreak h2\_\allowbreak photo\_\allowbreak channels.f90} & H$_2$ photoabsorption resolved into mutually exclusive final-state channels, with one stoichiometric table & 199 \\
\texttt{lower\_\allowbreak atmosphere/\allowbreak lyman\_\allowbreak werner.f90} & H$_2$ photodissociation in the Lyman and Werner bands (the Solomon process) & 195 \\
\texttt{radiation/\allowbreak opacity\_\allowbreak models.f90} & opacity-model dispatcher for the photoionization cross sections (A/C/P/T) & 180 \\
\texttt{files\_\allowbreak IO/\allowbreak metals\_\allowbreak input\_\allowbreak read.f90} & parse \texttt{metals.inp}: trace-metal abundances at runtime & 177 \\
\texttt{functions/\allowbreak composition.f90} & single point turning $(\rho, f_{\rm sp})$ into every species density, $n_e$ and $n_{\rm tot}$, and converting pressure and temperature & 162 \\
\texttt{flux/\allowbreak species\_\allowbreak face\_\allowbreak flux.f90} & advection of the species mass fractions on the same faces, with the same mass flux, as the density & 162 \\
\texttt{time\_\allowbreak step/\allowbreak hydrodynamic\_\allowbreak rows.f90} & the hydrodynamic rows of the stationary residual in double precision, instantiated from the kind-generic text & 161 \\
\texttt{nonlinear\_\allowbreak system\_\allowbreak solver/\allowbreak System\_\allowbreak HeH\_\allowbreak metals.f90} & coupled H/He$+$metals ionization ($+$\,analytic Jacobian) & 150 \\
\texttt{lower\_\allowbreak atmosphere/\allowbreak mol\_\allowbreak rates.f90} & molecular reaction rates (Koskinen et al.\ 2022, Table 1) & 150 \\
\texttt{lower\_\allowbreak atmosphere/\allowbreak h3p\_\allowbreak cooling.f90} & optically thin H$_3^+$ infrared cooling (Miller et al.\ 2013) & 148 \\
\texttt{nonlinear\_\allowbreak system\_\allowbreak solver/\allowbreak System\_\allowbreak HeH\_\allowbreak mol\_\allowbreak metals.f90} & molecular $+$ metals merged ionization system & 141 \\
\texttt{nonlinear\_\allowbreak system\_\allowbreak solver/\allowbreak newton\_\allowbreak solver.f90} & damped Newton with an analytic Jacobian for the ionization systems, with the MINPACK \texttt{hybrd1} retry & 128 \\
\texttt{init/\allowbreak species\_\allowbreak table.f90} & metadata table of the metal stages and of the base, molecular and oxygen species & 128 \\
\texttt{lower\_\allowbreak atmosphere/\allowbreak co\_\allowbreak photodissociation.f90} & CO photodissociation on the 912--1201\,\AA\ band & 127 \\
\texttt{nonlinear\_\allowbreak system\_\allowbreak solver/\allowbreak ion\_\allowbreak residual\_\allowbreak core.f90} & shared ionization-residual rows and analytic-Jacobian pieces of the \texttt{System\_*} family & 127 \\
\texttt{lower\_\allowbreak atmosphere/\allowbreak water\_\allowbreak photolysis.f90} & H$_2$O and OH photodissociation in the four far-ultraviolet bands of the oxygen option & 111 \\
\texttt{lower\_\allowbreak atmosphere/\allowbreak co\_\allowbreak self\_\allowbreak shielding\_\allowbreak table.f90} & CO self-shielding function $\Theta[N({\rm CO}), N({\rm H_2})]$ (Visser, van Dishoeck \& Black 2009) & 110 \\
\texttt{nonlinear\_\allowbreak system\_\allowbreak solver/\allowbreak ion\_\allowbreak cell\_\allowbreak state.f90} & named-field cell state of the ionization residuals, in place of the position-packed \texttt{params} & 93 \\
\texttt{lower\_\allowbreak atmosphere/\allowbreak lower\_\allowbreak column.f90} & analytic lower-atmosphere column from the 1-bar level to the base pressure (Koskinen et al.\ 2022) & 90 \\
\texttt{flux/\allowbreak low\_\allowbreak mach\_\allowbreak dissipation.f90} & gated fourth-difference stress damping the two-cell velocity mode a contact-resolving flux leaves undamped & 87 \\
\texttt{files\_\allowbreak IO/\allowbreak opacity\_\allowbreak input\_\allowbreak read.f90} & parse \texttt{opacity.inp} & 85 \\
\texttt{radiation/\allowbreak hydrogen\_\allowbreak n2\_\allowbreak rates.f90} & atomic data and rate coefficients of the H $n{=}2$ manifold and of the Ly$\alpha$ line, in one place & 76 \\
\texttt{nonlinear\_\allowbreak system\_\allowbreak solver/\allowbreak System\_\allowbreak HeH\_\allowbreak TR\_\allowbreak metals.f90} & merged He\,2$^3$S$+$metals system & 68 \\
\texttt{lower\_\allowbreak atmosphere/\allowbreak h2\_\allowbreak vibrational\_\allowbreak relaxation.f90} & the fate and the size of the vibrational energy left in an H$_2$ molecule, for both of its drivers & 55 \\
\texttt{nonlinear\_\allowbreak system\_\allowbreak solver/\allowbreak constrained\_\allowbreak equilibrium\_\allowbreak probe.f90} & standalone re-evaluation of one saved constrained-equilibrium state, without the hydrodynamics & 15 \\
\midrule
\multicolumn{2}{r}{production new-module total} & 30\,003 \\
\bottomrule
\end{longtable}
}

\noindent
The other 74 are the standalone test drivers under \texttt{src/tests/}, which
did not exist in ATES at all. They state what a routine is allowed to return
and are run by suite, not by the production binary:

{\footnotesize
\begin{longtable}{@{}>{\raggedright\arraybackslash}p{0.36\linewidth} >{\raggedright\arraybackslash}p{0.50\linewidth} r@{}}
\toprule
Test driver & What it asserts & SLOC \\
\midrule
\endfirsthead
\toprule
Test driver & What it asserts & SLOC \\
\midrule
\endhead
\texttt{tests/\allowbreak krylov\_\allowbreak and\_\allowbreak dogleg/\allowbreak krylov\_\allowbreak and\_\allowbreak dogleg\_\allowbreak tests.f90} & what the linear solve and the step geometry of the stationary solver are allowed to return & 2675 \\
\texttt{tests/\allowbreak steady\_\allowbreak species\_\allowbreak rows/\allowbreak steady\_\allowbreak species\_\allowbreak rows\_\allowbreak tests.f90} & the unknown space of the stationary system, one row per transported balance & 1806 \\
\texttt{tests/\allowbreak diffusion\_\allowbreak tests.f90} & the binary H/He element-diffusion operator: analytic limits and conservation & 1051 \\
\texttt{tests/\allowbreak grid\_\allowbreak and\_\allowbreak gates/\allowbreak hydrostatic\_\allowbreak residual.f90} & the momentum row returns zero on a hydrostatic state & 967 \\
\texttt{tests/\allowbreak carrier\_\allowbreak retry/\allowbreak carrier\_\allowbreak retry.f90} & the carrier-local checkpoint and the rejection-and-retry controller, on a real molecular column & 842 \\
\texttt{tests/\allowbreak certification/\allowbreak certification\_\allowbreak contexts.f90} & the certification evaluator and its contexts, against the predicates they replace & 790 \\
\texttt{tests/\allowbreak species\_\allowbreak face\_\allowbreak flux/\allowbreak species\_\allowbreak face\_\allowbreak flux\_\allowbreak tests.f90} & species advection on the face mass fluxes, and the element rows of the diffusion operator & 557 \\
\texttt{tests/\allowbreak physics\_\allowbreak probe/\allowbreak weno3\_\allowbreak reconstruction\_\allowbreak order.f90} & convergence order of the face reconstruction against the exact face value & 482 \\
\texttt{tests/\allowbreak element\_\allowbreak operator/\allowbreak element\_\allowbreak operator\_\allowbreak tests.f90} & the boundary the element transport operator is measured under & 471 \\
\texttt{tests/\allowbreak spectrum\_\allowbreak type/\allowbreak loaded\_\allowbreak sed\_\allowbreak table\_\allowbreak semantics.f90} & what a row of a loaded SED table means, and what the production consumers make of it & 417 \\
\texttt{tests/\allowbreak attempted\_\allowbreak step/\allowbreak attempted\_\allowbreak step\_\allowbreak tests.f90} & the checkpoint and the step-size policy of the attempted-step controller & 415 \\
\texttt{tests/\allowbreak a2\_\allowbreak m1/\allowbreak a2\_\allowbreak m1\_\allowbreak rate\_\allowbreak check.f90} & the oxygen rate coefficients against their published sources & 406 \\
\texttt{tests/\allowbreak spectrum\_\allowbreak type/\allowbreak balmer\_\allowbreak band\_\allowbreak quadrature.f90} & quadrature of the Balmer-continuum band integrals of \texttt{excited\_hydrogen} & 367 \\
\texttt{tests/\allowbreak energy\_\allowbreak update/\allowbreak conduction\_\allowbreak floor\_\allowbreak tests.f90} & what the Crank--Nicolson transport stage returns, refuses and conserves & 360 \\
\texttt{tests/\allowbreak physics\_\allowbreak probe/\allowbreak metal\_\allowbreak photoionization\_\allowbreak fits.f90} & every metal photoionization cross section against the published fit, and only where it is valid & 343 \\
\texttt{tests/\allowbreak element\_\allowbreak census\_\allowbreak tests.f90} & the elemental and charge invariants of a composition state & 340 \\
\texttt{tests/\allowbreak energy\_\allowbreak update/\allowbreak energy\_\allowbreak update\_\allowbreak tests.f90} & the residual-controlled temperature solve of the implicit energy source step & 334 \\
\texttt{tests/\allowbreak spectrum\_\allowbreak type/\allowbreak loaded\_\allowbreak sed\_\allowbreak threshold\_\allowbreak edges.f90} & the photon grid built from a loaded SED, and the rates formed on it & 304 \\
\texttt{tests/\allowbreak physics\_\allowbreak probe/\allowbreak species\_\allowbreak formation\_\allowbreak energy\_\allowbreak table.f90} & the formation-energy table of the energy ledger, and the reaction energies that are its differences & 275 \\
\texttt{tests/\allowbreak physics\_\allowbreak probe/\allowbreak lya\_\allowbreak band\_\allowbreak transmission.f90} & the stellar Ly$\alpha$ band reaching a cell through the atomic hydrogen column above it & 272 \\
\texttt{tests/\allowbreak spectrum\_\allowbreak type/\allowbreak wasp121b\_\allowbreak sed\_\allowbreak sensitivity.f90} & the production band quantities of the three WASP-121\,b spectrum candidates & 262 \\
\texttt{tests/\allowbreak coupled\_\allowbreak source\_\allowbreak step/\allowbreak coupled\_\allowbreak source\_\allowbreak tests.f90} & the coupled temperature-composition source step & 260 \\
\texttt{tests/\allowbreak physics\_\allowbreak probe/\allowbreak lyman\_\allowbreak werner\_\allowbreak cell\_\allowbreak mean.f90} & the Lyman--Werner rate a cell carries is the mean of the local rate over that cell & 234 \\
\texttt{tests/\allowbreak physics\_\allowbreak probe/\allowbreak lya\_\allowbreak beam\_\allowbreak cell\_\allowbreak mean.f90} & the stellar Ly$\alpha$ field of a cell is the beam mean over the cell's own optical depth & 229 \\
\texttt{tests/\allowbreak spectrum\_\allowbreak type/\allowbreak balmer\_\allowbreak continuum\_\allowbreak field.f90} & the two scalar Balmer-continuum rates of H($n{=}2$) & 227 \\
\texttt{tests/\allowbreak physics\_\allowbreak probe/\allowbreak positivity\_\allowbreak limiter\_\allowbreak scaling.f90} & the positivity scaling of the reconstructed face states (Zhang \& Shu 2010) & 223 \\
\texttt{tests/\allowbreak carrier\_\allowbreak returned\_\allowbreak state\_\allowbreak acceptance/\allowbreak carrier\_\allowbreak returned\_\allowbreak state\_\allowbreak acceptance.f90} & what makes a returned carrier state a solution, and what does not & 222 \\
\texttt{tests/\allowbreak grid\_\allowbreak and\_\allowbreak gates/\allowbreak photon\_\allowbreak grid\_\allowbreak threshold\_\allowbreak edges.f90} & every active absorber's threshold against the bin partition of the photon grid & 213 \\
\texttt{tests/\allowbreak physics\_\allowbreak probe/\allowbreak photoionization\_\allowbreak field\_\allowbreak substitution.f90} & the order in which the field and the composition that attenuates it are brought into agreement & 209 \\
\texttt{tests/\allowbreak adv\_\allowbreak static\_\allowbreak limit/\allowbreak stationarity\_\allowbreak by\_\allowbreak the\_\allowbreak mass\_\allowbreak row.f90} & which test decides that the state the post-process was handed is stationary in a cell & 199 \\
\texttt{tests/\allowbreak physics\_\allowbreak probe/\allowbreak co\_\allowbreak destruction.f90} & the CO destruction rates, the shielding table, the deposited energies and the cell mean & 195 \\
\texttt{tests/\allowbreak physics\_\allowbreak probe/\allowbreak recombination\_\allowbreak coefficient\_\allowbreak fits.f90} & every recombination coefficient against the published fit & 189 \\
\texttt{tests/\allowbreak spectrum\_\allowbreak type/\allowbreak photon\_\allowbreak grid\_\allowbreak n2\_\allowbreak floor.f90} & the photon grid built when the excited-hydrogen coupling is armed & 184 \\
\texttt{tests/\allowbreak carrier\_\allowbreak reference\_\allowbreak scales/\allowbreak carrier\_\allowbreak reference\_\allowbreak scales.f90} & the element a transported carrier is made of, and whether its chemistry rows can be differentiated & 181 \\
\texttt{tests/\allowbreak carrier\_\allowbreak constraint\_\allowbreak attribution/\allowbreak carrier\_\allowbreak constraint\_\allowbreak attribution.f90} & which residuals belong to the solver and which to a constraint & 179 \\
\texttt{tests/\allowbreak adv\_\allowbreak static\_\allowbreak limit/\allowbreak enthalpy\_\allowbreak flux\_\allowbreak of\_\allowbreak mass\_\allowbreak divergence.f90} & the third term of the advected energy equation & 178 \\
\texttt{tests/\allowbreak grid\_\allowbreak and\_\allowbreak gates/\allowbreak free\_\allowbreak outflow\_\allowbreak boundary.f90} & the outer free-outflow boundary on a supersonic wind, and its ghost state & 174 \\
\texttt{tests/\allowbreak spectrum\_\allowbreak type/\allowbreak planck\_\allowbreak photon\_\allowbreak field.f90} & the incident flux put on the photon grid for the photospheric blackbody type & 173 \\
\texttt{tests/\allowbreak a2\_\allowbreak m2/\allowbreak a2\_\allowbreak m2\_\allowbreak kinetics\_\allowbreak check.f90} & the kinetics quantities the oxygen option adds on top of the rate audit & 163 \\
\texttt{tests/\allowbreak physics\_\allowbreak probe/\allowbreak lya\_\allowbreak escape\_\allowbreak probability.f90} & the escape probability of the two-level closure against the published slab solution & 163 \\
\texttt{tests/\allowbreak e1\_\allowbreak h2/\allowbreak e1\_\allowbreak h2\_\allowbreak channel\_\allowbreak check.f90} & the mutually exclusive H$_2$ photoabsorption channels & 161 \\
\texttt{tests/\allowbreak adv\_\allowbreak static\_\allowbreak limit/\allowbreak upstream\_\allowbreak caloric\_\allowbreak energy.f90} & the upstream energy of the advected energy equation & 154 \\
\texttt{tests/\allowbreak physics\_\allowbreak probe/\allowbreak photoionization\_\allowbreak field\_\allowbreak self\_\allowbreak consistency.f90} & the field one cell is solved in, against the composition that cell returns & 148 \\
\texttt{tests/\allowbreak grid\_\allowbreak and\_\allowbreak gates/\allowbreak photon\_\allowbreak grid\_\allowbreak quadrature.f90} & the photoionization rates of one unattenuated cell against the band quadrature & 146 \\
\texttt{tests/\allowbreak physics\_\allowbreak probe/\allowbreak co\_\allowbreak helium\_\allowbreak ion\_\allowbreak sink.f90} & He$^+$ $+$ CO as a species sink, not only as a heat source & 145 \\
\texttt{tests/\allowbreak physics\_\allowbreak probe/\allowbreak riemann\_\allowbreak wave\_\allowbreak speeds.f90} & signal speeds at a face, and the state one conservative update with that flux leaves behind & 143 \\
\texttt{tests/\allowbreak physics\_\allowbreak probe/\allowbreak helium\_\allowbreak level\_\allowbreak cooling\_\allowbreak ledger.f90} & which helium level each cooling channel of \texttt{eval\_cool} is charged to & 139 \\
\texttt{tests/\allowbreak physics\_\allowbreak probe/\allowbreak heating\_\allowbreak channel\_\allowbreak closure.f90} & one assembly of the volumetric heating rate, and the closure of its channels & 136 \\
\texttt{tests/\allowbreak acceptance\_\allowbreak classes/\allowbreak acceptance\_\allowbreak class\_\allowbreak contract.f90} & one meaning of the ionization acceptance classes for every consumer & 132 \\
\texttt{tests/\allowbreak physics\_\allowbreak probe/\allowbreak opacity\_\allowbreak model\_\allowbreak p\_\allowbreak ledger.f90} & the radiative ledger of a cell whose cross sections are broadened by opacity model P & 124 \\
\texttt{tests/\allowbreak physics\_\allowbreak probe/\allowbreak photoevent\_\allowbreak energy\_\allowbreak ledger.f90} & the energy an isolated photoabsorption event has to account for & 117 \\
\texttt{tests/\allowbreak physics\_\allowbreak probe/\allowbreak xuv\_\allowbreak cell\_\allowbreak mean\_\allowbreak attenuation.f90} & the beam attenuation the photoionization rate of a cell is evaluated with & 115 \\
\texttt{tests/\allowbreak grid\_\allowbreak and\_\allowbreak gates/\allowbreak grid\_\allowbreak window\_\allowbreak indices.f90} & the two index statements that decide which cells a quantity is read from & 101 \\
\texttt{tests/\allowbreak physics\_\allowbreak probe/\allowbreak h2\_\allowbreak rovibrational\_\allowbreak identity.f90} & the identities tying the H$_2$ partition function of the chemistry to the internal energy of the caloric EOS & 96 \\
\texttt{tests/\allowbreak adv\_\allowbreak static\_\allowbreak limit/\allowbreak stationary\_\allowbreak flux\_\allowbreak refusal.f90} & the sensitivity of the advected energy correction to a non-stationary mass flux & 91 \\
\texttt{tests/\allowbreak physics\_\allowbreak probe/\allowbreak ionization\_\allowbreak threshold\_\allowbreak turn\_\allowbreak on.f90} & one definition of each ionization threshold, and every cross section turning on exactly there & 90 \\
\texttt{tests/\allowbreak grid\_\allowbreak and\_\allowbreak gates/\allowbreak grid\_\allowbreak width\_\allowbreak identity.f90} & the cell-width and total-width identities of the radial grid & 83 \\
\texttt{tests/\allowbreak physics\_\allowbreak probe/\allowbreak h2\_\allowbreak channel\_\allowbreak detector\_\allowbreak ratio.f90} & the measured proton-to-H$_2^+$ ratio of Chung et al.\ (1993) against the four channels & 83 \\
\texttt{tests/\allowbreak physics\_\allowbreak probe/\allowbreak fine\_\allowbreak structure\_\allowbreak escape\_\allowbreak probability.f90} & the escape probability of the ground-term fine-structure lines, and the depth it is fed & 80 \\
\texttt{tests/\allowbreak physics\_\allowbreak probe/\allowbreak atomic\_\allowbreak mass\_\allowbreak and\_\allowbreak radius\_\allowbreak constants.f90} & the helium mass and the astronomical constants of the input readers, one definition each & 79 \\
\texttt{tests/\allowbreak physics\_\allowbreak probe/\allowbreak h3p\_\allowbreak cooling\_\allowbreak limits.f90} & the limits, the continuities and the published values of the H$_3^+$ cooling rate & 79 \\
\texttt{tests/\allowbreak physics\_\allowbreak probe/\allowbreak assertion\_\allowbreak report.f90} & one printed verdict line per assertion, and the failure count the driver exits on & 74 \\
\texttt{tests/\allowbreak adv\_\allowbreak static\_\allowbreak limit/\allowbreak thermal\_\allowbreak damkohler\_\allowbreak limit.f90} & the validity condition of the advection-corrected energy solve, as a function of the flow speed & 72 \\
\texttt{tests/\allowbreak a2\_\allowbreak m2/\allowbreak a2\_\allowbreak m2\_\allowbreak g3\_\allowbreak run\_\allowbreak check.f90} & with every photolysis rate at zero, the solved O/OH/H$_2$O partition is the chemical equilibrium & 71 \\
\texttt{tests/\allowbreak physics\_\allowbreak probe/\allowbreak absorbed\_\allowbreak fraction\_\allowbreak switch.f90} & $(1-e^{-d})/d$, the fraction of the photons entering a cell that the cell absorbs & 71 \\
\texttt{tests/\allowbreak constrained\_\allowbreak network\_\allowbreak layout/\allowbreak transported\_\allowbreak proton\_\allowbreak constraint\_\allowbreak layout.f90} & which species are unknowns of the continuation, and whether the system built is square & 71 \\
\texttt{tests/\allowbreak test\_\allowbreak columns.f90} & synthetic columns whose species carry the mass density they are measured against & 67 \\
\texttt{tests/\allowbreak physics\_\allowbreak probe/\allowbreak h2\_\allowbreak infrared\_\allowbreak line\_\allowbreak populations.f90} & the LTE populations weighting the H$_2$ quadrupole lines are those of the caloric EOS ladder & 65 \\
\texttt{tests/\allowbreak physics\_\allowbreak probe/\allowbreak wind\_\allowbreak ae\_\allowbreak bridge\_\allowbreak composition.f90} & the composition an EXHALE \texttt{input.inp} hands the Wind-AE solver & 54 \\
\texttt{tests/\allowbreak physics\_\allowbreak probe/\allowbreak voronov\_\allowbreak carbon\_\allowbreak ionization.f90} & the electron-impact ionization of C$^0$ and C$^+$ against the published Voronov (1997) rows & 51 \\
\texttt{tests/\allowbreak physics\_\allowbreak probe/\allowbreak water\_\allowbreak photolysis\_\allowbreak lyman\_\allowbreak alpha\_\allowbreak yields.f90} & the H$_2$O branching on the band that contains Ly$\alpha$, against the published yields & 49 \\
\texttt{tests/\allowbreak physics\_\allowbreak probe/\allowbreak h3p\_\allowbreak recombination\_\allowbreak branching.f90} & the two product channels of H$_3^+$ dissociative recombination & 44 \\
\texttt{tests/\allowbreak energy\_\allowbreak update/\allowbreak cooling\_\allowbreak models.f90} & analytic cooling laws for the tests of the implicit energy update & 43 \\
\texttt{tests/\allowbreak species\_\allowbreak masses/\allowbreak carrier\_\allowbreak background\_\allowbreak masses.f90} & the mass the carrier transport gives each of its background collision partners & 36 \\
\midrule
\multicolumn{2}{r}{test-driver total} & 20\,741 \\
\midrule
\multicolumn{2}{r}{new-module total} & 50\,744 \\
\bottomrule
\end{longtable}
}

\noindent
So the 123 new core modules contribute 50\,744 of the 69\,179 added code lines;
the remaining 18\,435 added (and 1\,121 of the 1\,190 removed, the other 69
being the deleted HLLC estimator) are edits spread across the 38 modified
files. The ten heaviest, by added code lines, are the main driver
(\texttt{EXHALE\_\allowbreak main.f90}, 3\,498 added, 246 $\to$ 3\,625 SLOC: the marching
loop, the attempted-step adoption boundary, the run-mode contracts and the
written-state assertions), the ionization equilibrium and its diagnostics
(\texttt{ionization\_\allowbreak equilibrium.f90} 2\,274 and \texttt{util\_\allowbreak ion\_\allowbreak eq.f90}
1\,803: the metal, molecular and oxygen networks, the acceptance classes, the
photoionization field and the heating channels), the cooling module
(\texttt{Cool\_\allowbreak coeff.f90}, 2\,250 added, 442 $\to$ 3\,966 raw lines: CHIANTI
metal-line fits, Fe tables, fine-structure statistical equilibrium, Penning
and molecular rate coefficients), input parsing
(\texttt{input\_\allowbreak read.f90}, 1\,753), the state reader and writer
(\texttt{load\_\allowbreak IC.f90} 1\,063, \texttt{write\_\allowbreak output.f90} 928,
\texttt{write\_\allowbreak setup\_\allowbreak report.f90} 751), the advection-corrected post-processor
(\texttt{post\_\allowbreak process\_\allowbreak adv.f90}, 647) and the cross sections
(\texttt{cross\_\allowbreak sec.f90}, 438). Those ten carry 15\,405 of the 18\,435. The rest
of the edits are in the reconstruction, the fluxes, the boundary conditions,
the grid, the time step, the energy equation and the \texttt{System\_*}
ionization systems, which wire the metals, the molecules, the opacity
dispatcher, the Roche potential, the transonic/auto/Wind-AE IC, the
excited-H and Ly$\alpha$ physics, the element diffusion and the Brent/Newton
solver upgrades into the original solver.

In total the EXHALE core is now $\approx$20.3$\times$ the original code base
(3\,514\,$\to$\,71\,503 SLOC), or $\approx$14.4$\times$ counting only the
production build (50\,762 SLOC), touching almost all of the original files (38
of 47 modified, one deleted, plus 123 new); with the ported Wind-AE IC
generator the full Fortran tree is 205 files / 76\,259 SLOC
($\approx$21.7$\times$). Beyond Fortran, the tree also carries the Python
tooling (GUI, the \texttt{EXHALE\_transit.py} transmission post-processor,
plotting/analysis, the lower-atmosphere adapters and the elemental-flux
closure driver) and 181\,816 lines of documentation (187 md/tex files: the
memos, the manual and this log) under \texttt{docs/}.


% ===================================================================== %
\section{Fortran source inherited unchanged from ATES}\label{sec:untouched}
% ===================================================================== %
Of the 47 upstream Fortran files (\texttt{ATES/ATES-Code-main}), \textbf{8 are
byte-for-byte identical} in EXHALE (re-measured 2026-09-11): the numerics
EXHALE reuses without touching, which are now the linear-algebra part of the
included \textsc{MINPACK} nonlinear solver plus the gravity pre-evaluation.
Table~\ref{tab:untouched} lists them with their role. No file is identical
after whitespace and case normalization but not byte-identical, so the
normalized-unchanged set and the byte-identical set are the same eight.

Five files left the byte-identical set since the 2026-08-30 edition of this
table, each for a documented reason rather than drift.

\begin{itemize}\itemsep2pt
\item \texttt{nonlinear\_\allowbreak system\_\allowbreak solver/\allowbreak hybrd.f90},
      \texttt{hybrd1.f90} and \texttt{fdjac1.f90} carry a new
      \texttt{unit\_step\_floor} argument threaded from \texttt{hybrd} to
      \texttt{fdjac1}, which selects the forward-difference step of the
      Jacobian: \texttt{.false.} is MINPACK's own $h = \varepsilon|x_j|$ with
      the exact-zero fallback, \texttt{.true.} floors it at unit magnitude,
      $h = \varepsilon\max(1,|x_j|)$, which is the rule for unknowns already
      scaled so that unity is their natural size (the logarithmic unknowns of
      the constrained chemical equilibrium). The \texttt{params} dummy also
      became assumed-size, so one transport array serves systems of different
      length. \texttt{hybrd1} passes \texttt{.false.} and is therefore
      unchanged in behaviour.
\item \texttt{states/\allowbreak Source.f90} gained the well-balanced arm of
      2026-09-10 (item N37 of the stage-2 solver program,
      \texttt{docs/ISSUES\_20260909.md}): under
      \texttt{Well balanced: True} the gravitational source and the geometric
      pressure term are not evaluated here at all, because they cancel the
      equilibrium part of the momentum flux difference identically and
      \texttt{RK\_rhs} assembles the row from what is left (K\"appeli \&
      Mishra 2014, J. Comput. Phys. 259, 199, eq. 2.26). The file also reads
      its pressure through \texttt{caloric\_eos} instead of a constant
      $\gamma$.
\item \texttt{radiation/\allowbreak J\_\allowbreak inc.f90} carries the single-spectrum-type decision:
      \texttt{Spectrum type:} now states the spectrum of the whole photon
      grid, the XUV and the part below 13.6\,eV alike, and no band is filled
      from a type the input did not select. The file also became the one
      definition of the eV-to-Hz conversion and of the Balmer-continuum head
      $e_{\rm th,HI}/4$.
\end{itemize}

\noindent
The two wave-speed estimators are no longer in the ``effectively untouched''
category either. \texttt{flux/\allowbreak speed\_\allowbreak estimate\_\allowbreak HLLC.f90} was deleted as dead
code, as recorded above. \texttt{flux/\allowbreak speed\_\allowbreak estimate\_\allowbreak ROE.f90} was rewritten
(99 code lines added, 43 removed) into a star-state estimate that returns a
status: the primitive-variable, two-rarefaction and two-shock branches of
Toro (2009) chapter 9 are selected by the pressure ratio, the vacuum
condition of his eq. 4.76 is tested first, and a face with no admissible star
state is reported rather than given a square root of a negative ratio.

\begin{table}[h]
\centering\small
\caption{The 8 Fortran files byte-identical between the original ATES
(\texttt{ATES-Code-main}) and EXHALE (measured 2026-09-11), grouped by role.}
\label{tab:untouched}
\begin{tabular}{@{}l p{0.60\linewidth}@{}}
\toprule
File & Purpose \\
\midrule
\multicolumn{2}{@{}l}{\emph{MINPACK linear algebra and constants} (Argonne; dir \texttt{nonlinear\_system\_solver/})}\\
\texttt{qrfac.f90}   & QR factorization via Householder transformations \\
\texttt{qform.f90}   & Accumulate the orthogonal factor $Q$ from the QR factorization \\
\texttt{r1updt.f90}  & Rank-one update of the QR factorization \\
\texttt{r1mpyq.f90}  & Apply the Givens-rotation sequence of the rank-one update \\
\texttt{dogleg.f90}  & Dogleg (trust-region) step combining Gauss--Newton and steepest descent \\
\texttt{enorm.f90}   & Euclidean norm of a vector (overflow/underflow safe) \\
\texttt{dpmpar.f90}  & Machine floating-point precision / range constants \\
\midrule
\multicolumn{2}{@{}l}{\emph{Grid setup}}\\
\texttt{init/\allowbreak set\_\allowbreak gravity\_\allowbreak grid.f90}   & Pre-evaluate gravity at cell centers and interfaces \\
\bottomrule
\end{tabular}
\end{table}
 after the converted body, so this fragment
%% carries no \appendix, no \documentclass and no \begin{document} /
%% \end{document} of its own.  Edit this file directly, then regenerate the
%% typeset log.
%%
%% Every number below is reprinted by
%%     python3 src/utils/codesize.py
%% from the repository root, together with the byte-identical list of
%% Table tab:untouched.  Rerun it and update this file in the same pass.
%%
%% Moved here on 2026-09-11 from docs/Update_EXHALE_stage1.tex, where the two
%% sections were the appendices of the stage-1 log.

% ===================================================================== %
\section{Code-size summary: EXHALE vs.\ the original ATES (2026-06-07; updated 2026-09-11)}
% ===================================================================== %
A line-by-line comparison of the Fortran source against the pristine upstream code
(\texttt{ATES/ATES-Code-main}) quantifies the extent of the in-house extension.
\emph{Pure code lines (SLOC)} exclude blank lines and comments (the comment
stripper is string-aware, so a \texttt{!} inside a quoted string is not treated as
a comment); the line-by-line diff additionally normalizes whitespace and case, so
reformatting and comment-only edits are \emph{not} counted as changes. The
measurement script is \texttt{src/utils/codesize.py}, which reprints every
number of this appendix and the byte-identical list of
Table~\ref{tab:untouched}. Measured 2026-09-11, i.e.\ after the whole of
stage 1 and the stage-2 series through item P19 of section 8: the acceptance
and certification contracts, the attempted-step controller, the partitioned
stationary solver, the oxygen and carbon chemistry, the molecular infrared
bands and the test suites those series brought with them.

\paragraph{What the script counts as \emph{core}.} Every \texttt{.f90} file
under \texttt{src/} except those in \texttt{src/modules/wind\_ae/}, which is
the ported Wind-AE IC generator and is listed separately. That rule puts the
standalone drivers under \texttt{src/tests/} in the core figure, and they are
now a large part of it, so the table splits the core row into the production
build and those drivers. Only the production files are compiled into
\texttt{EXHALE.x}; each test suite builds its own driver from
\texttt{src/tests/<suite>/run.sh}.

\begin{center}\small
\begin{tabular}{@{}lcc@{}}
\toprule
 & Fortran files & code lines (SLOC) \\
\midrule
Original ATES (\texttt{ATES-Code-main}) & 47 & 3\,514 \\
EXHALE core (excl.\ Wind-AE port)       & 169 & 71\,503 \\
\quad of which the production build     & 95 & 50\,762 ($\approx14.4\times$) \\
\quad of which the \texttt{src/tests/} drivers & 74 & 20\,741 \\
\midrule
net change (core) & $+122$ & $+67\,989$ ($\approx20.3\times$) \\
\midrule
Wind-AE IC-generator port (\texttt{wind\_ae/}) & $+36$ & $+4\,756$ \\
EXHALE total & 205 & 76\,259 ($\approx21.7\times$) \\
\bottomrule
\end{tabular}
\end{center}

\noindent
\textbf{Line-by-line diff, core} (comments + blanks excluded; whitespace/case
normalized): \textbf{69\,179} code lines added in EXHALE (new modules + the new
side of modified lines) and \textbf{1\,190} removed from the original (deleted + the
old side of modified lines), for \textbf{70\,369} total changed lines. At the file
level, of the 47 original files \textbf{38 were modified}, \textbf{8 are
unchanged} (all eight byte-for-byte; the list is Table~\ref{tab:untouched}) and
\textbf{one was removed}: \texttt{flux/\allowbreak speed\_\allowbreak estimate\_\allowbreak HLLC.f90}, Toro's
adaptive $p_\star$ estimate, was \texttt{use}d by \texttt{Num\_\allowbreak Fluxes.f90} and
never called, because the HLLC branch builds its own Davis/Einfeldt speeds
inline. \textbf{123 new core modules} were added. Of those, 49 are part of the
simulation build:

{\footnotesize
\begin{longtable}{@{}>{\raggedright\arraybackslash}p{0.36\linewidth} >{\raggedright\arraybackslash}p{0.50\linewidth} r@{}}
\toprule
New module & Purpose & SLOC \\
\midrule
\endfirsthead
\toprule
New module & Purpose & SLOC \\
\midrule
\endhead
\texttt{time\_\allowbreak step/\allowbreak steady\_\allowbreak newton.f90} & stationary solve of the hydrodynamic rows: residual $F(Y)=0$ on the physical cells, Newton/Krylov with pseudo-transient continuation, line search, best-iterate ledger & 9435 \\
\texttt{lower\_\allowbreak atmosphere/\allowbreak diffusive\_\allowbreak photochemistry.f90} & vertical transport of the molecular carriers (H$_2$, OH, H$_2$O, CO) solved together with the chemistry that makes and destroys them & 2485 \\
\texttt{lower\_\allowbreak atmosphere/\allowbreak molecular\_\allowbreak infrared\_\allowbreak data.f90} & band and line data of the molecular infrared coolant (generated by \texttt{cooling\_data/molecular\_infrared\_bands.py}) & 2182 \\
\texttt{functions/\allowbreak binary\_\allowbreak element\_\allowbreak diffusion.f90} & conservative binary H/He element transport (stage-resolved pair coefficients, donor-cell drift, metals slaved to H) & 1898 \\
\texttt{time\_\allowbreak step/\allowbreak attempted\_\allowbreak step.f90} & the attempted-step controller: one checkpoint, one trial, one adoption boundary, and the rollback contract & 1608 \\
\texttt{time\_\allowbreak step/\allowbreak certification.f90} & one evaluator of the equations a run actually solves, and the contexts that read its verdict & 1273 \\
\texttt{nonlinear\_\allowbreak system\_\allowbreak solver/\allowbreak constrained\_\allowbreak chemical\_\allowbreak equilibrium.f90} & chemical equilibrium of the molecular H/He network in logarithmic number densities, element conservation as explicit residual rows & 1173 \\
\texttt{lower\_\allowbreak atmosphere/\allowbreak h2\_\allowbreak self\_\allowbreak shielding\_\allowbreak table.f90} & H$_2$ Lyman--Werner dissociation cross section with overlapping lines, and its two branching ratios & 952 \\
\texttt{time\_\allowbreak step/\allowbreak energy\_\allowbreak semi\_\allowbreak implicit.f90} & implicit source step for the thermal energy of a cell, at fixed composition and heating & 545 \\
\texttt{time\_\allowbreak step/\allowbreak steady\_\allowbreak residual.f90} & finite-volume stationary residual $R = \mathrm{d}u/\mathrm{d}t$ and its norms, shared by the diagnostic, the convergence monitor and the stationary solver & 509 \\
\texttt{time\_\allowbreak step/\allowbreak viscous\_\allowbreak conduction.f90} & Navier--Stokes viscous momentum force, its dissipation, and thermal conduction (Crank--Nicolson) & 504 \\
\texttt{radiation/\allowbreak excited\_\allowbreak hydrogen.f90} & non-LTE H($n{=}2$) 2s/2p equilibrium (Christie, Arras \& Li 2013) and the Balmer feedback & 474 \\
\texttt{lower\_\allowbreak atmosphere/\allowbreak element\_\allowbreak inventory.f90} & one stoichiometric map of the H, He, O and C nuclei of a cell, and the three constraints a composition state is judged by & 448 \\
\texttt{radiation/\allowbreak charge\_\allowbreak exchange.f90} & charge-exchange network (Huang et al.\ 2023, Table 4) & 446 \\
\texttt{files\_\allowbreak IO/\allowbreak lower\_\allowbreak atmosphere\_\allowbreak profile.f90} & lower-atmosphere profile reader: log-$p$ interpolation, elemental reservoirs, $K_{zz}(p)$ & 367 \\
\texttt{states/\allowbreak caloric\_\allowbreak eos.f90} & caloric equation of state: internal energy density, pressure and temperature, with the H$_2$ level ladder & 340 \\
\texttt{lower\_\allowbreak atmosphere/\allowbreak molecular\_\allowbreak reaction\_\allowbreak heat.f90} & net chemical heat of the collisional reactions of the H$_2$/He network & 339 \\
\texttt{lower\_\allowbreak atmosphere/\allowbreak oxygen\_\allowbreak rates.f90} & rate coefficients and thermodynamic data of the oxygen-hydrogen chemistry, with their validity ranges & 316 \\
\texttt{functions/\allowbreak element\_\allowbreak census.f90} & one definition of the elemental and charge content of a composition state, and the invariant checks built on it & 307 \\
\texttt{nonlinear\_\allowbreak system\_\allowbreak solver/\allowbreak System\_\allowbreak HeH\_\allowbreak mol.f90} & coupled H/He $+$ molecular (H$_2$/H$_2^+$/H$_3^+$/HeH$^+$) system, turnover-scaled rows & 303 \\
\texttt{radiation/\allowbreak electron\_\allowbreak energy\_\allowbreak degradation.f90} & degradation of a fast photoelectron in a partly ionized H/H$_2$/He gas: the heat fraction and the ionizations of each target & 287 \\
\texttt{radiation/\allowbreak lya\_\allowbreak rt.f90} & Ly$\alpha$ transfer by the escape-probability closure, evaluated in line each step & 257 \\
\texttt{lower\_\allowbreak atmosphere/\allowbreak molecular\_\allowbreak infrared\_\allowbreak cooling.f90} & infrared exchange of the molecular layer: H$_2$O and CO vibration-rotation bands, H$_2$ quadrupole and magnetic-dipole lines & 231 \\
\texttt{states/\allowbreak base\_\allowbreak boundary.f90} & the lower boundary as a characteristic condition at the face, the number of conditions set by the Mach number there & 199 \\
\texttt{functions/\allowbreak h2\_\allowbreak photo\_\allowbreak channels.f90} & H$_2$ photoabsorption resolved into mutually exclusive final-state channels, with one stoichiometric table & 199 \\
\texttt{lower\_\allowbreak atmosphere/\allowbreak lyman\_\allowbreak werner.f90} & H$_2$ photodissociation in the Lyman and Werner bands (the Solomon process) & 195 \\
\texttt{radiation/\allowbreak opacity\_\allowbreak models.f90} & opacity-model dispatcher for the photoionization cross sections (A/C/P/T) & 180 \\
\texttt{files\_\allowbreak IO/\allowbreak metals\_\allowbreak input\_\allowbreak read.f90} & parse \texttt{metals.inp}: trace-metal abundances at runtime & 177 \\
\texttt{functions/\allowbreak composition.f90} & single point turning $(\rho, f_{\rm sp})$ into every species density, $n_e$ and $n_{\rm tot}$, and converting pressure and temperature & 162 \\
\texttt{flux/\allowbreak species\_\allowbreak face\_\allowbreak flux.f90} & advection of the species mass fractions on the same faces, with the same mass flux, as the density & 162 \\
\texttt{time\_\allowbreak step/\allowbreak hydrodynamic\_\allowbreak rows.f90} & the hydrodynamic rows of the stationary residual in double precision, instantiated from the kind-generic text & 161 \\
\texttt{nonlinear\_\allowbreak system\_\allowbreak solver/\allowbreak System\_\allowbreak HeH\_\allowbreak metals.f90} & coupled H/He$+$metals ionization ($+$\,analytic Jacobian) & 150 \\
\texttt{lower\_\allowbreak atmosphere/\allowbreak mol\_\allowbreak rates.f90} & molecular reaction rates (Koskinen et al.\ 2022, Table 1) & 150 \\
\texttt{lower\_\allowbreak atmosphere/\allowbreak h3p\_\allowbreak cooling.f90} & optically thin H$_3^+$ infrared cooling (Miller et al.\ 2013) & 148 \\
\texttt{nonlinear\_\allowbreak system\_\allowbreak solver/\allowbreak System\_\allowbreak HeH\_\allowbreak mol\_\allowbreak metals.f90} & molecular $+$ metals merged ionization system & 141 \\
\texttt{nonlinear\_\allowbreak system\_\allowbreak solver/\allowbreak newton\_\allowbreak solver.f90} & damped Newton with an analytic Jacobian for the ionization systems, with the MINPACK \texttt{hybrd1} retry & 128 \\
\texttt{init/\allowbreak species\_\allowbreak table.f90} & metadata table of the metal stages and of the base, molecular and oxygen species & 128 \\
\texttt{lower\_\allowbreak atmosphere/\allowbreak co\_\allowbreak photodissociation.f90} & CO photodissociation on the 912--1201\,\AA\ band & 127 \\
\texttt{nonlinear\_\allowbreak system\_\allowbreak solver/\allowbreak ion\_\allowbreak residual\_\allowbreak core.f90} & shared ionization-residual rows and analytic-Jacobian pieces of the \texttt{System\_*} family & 127 \\
\texttt{lower\_\allowbreak atmosphere/\allowbreak water\_\allowbreak photolysis.f90} & H$_2$O and OH photodissociation in the four far-ultraviolet bands of the oxygen option & 111 \\
\texttt{lower\_\allowbreak atmosphere/\allowbreak co\_\allowbreak self\_\allowbreak shielding\_\allowbreak table.f90} & CO self-shielding function $\Theta[N({\rm CO}), N({\rm H_2})]$ (Visser, van Dishoeck \& Black 2009) & 110 \\
\texttt{nonlinear\_\allowbreak system\_\allowbreak solver/\allowbreak ion\_\allowbreak cell\_\allowbreak state.f90} & named-field cell state of the ionization residuals, in place of the position-packed \texttt{params} & 93 \\
\texttt{lower\_\allowbreak atmosphere/\allowbreak lower\_\allowbreak column.f90} & analytic lower-atmosphere column from the 1-bar level to the base pressure (Koskinen et al.\ 2022) & 90 \\
\texttt{flux/\allowbreak low\_\allowbreak mach\_\allowbreak dissipation.f90} & gated fourth-difference stress damping the two-cell velocity mode a contact-resolving flux leaves undamped & 87 \\
\texttt{files\_\allowbreak IO/\allowbreak opacity\_\allowbreak input\_\allowbreak read.f90} & parse \texttt{opacity.inp} & 85 \\
\texttt{radiation/\allowbreak hydrogen\_\allowbreak n2\_\allowbreak rates.f90} & atomic data and rate coefficients of the H $n{=}2$ manifold and of the Ly$\alpha$ line, in one place & 76 \\
\texttt{nonlinear\_\allowbreak system\_\allowbreak solver/\allowbreak System\_\allowbreak HeH\_\allowbreak TR\_\allowbreak metals.f90} & merged He\,2$^3$S$+$metals system & 68 \\
\texttt{lower\_\allowbreak atmosphere/\allowbreak h2\_\allowbreak vibrational\_\allowbreak relaxation.f90} & the fate and the size of the vibrational energy left in an H$_2$ molecule, for both of its drivers & 55 \\
\texttt{nonlinear\_\allowbreak system\_\allowbreak solver/\allowbreak constrained\_\allowbreak equilibrium\_\allowbreak probe.f90} & standalone re-evaluation of one saved constrained-equilibrium state, without the hydrodynamics & 15 \\
\midrule
\multicolumn{2}{r}{production new-module total} & 30\,003 \\
\bottomrule
\end{longtable}
}

\noindent
The other 74 are the standalone test drivers under \texttt{src/tests/}, which
did not exist in ATES at all. They state what a routine is allowed to return
and are run by suite, not by the production binary:

{\footnotesize
\begin{longtable}{@{}>{\raggedright\arraybackslash}p{0.36\linewidth} >{\raggedright\arraybackslash}p{0.50\linewidth} r@{}}
\toprule
Test driver & What it asserts & SLOC \\
\midrule
\endfirsthead
\toprule
Test driver & What it asserts & SLOC \\
\midrule
\endhead
\texttt{tests/\allowbreak krylov\_\allowbreak and\_\allowbreak dogleg/\allowbreak krylov\_\allowbreak and\_\allowbreak dogleg\_\allowbreak tests.f90} & what the linear solve and the step geometry of the stationary solver are allowed to return & 2675 \\
\texttt{tests/\allowbreak steady\_\allowbreak species\_\allowbreak rows/\allowbreak steady\_\allowbreak species\_\allowbreak rows\_\allowbreak tests.f90} & the unknown space of the stationary system, one row per transported balance & 1806 \\
\texttt{tests/\allowbreak diffusion\_\allowbreak tests.f90} & the binary H/He element-diffusion operator: analytic limits and conservation & 1051 \\
\texttt{tests/\allowbreak grid\_\allowbreak and\_\allowbreak gates/\allowbreak hydrostatic\_\allowbreak residual.f90} & the momentum row returns zero on a hydrostatic state & 967 \\
\texttt{tests/\allowbreak carrier\_\allowbreak retry/\allowbreak carrier\_\allowbreak retry.f90} & the carrier-local checkpoint and the rejection-and-retry controller, on a real molecular column & 842 \\
\texttt{tests/\allowbreak certification/\allowbreak certification\_\allowbreak contexts.f90} & the certification evaluator and its contexts, against the predicates they replace & 790 \\
\texttt{tests/\allowbreak species\_\allowbreak face\_\allowbreak flux/\allowbreak species\_\allowbreak face\_\allowbreak flux\_\allowbreak tests.f90} & species advection on the face mass fluxes, and the element rows of the diffusion operator & 557 \\
\texttt{tests/\allowbreak physics\_\allowbreak probe/\allowbreak weno3\_\allowbreak reconstruction\_\allowbreak order.f90} & convergence order of the face reconstruction against the exact face value & 482 \\
\texttt{tests/\allowbreak element\_\allowbreak operator/\allowbreak element\_\allowbreak operator\_\allowbreak tests.f90} & the boundary the element transport operator is measured under & 471 \\
\texttt{tests/\allowbreak spectrum\_\allowbreak type/\allowbreak loaded\_\allowbreak sed\_\allowbreak table\_\allowbreak semantics.f90} & what a row of a loaded SED table means, and what the production consumers make of it & 417 \\
\texttt{tests/\allowbreak attempted\_\allowbreak step/\allowbreak attempted\_\allowbreak step\_\allowbreak tests.f90} & the checkpoint and the step-size policy of the attempted-step controller & 415 \\
\texttt{tests/\allowbreak a2\_\allowbreak m1/\allowbreak a2\_\allowbreak m1\_\allowbreak rate\_\allowbreak check.f90} & the oxygen rate coefficients against their published sources & 406 \\
\texttt{tests/\allowbreak spectrum\_\allowbreak type/\allowbreak balmer\_\allowbreak band\_\allowbreak quadrature.f90} & quadrature of the Balmer-continuum band integrals of \texttt{excited\_hydrogen} & 367 \\
\texttt{tests/\allowbreak energy\_\allowbreak update/\allowbreak conduction\_\allowbreak floor\_\allowbreak tests.f90} & what the Crank--Nicolson transport stage returns, refuses and conserves & 360 \\
\texttt{tests/\allowbreak physics\_\allowbreak probe/\allowbreak metal\_\allowbreak photoionization\_\allowbreak fits.f90} & every metal photoionization cross section against the published fit, and only where it is valid & 343 \\
\texttt{tests/\allowbreak element\_\allowbreak census\_\allowbreak tests.f90} & the elemental and charge invariants of a composition state & 340 \\
\texttt{tests/\allowbreak energy\_\allowbreak update/\allowbreak energy\_\allowbreak update\_\allowbreak tests.f90} & the residual-controlled temperature solve of the implicit energy source step & 334 \\
\texttt{tests/\allowbreak spectrum\_\allowbreak type/\allowbreak loaded\_\allowbreak sed\_\allowbreak threshold\_\allowbreak edges.f90} & the photon grid built from a loaded SED, and the rates formed on it & 304 \\
\texttt{tests/\allowbreak physics\_\allowbreak probe/\allowbreak species\_\allowbreak formation\_\allowbreak energy\_\allowbreak table.f90} & the formation-energy table of the energy ledger, and the reaction energies that are its differences & 275 \\
\texttt{tests/\allowbreak physics\_\allowbreak probe/\allowbreak lya\_\allowbreak band\_\allowbreak transmission.f90} & the stellar Ly$\alpha$ band reaching a cell through the atomic hydrogen column above it & 272 \\
\texttt{tests/\allowbreak spectrum\_\allowbreak type/\allowbreak wasp121b\_\allowbreak sed\_\allowbreak sensitivity.f90} & the production band quantities of the three WASP-121\,b spectrum candidates & 262 \\
\texttt{tests/\allowbreak coupled\_\allowbreak source\_\allowbreak step/\allowbreak coupled\_\allowbreak source\_\allowbreak tests.f90} & the coupled temperature-composition source step & 260 \\
\texttt{tests/\allowbreak physics\_\allowbreak probe/\allowbreak lyman\_\allowbreak werner\_\allowbreak cell\_\allowbreak mean.f90} & the Lyman--Werner rate a cell carries is the mean of the local rate over that cell & 234 \\
\texttt{tests/\allowbreak physics\_\allowbreak probe/\allowbreak lya\_\allowbreak beam\_\allowbreak cell\_\allowbreak mean.f90} & the stellar Ly$\alpha$ field of a cell is the beam mean over the cell's own optical depth & 229 \\
\texttt{tests/\allowbreak spectrum\_\allowbreak type/\allowbreak balmer\_\allowbreak continuum\_\allowbreak field.f90} & the two scalar Balmer-continuum rates of H($n{=}2$) & 227 \\
\texttt{tests/\allowbreak physics\_\allowbreak probe/\allowbreak positivity\_\allowbreak limiter\_\allowbreak scaling.f90} & the positivity scaling of the reconstructed face states (Zhang \& Shu 2010) & 223 \\
\texttt{tests/\allowbreak carrier\_\allowbreak returned\_\allowbreak state\_\allowbreak acceptance/\allowbreak carrier\_\allowbreak returned\_\allowbreak state\_\allowbreak acceptance.f90} & what makes a returned carrier state a solution, and what does not & 222 \\
\texttt{tests/\allowbreak grid\_\allowbreak and\_\allowbreak gates/\allowbreak photon\_\allowbreak grid\_\allowbreak threshold\_\allowbreak edges.f90} & every active absorber's threshold against the bin partition of the photon grid & 213 \\
\texttt{tests/\allowbreak physics\_\allowbreak probe/\allowbreak photoionization\_\allowbreak field\_\allowbreak substitution.f90} & the order in which the field and the composition that attenuates it are brought into agreement & 209 \\
\texttt{tests/\allowbreak adv\_\allowbreak static\_\allowbreak limit/\allowbreak stationarity\_\allowbreak by\_\allowbreak the\_\allowbreak mass\_\allowbreak row.f90} & which test decides that the state the post-process was handed is stationary in a cell & 199 \\
\texttt{tests/\allowbreak physics\_\allowbreak probe/\allowbreak co\_\allowbreak destruction.f90} & the CO destruction rates, the shielding table, the deposited energies and the cell mean & 195 \\
\texttt{tests/\allowbreak physics\_\allowbreak probe/\allowbreak recombination\_\allowbreak coefficient\_\allowbreak fits.f90} & every recombination coefficient against the published fit & 189 \\
\texttt{tests/\allowbreak spectrum\_\allowbreak type/\allowbreak photon\_\allowbreak grid\_\allowbreak n2\_\allowbreak floor.f90} & the photon grid built when the excited-hydrogen coupling is armed & 184 \\
\texttt{tests/\allowbreak carrier\_\allowbreak reference\_\allowbreak scales/\allowbreak carrier\_\allowbreak reference\_\allowbreak scales.f90} & the element a transported carrier is made of, and whether its chemistry rows can be differentiated & 181 \\
\texttt{tests/\allowbreak carrier\_\allowbreak constraint\_\allowbreak attribution/\allowbreak carrier\_\allowbreak constraint\_\allowbreak attribution.f90} & which residuals belong to the solver and which to a constraint & 179 \\
\texttt{tests/\allowbreak adv\_\allowbreak static\_\allowbreak limit/\allowbreak enthalpy\_\allowbreak flux\_\allowbreak of\_\allowbreak mass\_\allowbreak divergence.f90} & the third term of the advected energy equation & 178 \\
\texttt{tests/\allowbreak grid\_\allowbreak and\_\allowbreak gates/\allowbreak free\_\allowbreak outflow\_\allowbreak boundary.f90} & the outer free-outflow boundary on a supersonic wind, and its ghost state & 174 \\
\texttt{tests/\allowbreak spectrum\_\allowbreak type/\allowbreak planck\_\allowbreak photon\_\allowbreak field.f90} & the incident flux put on the photon grid for the photospheric blackbody type & 173 \\
\texttt{tests/\allowbreak a2\_\allowbreak m2/\allowbreak a2\_\allowbreak m2\_\allowbreak kinetics\_\allowbreak check.f90} & the kinetics quantities the oxygen option adds on top of the rate audit & 163 \\
\texttt{tests/\allowbreak physics\_\allowbreak probe/\allowbreak lya\_\allowbreak escape\_\allowbreak probability.f90} & the escape probability of the two-level closure against the published slab solution & 163 \\
\texttt{tests/\allowbreak e1\_\allowbreak h2/\allowbreak e1\_\allowbreak h2\_\allowbreak channel\_\allowbreak check.f90} & the mutually exclusive H$_2$ photoabsorption channels & 161 \\
\texttt{tests/\allowbreak adv\_\allowbreak static\_\allowbreak limit/\allowbreak upstream\_\allowbreak caloric\_\allowbreak energy.f90} & the upstream energy of the advected energy equation & 154 \\
\texttt{tests/\allowbreak physics\_\allowbreak probe/\allowbreak photoionization\_\allowbreak field\_\allowbreak self\_\allowbreak consistency.f90} & the field one cell is solved in, against the composition that cell returns & 148 \\
\texttt{tests/\allowbreak grid\_\allowbreak and\_\allowbreak gates/\allowbreak photon\_\allowbreak grid\_\allowbreak quadrature.f90} & the photoionization rates of one unattenuated cell against the band quadrature & 146 \\
\texttt{tests/\allowbreak physics\_\allowbreak probe/\allowbreak co\_\allowbreak helium\_\allowbreak ion\_\allowbreak sink.f90} & He$^+$ $+$ CO as a species sink, not only as a heat source & 145 \\
\texttt{tests/\allowbreak physics\_\allowbreak probe/\allowbreak riemann\_\allowbreak wave\_\allowbreak speeds.f90} & signal speeds at a face, and the state one conservative update with that flux leaves behind & 143 \\
\texttt{tests/\allowbreak physics\_\allowbreak probe/\allowbreak helium\_\allowbreak level\_\allowbreak cooling\_\allowbreak ledger.f90} & which helium level each cooling channel of \texttt{eval\_cool} is charged to & 139 \\
\texttt{tests/\allowbreak physics\_\allowbreak probe/\allowbreak heating\_\allowbreak channel\_\allowbreak closure.f90} & one assembly of the volumetric heating rate, and the closure of its channels & 136 \\
\texttt{tests/\allowbreak acceptance\_\allowbreak classes/\allowbreak acceptance\_\allowbreak class\_\allowbreak contract.f90} & one meaning of the ionization acceptance classes for every consumer & 132 \\
\texttt{tests/\allowbreak physics\_\allowbreak probe/\allowbreak opacity\_\allowbreak model\_\allowbreak p\_\allowbreak ledger.f90} & the radiative ledger of a cell whose cross sections are broadened by opacity model P & 124 \\
\texttt{tests/\allowbreak physics\_\allowbreak probe/\allowbreak photoevent\_\allowbreak energy\_\allowbreak ledger.f90} & the energy an isolated photoabsorption event has to account for & 117 \\
\texttt{tests/\allowbreak physics\_\allowbreak probe/\allowbreak xuv\_\allowbreak cell\_\allowbreak mean\_\allowbreak attenuation.f90} & the beam attenuation the photoionization rate of a cell is evaluated with & 115 \\
\texttt{tests/\allowbreak grid\_\allowbreak and\_\allowbreak gates/\allowbreak grid\_\allowbreak window\_\allowbreak indices.f90} & the two index statements that decide which cells a quantity is read from & 101 \\
\texttt{tests/\allowbreak physics\_\allowbreak probe/\allowbreak h2\_\allowbreak rovibrational\_\allowbreak identity.f90} & the identities tying the H$_2$ partition function of the chemistry to the internal energy of the caloric EOS & 96 \\
\texttt{tests/\allowbreak adv\_\allowbreak static\_\allowbreak limit/\allowbreak stationary\_\allowbreak flux\_\allowbreak refusal.f90} & the sensitivity of the advected energy correction to a non-stationary mass flux & 91 \\
\texttt{tests/\allowbreak physics\_\allowbreak probe/\allowbreak ionization\_\allowbreak threshold\_\allowbreak turn\_\allowbreak on.f90} & one definition of each ionization threshold, and every cross section turning on exactly there & 90 \\
\texttt{tests/\allowbreak grid\_\allowbreak and\_\allowbreak gates/\allowbreak grid\_\allowbreak width\_\allowbreak identity.f90} & the cell-width and total-width identities of the radial grid & 83 \\
\texttt{tests/\allowbreak physics\_\allowbreak probe/\allowbreak h2\_\allowbreak channel\_\allowbreak detector\_\allowbreak ratio.f90} & the measured proton-to-H$_2^+$ ratio of Chung et al.\ (1993) against the four channels & 83 \\
\texttt{tests/\allowbreak physics\_\allowbreak probe/\allowbreak fine\_\allowbreak structure\_\allowbreak escape\_\allowbreak probability.f90} & the escape probability of the ground-term fine-structure lines, and the depth it is fed & 80 \\
\texttt{tests/\allowbreak physics\_\allowbreak probe/\allowbreak atomic\_\allowbreak mass\_\allowbreak and\_\allowbreak radius\_\allowbreak constants.f90} & the helium mass and the astronomical constants of the input readers, one definition each & 79 \\
\texttt{tests/\allowbreak physics\_\allowbreak probe/\allowbreak h3p\_\allowbreak cooling\_\allowbreak limits.f90} & the limits, the continuities and the published values of the H$_3^+$ cooling rate & 79 \\
\texttt{tests/\allowbreak physics\_\allowbreak probe/\allowbreak assertion\_\allowbreak report.f90} & one printed verdict line per assertion, and the failure count the driver exits on & 74 \\
\texttt{tests/\allowbreak adv\_\allowbreak static\_\allowbreak limit/\allowbreak thermal\_\allowbreak damkohler\_\allowbreak limit.f90} & the validity condition of the advection-corrected energy solve, as a function of the flow speed & 72 \\
\texttt{tests/\allowbreak a2\_\allowbreak m2/\allowbreak a2\_\allowbreak m2\_\allowbreak g3\_\allowbreak run\_\allowbreak check.f90} & with every photolysis rate at zero, the solved O/OH/H$_2$O partition is the chemical equilibrium & 71 \\
\texttt{tests/\allowbreak physics\_\allowbreak probe/\allowbreak absorbed\_\allowbreak fraction\_\allowbreak switch.f90} & $(1-e^{-d})/d$, the fraction of the photons entering a cell that the cell absorbs & 71 \\
\texttt{tests/\allowbreak constrained\_\allowbreak network\_\allowbreak layout/\allowbreak transported\_\allowbreak proton\_\allowbreak constraint\_\allowbreak layout.f90} & which species are unknowns of the continuation, and whether the system built is square & 71 \\
\texttt{tests/\allowbreak test\_\allowbreak columns.f90} & synthetic columns whose species carry the mass density they are measured against & 67 \\
\texttt{tests/\allowbreak physics\_\allowbreak probe/\allowbreak h2\_\allowbreak infrared\_\allowbreak line\_\allowbreak populations.f90} & the LTE populations weighting the H$_2$ quadrupole lines are those of the caloric EOS ladder & 65 \\
\texttt{tests/\allowbreak physics\_\allowbreak probe/\allowbreak wind\_\allowbreak ae\_\allowbreak bridge\_\allowbreak composition.f90} & the composition an EXHALE \texttt{input.inp} hands the Wind-AE solver & 54 \\
\texttt{tests/\allowbreak physics\_\allowbreak probe/\allowbreak voronov\_\allowbreak carbon\_\allowbreak ionization.f90} & the electron-impact ionization of C$^0$ and C$^+$ against the published Voronov (1997) rows & 51 \\
\texttt{tests/\allowbreak physics\_\allowbreak probe/\allowbreak water\_\allowbreak photolysis\_\allowbreak lyman\_\allowbreak alpha\_\allowbreak yields.f90} & the H$_2$O branching on the band that contains Ly$\alpha$, against the published yields & 49 \\
\texttt{tests/\allowbreak physics\_\allowbreak probe/\allowbreak h3p\_\allowbreak recombination\_\allowbreak branching.f90} & the two product channels of H$_3^+$ dissociative recombination & 44 \\
\texttt{tests/\allowbreak energy\_\allowbreak update/\allowbreak cooling\_\allowbreak models.f90} & analytic cooling laws for the tests of the implicit energy update & 43 \\
\texttt{tests/\allowbreak species\_\allowbreak masses/\allowbreak carrier\_\allowbreak background\_\allowbreak masses.f90} & the mass the carrier transport gives each of its background collision partners & 36 \\
\midrule
\multicolumn{2}{r}{test-driver total} & 20\,741 \\
\midrule
\multicolumn{2}{r}{new-module total} & 50\,744 \\
\bottomrule
\end{longtable}
}

\noindent
So the 123 new core modules contribute 50\,744 of the 69\,179 added code lines;
the remaining 18\,435 added (and 1\,121 of the 1\,190 removed, the other 69
being the deleted HLLC estimator) are edits spread across the 38 modified
files. The ten heaviest, by added code lines, are the main driver
(\texttt{EXHALE\_\allowbreak main.f90}, 3\,498 added, 246 $\to$ 3\,625 SLOC: the marching
loop, the attempted-step adoption boundary, the run-mode contracts and the
written-state assertions), the ionization equilibrium and its diagnostics
(\texttt{ionization\_\allowbreak equilibrium.f90} 2\,274 and \texttt{util\_\allowbreak ion\_\allowbreak eq.f90}
1\,803: the metal, molecular and oxygen networks, the acceptance classes, the
photoionization field and the heating channels), the cooling module
(\texttt{Cool\_\allowbreak coeff.f90}, 2\,250 added, 442 $\to$ 3\,966 raw lines: CHIANTI
metal-line fits, Fe tables, fine-structure statistical equilibrium, Penning
and molecular rate coefficients), input parsing
(\texttt{input\_\allowbreak read.f90}, 1\,753), the state reader and writer
(\texttt{load\_\allowbreak IC.f90} 1\,063, \texttt{write\_\allowbreak output.f90} 928,
\texttt{write\_\allowbreak setup\_\allowbreak report.f90} 751), the advection-corrected post-processor
(\texttt{post\_\allowbreak process\_\allowbreak adv.f90}, 647) and the cross sections
(\texttt{cross\_\allowbreak sec.f90}, 438). Those ten carry 15\,405 of the 18\,435. The rest
of the edits are in the reconstruction, the fluxes, the boundary conditions,
the grid, the time step, the energy equation and the \texttt{System\_*}
ionization systems, which wire the metals, the molecules, the opacity
dispatcher, the Roche potential, the transonic/auto/Wind-AE IC, the
excited-H and Ly$\alpha$ physics, the element diffusion and the Brent/Newton
solver upgrades into the original solver.

In total the EXHALE core is now $\approx$20.3$\times$ the original code base
(3\,514\,$\to$\,71\,503 SLOC), or $\approx$14.4$\times$ counting only the
production build (50\,762 SLOC), touching almost all of the original files (38
of 47 modified, one deleted, plus 123 new); with the ported Wind-AE IC
generator the full Fortran tree is 205 files / 76\,259 SLOC
($\approx$21.7$\times$). Beyond Fortran, the tree also carries the Python
tooling (GUI, the \texttt{EXHALE\_transit.py} transmission post-processor,
plotting/analysis, the lower-atmosphere adapters and the elemental-flux
closure driver) and 181\,816 lines of documentation (187 md/tex files: the
memos, the manual and this log) under \texttt{docs/}.


% ===================================================================== %
\section{Fortran source inherited unchanged from ATES}\label{sec:untouched}
% ===================================================================== %
Of the 47 upstream Fortran files (\texttt{ATES/ATES-Code-main}), \textbf{8 are
byte-for-byte identical} in EXHALE (re-measured 2026-09-11): the numerics
EXHALE reuses without touching, which are now the linear-algebra part of the
included \textsc{MINPACK} nonlinear solver plus the gravity pre-evaluation.
Table~\ref{tab:untouched} lists them with their role. No file is identical
after whitespace and case normalization but not byte-identical, so the
normalized-unchanged set and the byte-identical set are the same eight.

Five files left the byte-identical set since the 2026-08-30 edition of this
table, each for a documented reason rather than drift.

\begin{itemize}\itemsep2pt
\item \texttt{nonlinear\_\allowbreak system\_\allowbreak solver/\allowbreak hybrd.f90},
      \texttt{hybrd1.f90} and \texttt{fdjac1.f90} carry a new
      \texttt{unit\_step\_floor} argument threaded from \texttt{hybrd} to
      \texttt{fdjac1}, which selects the forward-difference step of the
      Jacobian: \texttt{.false.} is MINPACK's own $h = \varepsilon|x_j|$ with
      the exact-zero fallback, \texttt{.true.} floors it at unit magnitude,
      $h = \varepsilon\max(1,|x_j|)$, which is the rule for unknowns already
      scaled so that unity is their natural size (the logarithmic unknowns of
      the constrained chemical equilibrium). The \texttt{params} dummy also
      became assumed-size, so one transport array serves systems of different
      length. \texttt{hybrd1} passes \texttt{.false.} and is therefore
      unchanged in behaviour.
\item \texttt{states/\allowbreak Source.f90} gained the well-balanced arm of
      2026-09-10 (item N37 of the stage-2 solver program,
      \texttt{docs/ISSUES\_20260909.md}): under
      \texttt{Well balanced: True} the gravitational source and the geometric
      pressure term are not evaluated here at all, because they cancel the
      equilibrium part of the momentum flux difference identically and
      \texttt{RK\_rhs} assembles the row from what is left (K\"appeli \&
      Mishra 2014, J. Comput. Phys. 259, 199, eq. 2.26). The file also reads
      its pressure through \texttt{caloric\_eos} instead of a constant
      $\gamma$.
\item \texttt{radiation/\allowbreak J\_\allowbreak inc.f90} carries the single-spectrum-type decision:
      \texttt{Spectrum type:} now states the spectrum of the whole photon
      grid, the XUV and the part below 13.6\,eV alike, and no band is filled
      from a type the input did not select. The file also became the one
      definition of the eV-to-Hz conversion and of the Balmer-continuum head
      $e_{\rm th,HI}/4$.
\end{itemize}

\noindent
The two wave-speed estimators are no longer in the ``effectively untouched''
category either. \texttt{flux/\allowbreak speed\_\allowbreak estimate\_\allowbreak HLLC.f90} was deleted as dead
code, as recorded above. \texttt{flux/\allowbreak speed\_\allowbreak estimate\_\allowbreak ROE.f90} was rewritten
(99 code lines added, 43 removed) into a star-state estimate that returns a
status: the primitive-variable, two-rarefaction and two-shock branches of
Toro (2009) chapter 9 are selected by the pressure ratio, the vacuum
condition of his eq. 4.76 is tested first, and a face with no admissible star
state is reported rather than given a square root of a negative ratio.

\begin{table}[h]
\centering\small
\caption{The 8 Fortran files byte-identical between the original ATES
(\texttt{ATES-Code-main}) and EXHALE (measured 2026-09-11), grouped by role.}
\label{tab:untouched}
\begin{tabular}{@{}l p{0.60\linewidth}@{}}
\toprule
File & Purpose \\
\midrule
\multicolumn{2}{@{}l}{\emph{MINPACK linear algebra and constants} (Argonne; dir \texttt{nonlinear\_system\_solver/})}\\
\texttt{qrfac.f90}   & QR factorization via Householder transformations \\
\texttt{qform.f90}   & Accumulate the orthogonal factor $Q$ from the QR factorization \\
\texttt{r1updt.f90}  & Rank-one update of the QR factorization \\
\texttt{r1mpyq.f90}  & Apply the Givens-rotation sequence of the rank-one update \\
\texttt{dogleg.f90}  & Dogleg (trust-region) step combining Gauss--Newton and steepest descent \\
\texttt{enorm.f90}   & Euclidean norm of a vector (overflow/underflow safe) \\
\texttt{dpmpar.f90}  & Machine floating-point precision / range constants \\
\midrule
\multicolumn{2}{@{}l}{\emph{Grid setup}}\\
\texttt{init/\allowbreak set\_\allowbreak gravity\_\allowbreak grid.f90}   & Pre-evaluate gravity at cell centers and interfaces \\
\bottomrule
\end{tabular}
\end{table}
 after the converted body, so this fragment
%% carries no \appendix, no \documentclass and no \begin{document} /
%% \end{document} of its own.  Edit this file directly, then regenerate the
%% typeset log.
%%
%% Every number below is reprinted by
%%     python3 src/utils/codesize.py
%% from the repository root, together with the byte-identical list of
%% Table tab:untouched.  Rerun it and update this file in the same pass.
%%
%% Moved here on 2026-09-11 from docs/Update_EXHALE_stage1.tex, where the two
%% sections were the appendices of the stage-1 log.

% ===================================================================== %
\section{Code-size summary: EXHALE vs.\ the original ATES (2026-06-07; updated 2026-09-11)}
% ===================================================================== %
A line-by-line comparison of the Fortran source against the pristine upstream code
(\texttt{ATES/ATES-Code-main}) quantifies the extent of the in-house extension.
\emph{Pure code lines (SLOC)} exclude blank lines and comments (the comment
stripper is string-aware, so a \texttt{!} inside a quoted string is not treated as
a comment); the line-by-line diff additionally normalizes whitespace and case, so
reformatting and comment-only edits are \emph{not} counted as changes. The
measurement script is \texttt{src/utils/codesize.py}, which reprints every
number of this appendix and the byte-identical list of
Table~\ref{tab:untouched}. Measured 2026-09-11, i.e.\ after the whole of
stage 1 and the stage-2 series through item P19 of section 8: the acceptance
and certification contracts, the attempted-step controller, the partitioned
stationary solver, the oxygen and carbon chemistry, the molecular infrared
bands and the test suites those series brought with them.

\paragraph{What the script counts as \emph{core}.} Every \texttt{.f90} file
under \texttt{src/} except those in \texttt{src/modules/wind\_ae/}, which is
the ported Wind-AE IC generator and is listed separately. That rule puts the
standalone drivers under \texttt{src/tests/} in the core figure, and they are
now a large part of it, so the table splits the core row into the production
build and those drivers. Only the production files are compiled into
\texttt{EXHALE.x}; each test suite builds its own driver from
\texttt{src/tests/<suite>/run.sh}.

\begin{center}\small
\begin{tabular}{@{}lcc@{}}
\toprule
 & Fortran files & code lines (SLOC) \\
\midrule
Original ATES (\texttt{ATES-Code-main}) & 47 & 3\,514 \\
EXHALE core (excl.\ Wind-AE port)       & 169 & 71\,503 \\
\quad of which the production build     & 95 & 50\,762 ($\approx14.4\times$) \\
\quad of which the \texttt{src/tests/} drivers & 74 & 20\,741 \\
\midrule
net change (core) & $+122$ & $+67\,989$ ($\approx20.3\times$) \\
\midrule
Wind-AE IC-generator port (\texttt{wind\_ae/}) & $+36$ & $+4\,756$ \\
EXHALE total & 205 & 76\,259 ($\approx21.7\times$) \\
\bottomrule
\end{tabular}
\end{center}

\noindent
\textbf{Line-by-line diff, core} (comments + blanks excluded; whitespace/case
normalized): \textbf{69\,179} code lines added in EXHALE (new modules + the new
side of modified lines) and \textbf{1\,190} removed from the original (deleted + the
old side of modified lines), for \textbf{70\,369} total changed lines. At the file
level, of the 47 original files \textbf{38 were modified}, \textbf{8 are
unchanged} (all eight byte-for-byte; the list is Table~\ref{tab:untouched}) and
\textbf{one was removed}: \texttt{flux/\allowbreak speed\_\allowbreak estimate\_\allowbreak HLLC.f90}, Toro's
adaptive $p_\star$ estimate, was \texttt{use}d by \texttt{Num\_\allowbreak Fluxes.f90} and
never called, because the HLLC branch builds its own Davis/Einfeldt speeds
inline. \textbf{123 new core modules} were added. Of those, 49 are part of the
simulation build:

{\footnotesize
\begin{longtable}{@{}>{\raggedright\arraybackslash}p{0.36\linewidth} >{\raggedright\arraybackslash}p{0.50\linewidth} r@{}}
\toprule
New module & Purpose & SLOC \\
\midrule
\endfirsthead
\toprule
New module & Purpose & SLOC \\
\midrule
\endhead
\texttt{time\_\allowbreak step/\allowbreak steady\_\allowbreak newton.f90} & stationary solve of the hydrodynamic rows: residual $F(Y)=0$ on the physical cells, Newton/Krylov with pseudo-transient continuation, line search, best-iterate ledger & 9435 \\
\texttt{lower\_\allowbreak atmosphere/\allowbreak diffusive\_\allowbreak photochemistry.f90} & vertical transport of the molecular carriers (H$_2$, OH, H$_2$O, CO) solved together with the chemistry that makes and destroys them & 2485 \\
\texttt{lower\_\allowbreak atmosphere/\allowbreak molecular\_\allowbreak infrared\_\allowbreak data.f90} & band and line data of the molecular infrared coolant (generated by \texttt{cooling\_data/molecular\_infrared\_bands.py}) & 2182 \\
\texttt{functions/\allowbreak binary\_\allowbreak element\_\allowbreak diffusion.f90} & conservative binary H/He element transport (stage-resolved pair coefficients, donor-cell drift, metals slaved to H) & 1898 \\
\texttt{time\_\allowbreak step/\allowbreak attempted\_\allowbreak step.f90} & the attempted-step controller: one checkpoint, one trial, one adoption boundary, and the rollback contract & 1608 \\
\texttt{time\_\allowbreak step/\allowbreak certification.f90} & one evaluator of the equations a run actually solves, and the contexts that read its verdict & 1273 \\
\texttt{nonlinear\_\allowbreak system\_\allowbreak solver/\allowbreak constrained\_\allowbreak chemical\_\allowbreak equilibrium.f90} & chemical equilibrium of the molecular H/He network in logarithmic number densities, element conservation as explicit residual rows & 1173 \\
\texttt{lower\_\allowbreak atmosphere/\allowbreak h2\_\allowbreak self\_\allowbreak shielding\_\allowbreak table.f90} & H$_2$ Lyman--Werner dissociation cross section with overlapping lines, and its two branching ratios & 952 \\
\texttt{time\_\allowbreak step/\allowbreak energy\_\allowbreak semi\_\allowbreak implicit.f90} & implicit source step for the thermal energy of a cell, at fixed composition and heating & 545 \\
\texttt{time\_\allowbreak step/\allowbreak steady\_\allowbreak residual.f90} & finite-volume stationary residual $R = \mathrm{d}u/\mathrm{d}t$ and its norms, shared by the diagnostic, the convergence monitor and the stationary solver & 509 \\
\texttt{time\_\allowbreak step/\allowbreak viscous\_\allowbreak conduction.f90} & Navier--Stokes viscous momentum force, its dissipation, and thermal conduction (Crank--Nicolson) & 504 \\
\texttt{radiation/\allowbreak excited\_\allowbreak hydrogen.f90} & non-LTE H($n{=}2$) 2s/2p equilibrium (Christie, Arras \& Li 2013) and the Balmer feedback & 474 \\
\texttt{lower\_\allowbreak atmosphere/\allowbreak element\_\allowbreak inventory.f90} & one stoichiometric map of the H, He, O and C nuclei of a cell, and the three constraints a composition state is judged by & 448 \\
\texttt{radiation/\allowbreak charge\_\allowbreak exchange.f90} & charge-exchange network (Huang et al.\ 2023, Table 4) & 446 \\
\texttt{files\_\allowbreak IO/\allowbreak lower\_\allowbreak atmosphere\_\allowbreak profile.f90} & lower-atmosphere profile reader: log-$p$ interpolation, elemental reservoirs, $K_{zz}(p)$ & 367 \\
\texttt{states/\allowbreak caloric\_\allowbreak eos.f90} & caloric equation of state: internal energy density, pressure and temperature, with the H$_2$ level ladder & 340 \\
\texttt{lower\_\allowbreak atmosphere/\allowbreak molecular\_\allowbreak reaction\_\allowbreak heat.f90} & net chemical heat of the collisional reactions of the H$_2$/He network & 339 \\
\texttt{lower\_\allowbreak atmosphere/\allowbreak oxygen\_\allowbreak rates.f90} & rate coefficients and thermodynamic data of the oxygen-hydrogen chemistry, with their validity ranges & 316 \\
\texttt{functions/\allowbreak element\_\allowbreak census.f90} & one definition of the elemental and charge content of a composition state, and the invariant checks built on it & 307 \\
\texttt{nonlinear\_\allowbreak system\_\allowbreak solver/\allowbreak System\_\allowbreak HeH\_\allowbreak mol.f90} & coupled H/He $+$ molecular (H$_2$/H$_2^+$/H$_3^+$/HeH$^+$) system, turnover-scaled rows & 303 \\
\texttt{radiation/\allowbreak electron\_\allowbreak energy\_\allowbreak degradation.f90} & degradation of a fast photoelectron in a partly ionized H/H$_2$/He gas: the heat fraction and the ionizations of each target & 287 \\
\texttt{radiation/\allowbreak lya\_\allowbreak rt.f90} & Ly$\alpha$ transfer by the escape-probability closure, evaluated in line each step & 257 \\
\texttt{lower\_\allowbreak atmosphere/\allowbreak molecular\_\allowbreak infrared\_\allowbreak cooling.f90} & infrared exchange of the molecular layer: H$_2$O and CO vibration-rotation bands, H$_2$ quadrupole and magnetic-dipole lines & 231 \\
\texttt{states/\allowbreak base\_\allowbreak boundary.f90} & the lower boundary as a characteristic condition at the face, the number of conditions set by the Mach number there & 199 \\
\texttt{functions/\allowbreak h2\_\allowbreak photo\_\allowbreak channels.f90} & H$_2$ photoabsorption resolved into mutually exclusive final-state channels, with one stoichiometric table & 199 \\
\texttt{lower\_\allowbreak atmosphere/\allowbreak lyman\_\allowbreak werner.f90} & H$_2$ photodissociation in the Lyman and Werner bands (the Solomon process) & 195 \\
\texttt{radiation/\allowbreak opacity\_\allowbreak models.f90} & opacity-model dispatcher for the photoionization cross sections (A/C/P/T) & 180 \\
\texttt{files\_\allowbreak IO/\allowbreak metals\_\allowbreak input\_\allowbreak read.f90} & parse \texttt{metals.inp}: trace-metal abundances at runtime & 177 \\
\texttt{functions/\allowbreak composition.f90} & single point turning $(\rho, f_{\rm sp})$ into every species density, $n_e$ and $n_{\rm tot}$, and converting pressure and temperature & 162 \\
\texttt{flux/\allowbreak species\_\allowbreak face\_\allowbreak flux.f90} & advection of the species mass fractions on the same faces, with the same mass flux, as the density & 162 \\
\texttt{time\_\allowbreak step/\allowbreak hydrodynamic\_\allowbreak rows.f90} & the hydrodynamic rows of the stationary residual in double precision, instantiated from the kind-generic text & 161 \\
\texttt{nonlinear\_\allowbreak system\_\allowbreak solver/\allowbreak System\_\allowbreak HeH\_\allowbreak metals.f90} & coupled H/He$+$metals ionization ($+$\,analytic Jacobian) & 150 \\
\texttt{lower\_\allowbreak atmosphere/\allowbreak mol\_\allowbreak rates.f90} & molecular reaction rates (Koskinen et al.\ 2022, Table 1) & 150 \\
\texttt{lower\_\allowbreak atmosphere/\allowbreak h3p\_\allowbreak cooling.f90} & optically thin H$_3^+$ infrared cooling (Miller et al.\ 2013) & 148 \\
\texttt{nonlinear\_\allowbreak system\_\allowbreak solver/\allowbreak System\_\allowbreak HeH\_\allowbreak mol\_\allowbreak metals.f90} & molecular $+$ metals merged ionization system & 141 \\
\texttt{nonlinear\_\allowbreak system\_\allowbreak solver/\allowbreak newton\_\allowbreak solver.f90} & damped Newton with an analytic Jacobian for the ionization systems, with the MINPACK \texttt{hybrd1} retry & 128 \\
\texttt{init/\allowbreak species\_\allowbreak table.f90} & metadata table of the metal stages and of the base, molecular and oxygen species & 128 \\
\texttt{lower\_\allowbreak atmosphere/\allowbreak co\_\allowbreak photodissociation.f90} & CO photodissociation on the 912--1201\,\AA\ band & 127 \\
\texttt{nonlinear\_\allowbreak system\_\allowbreak solver/\allowbreak ion\_\allowbreak residual\_\allowbreak core.f90} & shared ionization-residual rows and analytic-Jacobian pieces of the \texttt{System\_*} family & 127 \\
\texttt{lower\_\allowbreak atmosphere/\allowbreak water\_\allowbreak photolysis.f90} & H$_2$O and OH photodissociation in the four far-ultraviolet bands of the oxygen option & 111 \\
\texttt{lower\_\allowbreak atmosphere/\allowbreak co\_\allowbreak self\_\allowbreak shielding\_\allowbreak table.f90} & CO self-shielding function $\Theta[N({\rm CO}), N({\rm H_2})]$ (Visser, van Dishoeck \& Black 2009) & 110 \\
\texttt{nonlinear\_\allowbreak system\_\allowbreak solver/\allowbreak ion\_\allowbreak cell\_\allowbreak state.f90} & named-field cell state of the ionization residuals, in place of the position-packed \texttt{params} & 93 \\
\texttt{lower\_\allowbreak atmosphere/\allowbreak lower\_\allowbreak column.f90} & analytic lower-atmosphere column from the 1-bar level to the base pressure (Koskinen et al.\ 2022) & 90 \\
\texttt{flux/\allowbreak low\_\allowbreak mach\_\allowbreak dissipation.f90} & gated fourth-difference stress damping the two-cell velocity mode a contact-resolving flux leaves undamped & 87 \\
\texttt{files\_\allowbreak IO/\allowbreak opacity\_\allowbreak input\_\allowbreak read.f90} & parse \texttt{opacity.inp} & 85 \\
\texttt{radiation/\allowbreak hydrogen\_\allowbreak n2\_\allowbreak rates.f90} & atomic data and rate coefficients of the H $n{=}2$ manifold and of the Ly$\alpha$ line, in one place & 76 \\
\texttt{nonlinear\_\allowbreak system\_\allowbreak solver/\allowbreak System\_\allowbreak HeH\_\allowbreak TR\_\allowbreak metals.f90} & merged He\,2$^3$S$+$metals system & 68 \\
\texttt{lower\_\allowbreak atmosphere/\allowbreak h2\_\allowbreak vibrational\_\allowbreak relaxation.f90} & the fate and the size of the vibrational energy left in an H$_2$ molecule, for both of its drivers & 55 \\
\texttt{nonlinear\_\allowbreak system\_\allowbreak solver/\allowbreak constrained\_\allowbreak equilibrium\_\allowbreak probe.f90} & standalone re-evaluation of one saved constrained-equilibrium state, without the hydrodynamics & 15 \\
\midrule
\multicolumn{2}{r}{production new-module total} & 30\,003 \\
\bottomrule
\end{longtable}
}

\noindent
The other 74 are the standalone test drivers under \texttt{src/tests/}, which
did not exist in ATES at all. They state what a routine is allowed to return
and are run by suite, not by the production binary:

{\footnotesize
\begin{longtable}{@{}>{\raggedright\arraybackslash}p{0.36\linewidth} >{\raggedright\arraybackslash}p{0.50\linewidth} r@{}}
\toprule
Test driver & What it asserts & SLOC \\
\midrule
\endfirsthead
\toprule
Test driver & What it asserts & SLOC \\
\midrule
\endhead
\texttt{tests/\allowbreak krylov\_\allowbreak and\_\allowbreak dogleg/\allowbreak krylov\_\allowbreak and\_\allowbreak dogleg\_\allowbreak tests.f90} & what the linear solve and the step geometry of the stationary solver are allowed to return & 2675 \\
\texttt{tests/\allowbreak steady\_\allowbreak species\_\allowbreak rows/\allowbreak steady\_\allowbreak species\_\allowbreak rows\_\allowbreak tests.f90} & the unknown space of the stationary system, one row per transported balance & 1806 \\
\texttt{tests/\allowbreak diffusion\_\allowbreak tests.f90} & the binary H/He element-diffusion operator: analytic limits and conservation & 1051 \\
\texttt{tests/\allowbreak grid\_\allowbreak and\_\allowbreak gates/\allowbreak hydrostatic\_\allowbreak residual.f90} & the momentum row returns zero on a hydrostatic state & 967 \\
\texttt{tests/\allowbreak carrier\_\allowbreak retry/\allowbreak carrier\_\allowbreak retry.f90} & the carrier-local checkpoint and the rejection-and-retry controller, on a real molecular column & 842 \\
\texttt{tests/\allowbreak certification/\allowbreak certification\_\allowbreak contexts.f90} & the certification evaluator and its contexts, against the predicates they replace & 790 \\
\texttt{tests/\allowbreak species\_\allowbreak face\_\allowbreak flux/\allowbreak species\_\allowbreak face\_\allowbreak flux\_\allowbreak tests.f90} & species advection on the face mass fluxes, and the element rows of the diffusion operator & 557 \\
\texttt{tests/\allowbreak physics\_\allowbreak probe/\allowbreak weno3\_\allowbreak reconstruction\_\allowbreak order.f90} & convergence order of the face reconstruction against the exact face value & 482 \\
\texttt{tests/\allowbreak element\_\allowbreak operator/\allowbreak element\_\allowbreak operator\_\allowbreak tests.f90} & the boundary the element transport operator is measured under & 471 \\
\texttt{tests/\allowbreak spectrum\_\allowbreak type/\allowbreak loaded\_\allowbreak sed\_\allowbreak table\_\allowbreak semantics.f90} & what a row of a loaded SED table means, and what the production consumers make of it & 417 \\
\texttt{tests/\allowbreak attempted\_\allowbreak step/\allowbreak attempted\_\allowbreak step\_\allowbreak tests.f90} & the checkpoint and the step-size policy of the attempted-step controller & 415 \\
\texttt{tests/\allowbreak a2\_\allowbreak m1/\allowbreak a2\_\allowbreak m1\_\allowbreak rate\_\allowbreak check.f90} & the oxygen rate coefficients against their published sources & 406 \\
\texttt{tests/\allowbreak spectrum\_\allowbreak type/\allowbreak balmer\_\allowbreak band\_\allowbreak quadrature.f90} & quadrature of the Balmer-continuum band integrals of \texttt{excited\_hydrogen} & 367 \\
\texttt{tests/\allowbreak energy\_\allowbreak update/\allowbreak conduction\_\allowbreak floor\_\allowbreak tests.f90} & what the Crank--Nicolson transport stage returns, refuses and conserves & 360 \\
\texttt{tests/\allowbreak physics\_\allowbreak probe/\allowbreak metal\_\allowbreak photoionization\_\allowbreak fits.f90} & every metal photoionization cross section against the published fit, and only where it is valid & 343 \\
\texttt{tests/\allowbreak element\_\allowbreak census\_\allowbreak tests.f90} & the elemental and charge invariants of a composition state & 340 \\
\texttt{tests/\allowbreak energy\_\allowbreak update/\allowbreak energy\_\allowbreak update\_\allowbreak tests.f90} & the residual-controlled temperature solve of the implicit energy source step & 334 \\
\texttt{tests/\allowbreak spectrum\_\allowbreak type/\allowbreak loaded\_\allowbreak sed\_\allowbreak threshold\_\allowbreak edges.f90} & the photon grid built from a loaded SED, and the rates formed on it & 304 \\
\texttt{tests/\allowbreak physics\_\allowbreak probe/\allowbreak species\_\allowbreak formation\_\allowbreak energy\_\allowbreak table.f90} & the formation-energy table of the energy ledger, and the reaction energies that are its differences & 275 \\
\texttt{tests/\allowbreak physics\_\allowbreak probe/\allowbreak lya\_\allowbreak band\_\allowbreak transmission.f90} & the stellar Ly$\alpha$ band reaching a cell through the atomic hydrogen column above it & 272 \\
\texttt{tests/\allowbreak spectrum\_\allowbreak type/\allowbreak wasp121b\_\allowbreak sed\_\allowbreak sensitivity.f90} & the production band quantities of the three WASP-121\,b spectrum candidates & 262 \\
\texttt{tests/\allowbreak coupled\_\allowbreak source\_\allowbreak step/\allowbreak coupled\_\allowbreak source\_\allowbreak tests.f90} & the coupled temperature-composition source step & 260 \\
\texttt{tests/\allowbreak physics\_\allowbreak probe/\allowbreak lyman\_\allowbreak werner\_\allowbreak cell\_\allowbreak mean.f90} & the Lyman--Werner rate a cell carries is the mean of the local rate over that cell & 234 \\
\texttt{tests/\allowbreak physics\_\allowbreak probe/\allowbreak lya\_\allowbreak beam\_\allowbreak cell\_\allowbreak mean.f90} & the stellar Ly$\alpha$ field of a cell is the beam mean over the cell's own optical depth & 229 \\
\texttt{tests/\allowbreak spectrum\_\allowbreak type/\allowbreak balmer\_\allowbreak continuum\_\allowbreak field.f90} & the two scalar Balmer-continuum rates of H($n{=}2$) & 227 \\
\texttt{tests/\allowbreak physics\_\allowbreak probe/\allowbreak positivity\_\allowbreak limiter\_\allowbreak scaling.f90} & the positivity scaling of the reconstructed face states (Zhang \& Shu 2010) & 223 \\
\texttt{tests/\allowbreak carrier\_\allowbreak returned\_\allowbreak state\_\allowbreak acceptance/\allowbreak carrier\_\allowbreak returned\_\allowbreak state\_\allowbreak acceptance.f90} & what makes a returned carrier state a solution, and what does not & 222 \\
\texttt{tests/\allowbreak grid\_\allowbreak and\_\allowbreak gates/\allowbreak photon\_\allowbreak grid\_\allowbreak threshold\_\allowbreak edges.f90} & every active absorber's threshold against the bin partition of the photon grid & 213 \\
\texttt{tests/\allowbreak physics\_\allowbreak probe/\allowbreak photoionization\_\allowbreak field\_\allowbreak substitution.f90} & the order in which the field and the composition that attenuates it are brought into agreement & 209 \\
\texttt{tests/\allowbreak adv\_\allowbreak static\_\allowbreak limit/\allowbreak stationarity\_\allowbreak by\_\allowbreak the\_\allowbreak mass\_\allowbreak row.f90} & which test decides that the state the post-process was handed is stationary in a cell & 199 \\
\texttt{tests/\allowbreak physics\_\allowbreak probe/\allowbreak co\_\allowbreak destruction.f90} & the CO destruction rates, the shielding table, the deposited energies and the cell mean & 195 \\
\texttt{tests/\allowbreak physics\_\allowbreak probe/\allowbreak recombination\_\allowbreak coefficient\_\allowbreak fits.f90} & every recombination coefficient against the published fit & 189 \\
\texttt{tests/\allowbreak spectrum\_\allowbreak type/\allowbreak photon\_\allowbreak grid\_\allowbreak n2\_\allowbreak floor.f90} & the photon grid built when the excited-hydrogen coupling is armed & 184 \\
\texttt{tests/\allowbreak carrier\_\allowbreak reference\_\allowbreak scales/\allowbreak carrier\_\allowbreak reference\_\allowbreak scales.f90} & the element a transported carrier is made of, and whether its chemistry rows can be differentiated & 181 \\
\texttt{tests/\allowbreak carrier\_\allowbreak constraint\_\allowbreak attribution/\allowbreak carrier\_\allowbreak constraint\_\allowbreak attribution.f90} & which residuals belong to the solver and which to a constraint & 179 \\
\texttt{tests/\allowbreak adv\_\allowbreak static\_\allowbreak limit/\allowbreak enthalpy\_\allowbreak flux\_\allowbreak of\_\allowbreak mass\_\allowbreak divergence.f90} & the third term of the advected energy equation & 178 \\
\texttt{tests/\allowbreak grid\_\allowbreak and\_\allowbreak gates/\allowbreak free\_\allowbreak outflow\_\allowbreak boundary.f90} & the outer free-outflow boundary on a supersonic wind, and its ghost state & 174 \\
\texttt{tests/\allowbreak spectrum\_\allowbreak type/\allowbreak planck\_\allowbreak photon\_\allowbreak field.f90} & the incident flux put on the photon grid for the photospheric blackbody type & 173 \\
\texttt{tests/\allowbreak a2\_\allowbreak m2/\allowbreak a2\_\allowbreak m2\_\allowbreak kinetics\_\allowbreak check.f90} & the kinetics quantities the oxygen option adds on top of the rate audit & 163 \\
\texttt{tests/\allowbreak physics\_\allowbreak probe/\allowbreak lya\_\allowbreak escape\_\allowbreak probability.f90} & the escape probability of the two-level closure against the published slab solution & 163 \\
\texttt{tests/\allowbreak e1\_\allowbreak h2/\allowbreak e1\_\allowbreak h2\_\allowbreak channel\_\allowbreak check.f90} & the mutually exclusive H$_2$ photoabsorption channels & 161 \\
\texttt{tests/\allowbreak adv\_\allowbreak static\_\allowbreak limit/\allowbreak upstream\_\allowbreak caloric\_\allowbreak energy.f90} & the upstream energy of the advected energy equation & 154 \\
\texttt{tests/\allowbreak physics\_\allowbreak probe/\allowbreak photoionization\_\allowbreak field\_\allowbreak self\_\allowbreak consistency.f90} & the field one cell is solved in, against the composition that cell returns & 148 \\
\texttt{tests/\allowbreak grid\_\allowbreak and\_\allowbreak gates/\allowbreak photon\_\allowbreak grid\_\allowbreak quadrature.f90} & the photoionization rates of one unattenuated cell against the band quadrature & 146 \\
\texttt{tests/\allowbreak physics\_\allowbreak probe/\allowbreak co\_\allowbreak helium\_\allowbreak ion\_\allowbreak sink.f90} & He$^+$ $+$ CO as a species sink, not only as a heat source & 145 \\
\texttt{tests/\allowbreak physics\_\allowbreak probe/\allowbreak riemann\_\allowbreak wave\_\allowbreak speeds.f90} & signal speeds at a face, and the state one conservative update with that flux leaves behind & 143 \\
\texttt{tests/\allowbreak physics\_\allowbreak probe/\allowbreak helium\_\allowbreak level\_\allowbreak cooling\_\allowbreak ledger.f90} & which helium level each cooling channel of \texttt{eval\_cool} is charged to & 139 \\
\texttt{tests/\allowbreak physics\_\allowbreak probe/\allowbreak heating\_\allowbreak channel\_\allowbreak closure.f90} & one assembly of the volumetric heating rate, and the closure of its channels & 136 \\
\texttt{tests/\allowbreak acceptance\_\allowbreak classes/\allowbreak acceptance\_\allowbreak class\_\allowbreak contract.f90} & one meaning of the ionization acceptance classes for every consumer & 132 \\
\texttt{tests/\allowbreak physics\_\allowbreak probe/\allowbreak opacity\_\allowbreak model\_\allowbreak p\_\allowbreak ledger.f90} & the radiative ledger of a cell whose cross sections are broadened by opacity model P & 124 \\
\texttt{tests/\allowbreak physics\_\allowbreak probe/\allowbreak photoevent\_\allowbreak energy\_\allowbreak ledger.f90} & the energy an isolated photoabsorption event has to account for & 117 \\
\texttt{tests/\allowbreak physics\_\allowbreak probe/\allowbreak xuv\_\allowbreak cell\_\allowbreak mean\_\allowbreak attenuation.f90} & the beam attenuation the photoionization rate of a cell is evaluated with & 115 \\
\texttt{tests/\allowbreak grid\_\allowbreak and\_\allowbreak gates/\allowbreak grid\_\allowbreak window\_\allowbreak indices.f90} & the two index statements that decide which cells a quantity is read from & 101 \\
\texttt{tests/\allowbreak physics\_\allowbreak probe/\allowbreak h2\_\allowbreak rovibrational\_\allowbreak identity.f90} & the identities tying the H$_2$ partition function of the chemistry to the internal energy of the caloric EOS & 96 \\
\texttt{tests/\allowbreak adv\_\allowbreak static\_\allowbreak limit/\allowbreak stationary\_\allowbreak flux\_\allowbreak refusal.f90} & the sensitivity of the advected energy correction to a non-stationary mass flux & 91 \\
\texttt{tests/\allowbreak physics\_\allowbreak probe/\allowbreak ionization\_\allowbreak threshold\_\allowbreak turn\_\allowbreak on.f90} & one definition of each ionization threshold, and every cross section turning on exactly there & 90 \\
\texttt{tests/\allowbreak grid\_\allowbreak and\_\allowbreak gates/\allowbreak grid\_\allowbreak width\_\allowbreak identity.f90} & the cell-width and total-width identities of the radial grid & 83 \\
\texttt{tests/\allowbreak physics\_\allowbreak probe/\allowbreak h2\_\allowbreak channel\_\allowbreak detector\_\allowbreak ratio.f90} & the measured proton-to-H$_2^+$ ratio of Chung et al.\ (1993) against the four channels & 83 \\
\texttt{tests/\allowbreak physics\_\allowbreak probe/\allowbreak fine\_\allowbreak structure\_\allowbreak escape\_\allowbreak probability.f90} & the escape probability of the ground-term fine-structure lines, and the depth it is fed & 80 \\
\texttt{tests/\allowbreak physics\_\allowbreak probe/\allowbreak atomic\_\allowbreak mass\_\allowbreak and\_\allowbreak radius\_\allowbreak constants.f90} & the helium mass and the astronomical constants of the input readers, one definition each & 79 \\
\texttt{tests/\allowbreak physics\_\allowbreak probe/\allowbreak h3p\_\allowbreak cooling\_\allowbreak limits.f90} & the limits, the continuities and the published values of the H$_3^+$ cooling rate & 79 \\
\texttt{tests/\allowbreak physics\_\allowbreak probe/\allowbreak assertion\_\allowbreak report.f90} & one printed verdict line per assertion, and the failure count the driver exits on & 74 \\
\texttt{tests/\allowbreak adv\_\allowbreak static\_\allowbreak limit/\allowbreak thermal\_\allowbreak damkohler\_\allowbreak limit.f90} & the validity condition of the advection-corrected energy solve, as a function of the flow speed & 72 \\
\texttt{tests/\allowbreak a2\_\allowbreak m2/\allowbreak a2\_\allowbreak m2\_\allowbreak g3\_\allowbreak run\_\allowbreak check.f90} & with every photolysis rate at zero, the solved O/OH/H$_2$O partition is the chemical equilibrium & 71 \\
\texttt{tests/\allowbreak physics\_\allowbreak probe/\allowbreak absorbed\_\allowbreak fraction\_\allowbreak switch.f90} & $(1-e^{-d})/d$, the fraction of the photons entering a cell that the cell absorbs & 71 \\
\texttt{tests/\allowbreak constrained\_\allowbreak network\_\allowbreak layout/\allowbreak transported\_\allowbreak proton\_\allowbreak constraint\_\allowbreak layout.f90} & which species are unknowns of the continuation, and whether the system built is square & 71 \\
\texttt{tests/\allowbreak test\_\allowbreak columns.f90} & synthetic columns whose species carry the mass density they are measured against & 67 \\
\texttt{tests/\allowbreak physics\_\allowbreak probe/\allowbreak h2\_\allowbreak infrared\_\allowbreak line\_\allowbreak populations.f90} & the LTE populations weighting the H$_2$ quadrupole lines are those of the caloric EOS ladder & 65 \\
\texttt{tests/\allowbreak physics\_\allowbreak probe/\allowbreak wind\_\allowbreak ae\_\allowbreak bridge\_\allowbreak composition.f90} & the composition an EXHALE \texttt{input.inp} hands the Wind-AE solver & 54 \\
\texttt{tests/\allowbreak physics\_\allowbreak probe/\allowbreak voronov\_\allowbreak carbon\_\allowbreak ionization.f90} & the electron-impact ionization of C$^0$ and C$^+$ against the published Voronov (1997) rows & 51 \\
\texttt{tests/\allowbreak physics\_\allowbreak probe/\allowbreak water\_\allowbreak photolysis\_\allowbreak lyman\_\allowbreak alpha\_\allowbreak yields.f90} & the H$_2$O branching on the band that contains Ly$\alpha$, against the published yields & 49 \\
\texttt{tests/\allowbreak physics\_\allowbreak probe/\allowbreak h3p\_\allowbreak recombination\_\allowbreak branching.f90} & the two product channels of H$_3^+$ dissociative recombination & 44 \\
\texttt{tests/\allowbreak energy\_\allowbreak update/\allowbreak cooling\_\allowbreak models.f90} & analytic cooling laws for the tests of the implicit energy update & 43 \\
\texttt{tests/\allowbreak species\_\allowbreak masses/\allowbreak carrier\_\allowbreak background\_\allowbreak masses.f90} & the mass the carrier transport gives each of its background collision partners & 36 \\
\midrule
\multicolumn{2}{r}{test-driver total} & 20\,741 \\
\midrule
\multicolumn{2}{r}{new-module total} & 50\,744 \\
\bottomrule
\end{longtable}
}

\noindent
So the 123 new core modules contribute 50\,744 of the 69\,179 added code lines;
the remaining 18\,435 added (and 1\,121 of the 1\,190 removed, the other 69
being the deleted HLLC estimator) are edits spread across the 38 modified
files. The ten heaviest, by added code lines, are the main driver
(\texttt{EXHALE\_\allowbreak main.f90}, 3\,498 added, 246 $\to$ 3\,625 SLOC: the marching
loop, the attempted-step adoption boundary, the run-mode contracts and the
written-state assertions), the ionization equilibrium and its diagnostics
(\texttt{ionization\_\allowbreak equilibrium.f90} 2\,274 and \texttt{util\_\allowbreak ion\_\allowbreak eq.f90}
1\,803: the metal, molecular and oxygen networks, the acceptance classes, the
photoionization field and the heating channels), the cooling module
(\texttt{Cool\_\allowbreak coeff.f90}, 2\,250 added, 442 $\to$ 3\,966 raw lines: CHIANTI
metal-line fits, Fe tables, fine-structure statistical equilibrium, Penning
and molecular rate coefficients), input parsing
(\texttt{input\_\allowbreak read.f90}, 1\,753), the state reader and writer
(\texttt{load\_\allowbreak IC.f90} 1\,063, \texttt{write\_\allowbreak output.f90} 928,
\texttt{write\_\allowbreak setup\_\allowbreak report.f90} 751), the advection-corrected post-processor
(\texttt{post\_\allowbreak process\_\allowbreak adv.f90}, 647) and the cross sections
(\texttt{cross\_\allowbreak sec.f90}, 438). Those ten carry 15\,405 of the 18\,435. The rest
of the edits are in the reconstruction, the fluxes, the boundary conditions,
the grid, the time step, the energy equation and the \texttt{System\_*}
ionization systems, which wire the metals, the molecules, the opacity
dispatcher, the Roche potential, the transonic/auto/Wind-AE IC, the
excited-H and Ly$\alpha$ physics, the element diffusion and the Brent/Newton
solver upgrades into the original solver.

In total the EXHALE core is now $\approx$20.3$\times$ the original code base
(3\,514\,$\to$\,71\,503 SLOC), or $\approx$14.4$\times$ counting only the
production build (50\,762 SLOC), touching almost all of the original files (38
of 47 modified, one deleted, plus 123 new); with the ported Wind-AE IC
generator the full Fortran tree is 205 files / 76\,259 SLOC
($\approx$21.7$\times$). Beyond Fortran, the tree also carries the Python
tooling (GUI, the \texttt{EXHALE\_transit.py} transmission post-processor,
plotting/analysis, the lower-atmosphere adapters and the elemental-flux
closure driver) and 181\,816 lines of documentation (187 md/tex files: the
memos, the manual and this log) under \texttt{docs/}.


% ===================================================================== %
\section{Fortran source inherited unchanged from ATES}\label{sec:untouched}
% ===================================================================== %
Of the 47 upstream Fortran files (\texttt{ATES/ATES-Code-main}), \textbf{8 are
byte-for-byte identical} in EXHALE (re-measured 2026-09-11): the numerics
EXHALE reuses without touching, which are now the linear-algebra part of the
included \textsc{MINPACK} nonlinear solver plus the gravity pre-evaluation.
Table~\ref{tab:untouched} lists them with their role. No file is identical
after whitespace and case normalization but not byte-identical, so the
normalized-unchanged set and the byte-identical set are the same eight.

Five files left the byte-identical set since the 2026-08-30 edition of this
table, each for a documented reason rather than drift.

\begin{itemize}\itemsep2pt
\item \texttt{nonlinear\_\allowbreak system\_\allowbreak solver/\allowbreak hybrd.f90},
      \texttt{hybrd1.f90} and \texttt{fdjac1.f90} carry a new
      \texttt{unit\_step\_floor} argument threaded from \texttt{hybrd} to
      \texttt{fdjac1}, which selects the forward-difference step of the
      Jacobian: \texttt{.false.} is MINPACK's own $h = \varepsilon|x_j|$ with
      the exact-zero fallback, \texttt{.true.} floors it at unit magnitude,
      $h = \varepsilon\max(1,|x_j|)$, which is the rule for unknowns already
      scaled so that unity is their natural size (the logarithmic unknowns of
      the constrained chemical equilibrium). The \texttt{params} dummy also
      became assumed-size, so one transport array serves systems of different
      length. \texttt{hybrd1} passes \texttt{.false.} and is therefore
      unchanged in behaviour.
\item \texttt{states/\allowbreak Source.f90} gained the well-balanced arm of
      2026-09-10 (item N37 of the stage-2 solver program,
      \texttt{docs/ISSUES\_20260909.md}): under
      \texttt{Well balanced: True} the gravitational source and the geometric
      pressure term are not evaluated here at all, because they cancel the
      equilibrium part of the momentum flux difference identically and
      \texttt{RK\_rhs} assembles the row from what is left (K\"appeli \&
      Mishra 2014, J. Comput. Phys. 259, 199, eq. 2.26). The file also reads
      its pressure through \texttt{caloric\_eos} instead of a constant
      $\gamma$.
\item \texttt{radiation/\allowbreak J\_\allowbreak inc.f90} carries the single-spectrum-type decision:
      \texttt{Spectrum type:} now states the spectrum of the whole photon
      grid, the XUV and the part below 13.6\,eV alike, and no band is filled
      from a type the input did not select. The file also became the one
      definition of the eV-to-Hz conversion and of the Balmer-continuum head
      $e_{\rm th,HI}/4$.
\end{itemize}

\noindent
The two wave-speed estimators are no longer in the ``effectively untouched''
category either. \texttt{flux/\allowbreak speed\_\allowbreak estimate\_\allowbreak HLLC.f90} was deleted as dead
code, as recorded above. \texttt{flux/\allowbreak speed\_\allowbreak estimate\_\allowbreak ROE.f90} was rewritten
(99 code lines added, 43 removed) into a star-state estimate that returns a
status: the primitive-variable, two-rarefaction and two-shock branches of
Toro (2009) chapter 9 are selected by the pressure ratio, the vacuum
condition of his eq. 4.76 is tested first, and a face with no admissible star
state is reported rather than given a square root of a negative ratio.

\begin{table}[h]
\centering\small
\caption{The 8 Fortran files byte-identical between the original ATES
(\texttt{ATES-Code-main}) and EXHALE (measured 2026-09-11), grouped by role.}
\label{tab:untouched}
\begin{tabular}{@{}l p{0.60\linewidth}@{}}
\toprule
File & Purpose \\
\midrule
\multicolumn{2}{@{}l}{\emph{MINPACK linear algebra and constants} (Argonne; dir \texttt{nonlinear\_system\_solver/})}\\
\texttt{qrfac.f90}   & QR factorization via Householder transformations \\
\texttt{qform.f90}   & Accumulate the orthogonal factor $Q$ from the QR factorization \\
\texttt{r1updt.f90}  & Rank-one update of the QR factorization \\
\texttt{r1mpyq.f90}  & Apply the Givens-rotation sequence of the rank-one update \\
\texttt{dogleg.f90}  & Dogleg (trust-region) step combining Gauss--Newton and steepest descent \\
\texttt{enorm.f90}   & Euclidean norm of a vector (overflow/underflow safe) \\
\texttt{dpmpar.f90}  & Machine floating-point precision / range constants \\
\midrule
\multicolumn{2}{@{}l}{\emph{Grid setup}}\\
\texttt{init/\allowbreak set\_\allowbreak gravity\_\allowbreak grid.f90}   & Pre-evaluate gravity at cell centers and interfaces \\
\bottomrule
\end{tabular}
\end{table}

\end{document}
